% Copyright 2025-2026 by Jonatan Perren
%
% This file may be distributed and/or modified
%
% 1. under the LaTeX Project Public License and/or
% 2. under the GNU Free Documentation License.


\documentclass[a4paper,doc2]{ltxdoc}



\input{../pgfgo.cfg}
% Copyright 2019 by Till Tantau
%
% This file may be distributed and/or modified
%
% 1. under the LaTeX Project Public License and/or
% 2. under the GNU Free Documentation License.
%
% See the file doc/generic/pgf/licenses/LICENSE for more details.


% pgf version is defined in \pgfversion in file
% generic/pgf/utilities/pgfrcs.code.tex

\def\xcolorversion{2.00}

\usepackage[version=latest]{pgf}

\usepackage{xkeyval,calc,listings,tikz,fp}

\usepackage{imakeidx}
\makeindex
\usepackage{hyperref}
\hypersetup{%
    colorlinks=false, % use true to enable colors below:
    linkcolor=blue,%red,
    filecolor=blue,%magenta,
    urlcolor=blue,%cyan,
    citecolor=blue,
    pdfborder=0 0 0,
}

% We need lots of libraries...
\usetikzlibrary{
  3d,
  arrows,
  arrows.spaced,
  arrows.meta,
  bending,
  babel,
  calc,
  fit,
  patterns,
  patterns.meta,
  plotmarks,
  shapes.geometric,
  shapes.misc,
  shapes.symbols,
  shapes.arrows,
  shapes.callouts,
  shapes.multipart,
  shapes.gates.logic.US,
  shapes.gates.logic.IEC,
  circuits.logic.US,
  circuits.logic.IEC,
  circuits.logic.CDH,
  circuits.ee.IEC,
  datavisualization,
  datavisualization.polar,
  datavisualization.formats.functions,
  er,
  automata,
  backgrounds,
  chains,
  topaths,
  trees,
  petri,
  mindmap,
  matrix,
  calendar,
  folding,
  fadings,
  shadings,
  spy,
  through,
  turtle,
  positioning,
  scopes,
  decorations.fractals,
  decorations.shapes,
  decorations.text,
  decorations.pathmorphing,
  decorations.pathreplacing,
  decorations.footprints,
  decorations.markings,
  shadows,
  lindenmayersystems,
  intersections,
  fixedpointarithmetic,
  fpu,
  svg.path,
  external,
  graphs,
  graphs.standard,
  quotes,
  math,
  angles,
  views,
  animations,
  rdf,
  perspective,
}

\usepackage{ifluatex}
\newif\ifgdccodebasic
\newif\ifgdccodeogdf

\ifluatex

  \usetikzlibrary{graphdrawing}
  \usegdlibrary{trees,circular,layered,examples,force,phylogenetics,routing}

  % Test whether C code is available:
  \directlua{
    if pcall(require,'pgf_gd_examples_c_SimpleDemoC') then
      tex.print('\string\\gdccodebasictrue')
    end
  }

  \ifgdccodebasic
    \usegdlibrary{pgf_gd_examples_c_SimpleDemoC,pgf_gd_examples_c_SimpleDemoCPlusPlus}

    % Test whether OGDF code is available:
    \directlua{
      if pcall(require,'pgf_gd_ogdf_c_SimpleDemoOGDF') then
        tex.print('\string\\gdccodeogdftrue')
      end
    }

    \ifgdccodeogdf
      \usegdlibrary{pgf_gd_ogdf_c_SimpleDemoOGDF,ogdf}
    \fi

  \fi

\fi

\def\LuaTeX{Lua\TeX}%


\newif\ifpgfmanualexternalize
\pgfmanualexternalizefalse
\ifpgfmanualexternalize
  \tikzexternalize[
    mode=list only,export=true,% simply skips EVERY picture -> good for debugging the text.
  ]{pgfmanual}

  \tikzifexternalizing{%
    \pgfkeys{/pdflinks/codeexample links=false}%
  }{}%
\fi


\usepackage[a4paper,left=2.25cm,right=2.25cm,top=2.5cm,bottom=2.5cm,nohead]{geometry}
\usepackage{amsmath,amssymb}
\usepackage{xxcolor}
\usepackage{pifont}

\ifluatex
%  \usepackage[no-math]{fontspec}
%  \usepackage{luatextra}

  % TT: I have commented luatextra since it loads fontspec, which
  % currently breaks "$\mathrm{\Omega}$" (nothing is
  % shown). Also, fontspec loads so much stuff, I'm not really sure
  % any of it is needed...

  % \filedescription is defined in expl3, required by fontspec,
  % required by luatextra. Needs to be \relaxed since
  % pgfmanual-en-macros.tex defines an environment named filedescription
  \let\filedescription\relax
  \usepackage[utf8]{luainputenc}
  % dvisvgm does not support OpenType fonts so we have to bite the bullet and
  % use T1 where quotes are weird.  In LuaTeX \outputmode=0 means DVI.
  \ifnum\outputmode=0
    \usepackage[T1]{fontenc}
  \fi
\else
  \usepackage[utf8]{inputenc}
  \usepackage[T1]{fontenc}
\fi

% Promote `Missing character` reports to full errors
% require texlive 2021 or above
\tracinglostchars=3

\graphicspath{{../../img/}}
\input{pgfmanual-en-macros}

\usepackage{todonotes}
\newcommand\todosp[2][]{%        % Stefan Pinnow
    \todo[
        disable,
        color=blue!40,
        #1]{#2}
}

\makeatletter
\renewcommand*\l@section[2]{%
  \ifnum \c@tocdepth >\z@
    \addpenalty\@secpenalty
    \addvspace{1.0em \@plus\p@}%
    \setlength\@tempdima{2.5em}%
    \begingroup
      \parindent \z@ \rightskip \@pnumwidth
      \parfillskip -\@pnumwidth
      \leavevmode \bfseries
      \advance\leftskip\@tempdima
      \hskip -\leftskip
      #1\nobreak\hfil \nobreak\hb@xt@\@pnumwidth{\hss #2}\par
    \endgroup
  \fi}
\renewcommand*\l@subsection{\@dottedtocline{2}{2.5em}{3.3em}}
\renewcommand*\l@subsubsection{\@dottedtocline{3}{5.8em}{4.2em}}
\def\@pnumwidth{2.2em}
\makeatother

%\includeonly{pgfmanual-en-library-profiler}

% Global styles:
\tikzset{
  every plot/.style={prefix=plots/pgf-},
  shape example/.style={
    color=black!30,
    draw,
    fill=yellow!30,
    line width=.5cm,
    inner xsep=2.5cm,
    inner ysep=0.5cm}
}

\index{Options for graphics|see{Graphic options and styles}}
\index{Styles for graphics|see{Graphic options and styles}}
\index{Options for packages|see{Package options}}
\index{Handlers for keys|see{Key handlers}}
\index{File|see{Packages and files}}
\index{Layout|see{Page layout}}
\index{Node|see{Predefined node}}
\index{Data formats|see{Formats}}


% Animations:

\tikzset{make snapshot if necessary}

% Detect changed labels
% by David Carlisle https://tex.stackexchange.com/a/169245
\makeatletter
\def\@testdef#1#2#3{%
  \def\reserved@a{#3}%
  \expandafter\ifx\csname #1@#2\endcsname\reserved@a\else
    \typeout{^^JLabel `#2' changed:^^J\meaning\reserved@a^^J\expandafter\meaning\csname #1@#2\endcsname^^J}%
    \@tempswatrue
  \fi}
\makeatother

%% -----------------------------------------------------------------------------
%\includeonly{
%%    pgfmanual-en-introduction,
%%    pgfmanual-en-tutorial,
%%    pgfmanual-en-tutorial-nodes,
%%    pgfmanual-en-tutorial-Euclid,
%%    pgfmanual-en-tutorial-chains,
%%    pgfmanual-en-tutorial-map,
%%    pgfmanual-en-guidelines,
%%    pgfmanual-en-installation,
%%    pgfmanual-en-license,
%%    pgfmanual-en-drivers,
%%    pgfmanual-en-tikz-design,
%%    pgfmanual-en-tikz-scopes,
%%    pgfmanual-en-tikz-coordinates,
%%    pgfmanual-en-tikz-paths,
%%    pgfmanual-en-tikz-actions,
%%    pgfmanual-en-tikz-arrows,
%%    pgfmanual-en-tikz-shapes,
%%    pgfmanual-en-tikz-pics,
%%    pgfmanual-en-tikz-graphs,
%%    pgfmanual-en-tikz-matrices,
%%    pgfmanual-en-tikz-trees,
%%    pgfmanual-en-tikz-plots,
%%    pgfmanual-en-tikz-transparency,
%%    pgfmanual-en-tikz-decorations,
%%    pgfmanual-en-tikz-transformations,
%%    pgfmanual-en-tikz-animations,
%%    pgfmanual-en-gd-overview,
%%    pgfmanual-en-gd-usage-tikz,
%%    pgfmanual-en-gd-usage-pgf,
%%    pgfmanual-en-gd-trees,
%%    pgfmanual-en-gd-layered,
%%    pgfmanual-en-gd-force,
%%    pgfmanual-en-gd-circular,
%%    pgfmanual-en-gd-phylogenetics,
%%    pgfmanual-en-gd-edge-routing,
%%%    pgfmanual-en-gd-ogdf,
%%    pgfmanual-en-gd-algorithm-layer,
%%    pgfmanual-en-gd-algorithms-in-c,
%%    pgfmanual-en-gd-display-layer,
%%    pgfmanual-en-gd-binding-layer,
%%    pgfmanual-en-library-angles,
%%    pgfmanual-en-library-arrows,
%%    pgfmanual-en-library-automata,
%%    pgfmanual-en-library-babel,
%%    pgfmanual-en-library-backgrounds,
%%    pgfmanual-en-library-calc,
%%    pgfmanual-en-library-calendar,
%%    pgfmanual-en-library-chains,
%%    pgfmanual-en-library-circuits,
%%    pgfmanual-en-library-decorations,
%%    pgfmanual-en-library-er,
%%    pgfmanual-en-library-external,
%%    pgfmanual-en-library-fadings,
%%    pgfmanual-en-library-fit,
%%    pgfmanual-en-library-fixedpoint,
%%    pgfmanual-en-library-fpu,
%%    pgfmanual-en-library-lsystems,
%%    pgfmanual-en-library-math,
%%    pgfmanual-en-library-matrices,
%%    pgfmanual-en-library-mindmaps,
%%    pgfmanual-en-library-folding,
%%    pgfmanual-en-library-patterns,
%%    pgfmanual-en-library-petri,
%%    pgfmanual-en-library-plot-handlers,
%%    pgfmanual-en-library-plot-marks,
%%    pgfmanual-en-library-profiler,
%%    pgfmanual-en-library-rdf,
%%    pgfmanual-en-library-shadings,
%%    pgfmanual-en-library-shadows,
%%    pgfmanual-en-library-shapes,
%%    pgfmanual-en-library-spy,
%%    pgfmanual-en-library-svg-path,
%%    pgfmanual-en-library-edges,
%%    pgfmanual-en-library-through,
%%    pgfmanual-en-library-trees,
%%    pgfmanual-en-library-turtle,
%%    pgfmanual-en-library-views,
%%    pgfmanual-en-dv-introduction,
%%    pgfmanual-en-dv-main,
%%    pgfmanual-en-dv-formats,
%%    pgfmanual-en-dv-axes,
%%    pgfmanual-en-dv-visualizers,
%%    pgfmanual-en-dv-stylesheets,
%%    pgfmanual-en-dv-polar,
%%    pgfmanual-en-dv-backend,
%%    pgfmanual-en-pgfkeys,
%%    pgfmanual-en-pgffor,
%%    pgfmanual-en-pgfcalendar,
%%    pgfmanual-en-pages,
%%    pgfmanual-en-xxcolor,
%%    pgfmanual-en-module-parser,
%%    pgfmanual-en-math-design,
%%    pgfmanual-en-math-parsing,
%%    pgfmanual-en-math-commands,
%%    pgfmanual-en-math-algorithms,
%%    pgfmanual-en-math-numberprinting,
%%    pgfmanual-en-oo,
%%    pgfmanual-en-base-design,
%%    pgfmanual-en-base-scopes,
%%    pgfmanual-en-base-points,
%%    pgfmanual-en-base-paths,
%%    pgfmanual-en-base-decorations,
%%    pgfmanual-en-base-actions,
%%    pgfmanual-en-base-arrows,
%%    pgfmanual-en-base-nodes,
%%    pgfmanual-en-base-matrices,
%%    pgfmanual-en-base-transformations,
%%    pgfmanual-en-base-patterns,
%%    pgfmanual-en-base-images,
%%    pgfmanual-en-base-external,
%%    pgfmanual-en-base-plots,
%%    pgfmanual-en-base-layers,
%%    pgfmanual-en-base-shadings,
%%    pgfmanual-en-base-transparency,
%%    pgfmanual-en-base-animations,
%%    pgfmanual-en-base-internalregisters,
%%    pgfmanual-en-base-quick,
%%    pgfmanual-en-pgfsys-overview,
%%    pgfmanual-en-pgfsys-commands,
%%    pgfmanual-en-pgfsys-paths,
%%    pgfmanual-en-pgfsys-protocol,
%%    pgfmanual-en-pgfsys-animations,
%}
%% -----------------------------------------------------------------------------


%%% Local Variables:
%%% mode: latex
%%% TeX-master: "pgfmanual"
%%% End:


\usepackage{animate}
\usepackage{../../latex/pgf-go}
%LOS ARCHIVOS SGF DE LOS EJEMPLOS, POR SU NOMBRE (LOS TUTORIALES Y EL CAPITULO SGF LOS PRESENTAN).
\pgfgodefinesgf{partida}{partida.sgf}
\pgfgodefinesgf{lee}{lee.sgf}
\pgfgodefinesgf{variantes}{variantes.sgf}
\pgfgodefinesgf{problema}{problema.sgf}
\pgfgodefinesgf{josekis}{josekis.sgf}
%LOS KANJI DEL TEXTO (VER label/style).
\font\pgfgokanji = {file:ipaexm.ttf} at 9pt
%LA MISMA FUENTE COMO FAMILIA, PARA format = kanji (VER mark/label/format).
\DeclareFontFamily{TU}{ipaex}{}
\DeclareFontShape{TU}{ipaex}{m}{n}{<->file:ipaexm.ttf}{}
\DeclareFontShape{TU}{ipaex}{bx}{n}{<->ssub*ipaex/m/n}{}
\def\pgfgoname{\textsc{pgfgo}}
\def\karu{K\={a}ru}

\usepackage[spanish]{babel}

\usepackage{float}
\usetikzlibrary{positioning, fit}
\usepackage{caption}



\definecolor{pgfgo}{RGB}{220,179,92}
\colorlet{tablero2}{red!10!blue!10}
\colorlet{tablero3}{blue!10!green!10}
\definecolor{tablero4}{HTML}{9B8465}

\def\pgfgoversion{\leavevmode\hbox{0.2026.10.14}}
\def\twodigits#1{\ifnum#1<10 0\fi\the#1}



\begin{document}
	\title{\Huge\bfseries pgf-go\\[1em]
		\normalsize
		
		\begin{center}
			\pgfgoset{profile = pgfgosmooth}
			\begin{goban}[partial = B2:E3]
				\stones[black]{C2,D2,E3}
				\stones[white] {B2,C3,D3}
				\mark[w, custom = {36}]{B2}
				\mark[custom = {a}] {E2}
				
				\marks[circle, line = {color = red!80!black}] {C3,D3}
				
			\end{goban}
		\end{center}
		
		\large Manual for version~\pgfgoversion}
		\author{Jonat\'an Perren Hug\\
		\texttt{jonatan.perren@gmail.com}
		}
	
	\maketitle
	
	
	
	\label{table-of-contents}
	\tableofcontents

	\clearpage 

	\part{Introducción}
	\section{Introducci\'on}

\subsection{Motivaci\'on}

El siguiente paquete ofrece una colección de macros diseñadas para representar tableros y posiciones de Go, implementadas íntegramente en \pgfname. Este proyecto, aunque aún en desarrollo, incorpora funcionalidades innovadoras que buscan simplificar y agilizar la creación de diagramas relacionados con el juego de Go:
\begin{itemize}
	\item Gestión de perfiles para tableros y jugadores.
	\item Carga de múltiples coordenadas simultáneamente mediante una sintaxis que algunos podrían encontrar familiar.
\end{itemize}

El desarrollo de \pgfgoname\ surge de la frustración por no encontrar paquetes que satisficieran mis necesidades. Los paquetes que exploré permiten representar tableros en blanco y negro, adecuados para impresión, pero poco prácticos al combinarlos con tecnolog\'ias como |beamer| o en documentos donde se desea una presentación a color.

Inicialmente, \pgfgoname\ se limitaría a estas tareas; sin embargo, ahora lo concibo más como una interfaz para gestionar estilos y estados, evitando el uso excesivo del formato directo. Esto se reflejará principalmente en la estructura del árbol de \emph{keys}.

En cuanto al nombre, aunque las macros están implementadas en \pgfname\ y L3, el tablero se construye sobre |\tikzpicture| para que los usuarios familiarizados con este entorno puedan aprovechar sus macros.


\subsection{Novedades}


\begin{itemize}
	\item \pgfgoversion:
		\begin{itemize}
			\item El cálculo (coordenadas, numeración de jugadas, memoria, etc.) se ha separado del dibujo. Lo hace un \emph{kernel}, que se elige con la opción |kernel| del paquete (pág.~\pageref{subsec:kernel}). Además del kernel escrito en L3, se ha añadido un kernel en Lua, que es el que se usa por defecto con Lua\LaTeX.
			\item Con el kernel en Lua se han añadido las capturas (|captures|, pág.~\pageref{subsec:capturas}) y la lectura de partidas en SGF (|\sgf|, pág.~\pageref{subsec:sgf}). Con el paquete |animate|, |\sgfanimation| permite recorrer una partida con botones.
			\item Se han añadido las notas «N en M» para las jugadas numeradas en puntos ya jugados (|notes| y |\pgfgomovenotes|, pág.~\pageref{subsec:notas}).
			\item SGF: variantes (|variation|), marcas del archivo (|marks|), los datos de la partida (|\sgfinfo|), la numeración de |MN| y |restart|, y el tamaño del tablero tomado del archivo (|sgf size|).
			\item SGF: colecciones de partidas (|game|), la vista (|VW|, |sgf view|), las figuras (|FG|) en |\sgfdiagrams|, las variantes como letras (|variations|), |PL| y los archivos en Latin-1.
			\item SGF: el nombre de los nodos y las anotaciones (|TE|, |BM|, |GB|, etc.) se escriben con los comentarios. |\sgfcomment| escribe el comentario de una jugada, y |\sgfdiagrams| puede escribir los comentarios debajo de cada diagrama (|comments|).
			\item Se han añadido flechas y líneas (|\arrow| y |\line|, pág.~\pageref{subsec:flechas}) y la marca |dim|. |\sgf| dibuja además |AR|, |LN|, |SL| y |DD|.
			\item El tablero dibuja los puntos de estrella (|hoshi|, pág.~\pageref{subsec:hoshi}).
			\item Se han añadido |\territory|, que marca el territorio (con un polígono, o solo con |auto| y el kernel en Lua), y |\pgfgoscore|, que cuenta los puntos.
			\item |\sgfanimation| arma la animación por capas (|layers|): cada piedra se dibuja una sola vez, y el PDF de una partida larga ocupa mucho menos.
			\item Con |captured marks = remove|, las marcas sobre una piedra capturada se van con ella.
			\item Se han añadido |\pgfgoprisoners|, que cuenta los prisioneros de cada jugador, y |\sgfdiagrams|, que dibuja una partida como serie de diagramas, como en los libros.
			\item Se ha añadido |\sgfproblem|, que dibuja un problema (tsumego): el enunciado y un diagrama por cada variante de la solución, con sus comentarios.
			\item SGF: la numeración que pide |PM| (|numbers| y |restart = auto|), y la vista (|VW|) cuando cambia durante la partida (|sgf view move|; |\sgfdiagrams| y |\sgfanimation| la usan solos). |\sgfdiagrams| escribe el resultado y el puntaje debajo del último diagrama (|result|).
			\item Las coordenadas del borde pueden ser letras, números o kanji en cada eje, contados desde cualquier lado (|columns|, |rows|, |style = japanese|, etc.).
		\end{itemize}
	\item 0.2025.10.13: Corrección de errores. Cada |goban| empieza en la jugada 1 con negras, |\mark| y |\marks| solo existen dentro de |goban|, y lo que cae fuera de un tablero parcial ya no se dibuja. La lista completa está en el archivo |CHANGELOG.md|.
	\item 0.2025.03.10:
		\begin{itemize}
			\item Se han a\~nadido las \emph{keys} |inner| y |outer| para el \emph{rounded corner} del tablero.
			\item Se han a\~nadido tres nuevos operadores de coordenadas: |*| para todas las intersecciones, |L| para toda la columna de coordenadas L y |N| para toda la fila de coordenadas N.
		\end{itemize}
	\item 0.2025.03.07: Primer lanzamiento.
\end{itemize}

	\part{Tutoriales}\label{part:tutoriales}
	\section{Los tableros de \karu}

\subsection{La primera partida}
\label{tut:primera}

\karu\ está escribiendo un artículo para compartir con sus estudiantes. Le preocupa el peso final de su PDF y ha notado que los \emph{softwares} para diagramar no ofrecen formatos ni resultados satisfactorios. Por ello, decide probar \pgfgoname. Lo primero que quiere mostrar es una partida corta que jugó en clase. La tiene guardada en un archivo SGF, |partida.sgf|: el formato en el que casi todos los programas y servidores de go guardan las partidas.
  
\begin{figure}[H]
	\centering
	\begin{goban}[captures, label = {at = {north, west}}]
		\sgf{partida}
	\end{goban}
\end{figure}

En \pgfgoname, todos los diagramas se dibujan dentro del entorno |goban|, y dentro de él, |\sgf| dibuja la partida del archivo. Para leer archivos SGF hay que compilar con Lua\LaTeX\ (pág.~\pageref{subsec:kernel}). Como va a usar el archivo muchas veces, primero le da un nombre con |\pgfgodefinesgf| (pág.~\pageref{subsec:sgfnombres}): desde ahí lo llama |partida|, y si el archivo cambia de lugar, lo corrige en un solo sitio\footnote{Desgraciadamente lo implement\'e muy avanzada la documentaci\'on.}.

\begin{codeexample}[width=4.8cm]
\pgfgodefinesgf{partida}{partida.sgf}

\begin{goban}
	\sgf{partida}
\end{goban}
\end{codeexample}

Las jugadas aparecen numeradas, como en los libros. Pero con la jugada 7 las negras capturaron la piedra blanca 2, que sigue en el tablero: \pgfgoname\ dibuja lo que se le pide, sin aplicar las reglas. Con la \emph{key} |captures|, el |goban| retira las piedras capturadas (pág.~\pageref{subsec:capturas}). Las opciones del tablero se escriben entre corchetes, después de |\begin{goban}|.

\begin{codeexample}[width=4.8cm]
\begin{goban}[captures]
	\sgf{partida}
\end{goban}
\end{codeexample}

Para comentar la partida en el texto, \karu\ la detiene en la jugada 7 con |to|. El archivo trae un comentario para esa jugada, y |\sgfcomment| lo escribe.

\begin{codeexample}[width=4.8cm]
\begin{goban}[captures]
	\sgf[to = 7]{partida}
\end{goban}

Jugada 7: \sgfcomment{partida}{7}
\end{codeexample}

Esta partida es de $9\times9$, el tamaño que tiene un |goban| por defecto. Para una partida de otro tamaño, |sgf size = partida| toma el tamaño del archivo; así lo hará \karu\ con una partida de $19\times19$ en la pág.~\pageref{tut:partida}. Allí también verá cómo dibujar una partida larga en una serie de diagramas.

\subsection{Una posición a mano}
\label{tut:posicion}

Para la introducción del artículo, \karu\ quiere una posición propia, que no está en ningún archivo: la de la Figura~\ref{fig:primer-ejemplo}. Tiene que dibujarla piedra por piedra.

\begin{figure}[H]
	\centering
	\pgfgoset{profile = pgfgo}
	\begin{goban}[label = {at = {north, west}}]
		\stones[black]{B2,E2,J2,C3,D3,G3,H3,C5,D5,E5,F5,G5,E6,G6}
		\stones[white]{D8,F8,E7,H7,C6,D6,F6,H6,B5,J5,B4,E4,J4,J3,H2}
	\end{goban}
	\caption{}
	\label{fig:primer-ejemplo}
\end{figure}

\karu\ empieza por el tablero vacío. Como en el tutorial anterior, todo se dibuja dentro del entorno |goban|, que se encarga del tamaño del tablero, su color, sus coordenadas y todo lo que se dibuja encima.

\begin{codeexample}[width=4.4cm]
\begin{goban}
\end{goban}
\end{codeexample}

El tablero es de $9\times 9$, justo lo que necesita, pero le faltan las coordenadas en los bordes. Las opciones del |goban| van entre corchetes, y las de cada parte del tablero se agrupan en tres \emph{keys}: |grid|, para la cuadrícula; |background|, para el borde y el relleno, y |label|, para las coordenadas (pág.~\pageref{sec:tablero}). Dentro de |label|, |at| dice en qué lados van: |north|, |south|, |east| y |west| (o |above|, |below|, |right| y |left|), |all| o |none|. \karu\ las quiere arriba y a la izquierda.

\begin{codeexample}[width=4.7cm]
\begin{goban}[label = {at = {north, west}}]
\end{goban}
\end{codeexample}

\karu\ nota que las columnas pasan de la |H| a la |J|: como es costumbre en el go, \pgfgoname\ omite la «I», tanto en las coordenadas como en las etiquetas laterales. Si alguna vez la necesita, la \emph{key} |goban/use i| (pág.~\pageref{subsec:usei}) la vuelve a poner.

Ahora sí, las piedras. La más simple es |\stone|, que recibe la intersección, como |D8|, y entre corchetes sus opciones: el jugador, el estilo, una marca. Sin jugador, la piedra es neutral; para las negras, \karu\ escribe |black|, o |b|.

\begin{codeexample}[width=4.8cm]
\begin{goban}[label = {at = {north, west}}]
	\stone[black]{D8}
	\stone[black]{F8}
	\stone[b]{E7}
	\stone[b]{H7}
\end{goban}
\end{codeexample}

\karu\ sospecha que esta no es la mejor opción si necesita añadir 200 piedras o más. Encuentra |\stones|, que recibe una lista de intersecciones separadas por comas y acepta las mismas opciones que |\stone|.\footnote{Entre \texttt{\textbackslash stone\{A1\}} y \texttt{\textbackslash stones\{A1\}}, la primera opción es mucho más rápida. Veremos más adelante que hay muchos procesos involucrados en la lectura del parámetro obligatorio de las macros plurales.}

\begin{codeexample}[width=4.8cm]
\begin{goban}[label = {at = {north, west}}]
	\stones[b]{B2,E2,J2,C3,D3,G3,H3,C5,D5,E5,F5,G5,E6,G6}
\end{goban}
\end{codeexample}

Las blancas llevan |white|, o |w|. En rigor, |black| y |white| no son colores sino jugadores, el primero y el segundo: cada piedra toma el estilo de su jugador (pág.~\pageref{sec:jugadores}).

\begin{codeexample}[width=4.8cm]
\begin{goban}[label = {at = {north, west}}]
	\stones[b]{B2,E2,J2,C3,D3,G3,H3,C5,D5,E5,F5,G5,E6,G6}
	\stones[w]{D8,F8,E7,H7,C6,D6,F6,H6,B5,J5,B4,E4,J4,J3,H2}
\end{goban}
\end{codeexample}

Falta el estilo. En el catálogo (pág.~\pageref{sec:catalogo}), \karu\ encuentra uno que le gusta: el perfil |pgfgo|. Lo aplica con |\pgfgoset{profile = pgfgo}| (pág.~\pageref{subsec:perfiles}), y desde ahí vale para todos los tableros que siguen.

\begin{codeexample}[width=4.8cm]
\pgfgoset{profile = pgfgo}

\begin{goban}[label = {at = {north, west}}]
	\stones[b]{B2,E2,J2,C3,D3,G3,H3,C5,D5,E5,F5,G5,E6,G6}
	\stones[w]{D8,F8,E7,H7,C6,D6,F6,H6,B5,J5,B4,E4,J4,J3,H2}
\end{goban}
\end{codeexample}

Las listas de este tutorial son largas, y la paciencia de \karu, corta. En el último tutorial (pág.~\pageref{tut:abreviar}), \karu\ aprende a escribirlas más cortas.

\subsection{Movimientos}
\label{tut:movimientos}

Para explicar cómo empezó una pelea, \karu\ quiere mostrar las piedras en el orden en que se jugaron. Con |\stone|, tendría que indicar el jugador de cada una:

\begin{codeexample}[code only]
\begin{goban}
	\stone[b]{F7}
	\stone[w]{E6}
	\stone[b]{G3}
	\stone[w]{D6}
\end{goban}
\end{codeexample}

Para eso está |\moves| (pág.~\pageref{subsec:movimientos}): recibe una lista como |\stones|, pero juega las intersecciones en el orden en que se escriben, alterna los colores y numera las jugadas.

\begin{codeexample}[width=5cm]
\begin{goban}[label = {at = all}]
	\moves{F7, E6, G3, D6}
\end{goban}
\end{codeexample}

La pelea empezó con una negra en F7, que \karu\ pone con |\stone|. Las jugadas que siguen no las va a comentar: con |numbers = false|, se dibujan sin números. Con |white|, la secuencia empieza con blancas: en |\moves|, |black| y |white| solo dicen quién juega primero.

\begin{codeexample}[width=5cm]
\begin{goban}[label = {at = all}]
	\stone[b] {F7}
	\moves[numbers = false, white]{E6,G3,D6,F5,E7,E2,F8,E6,D8,D3,
		G7,C5,C6,D5,B5,G5,B4,B2,H6,C3,
		J3,J2,J5,H3}
\end{goban}
\end{codeexample}

Las últimas jugadas sí las comenta, así que van con número. La numeración y el turno continúan de un |\moves| al siguiente, aunque las jugadas anteriores no lleven números, y con |from| se elige desde qué número se sigue.

\begin{codeexample}[width=5cm]
\begin{goban}[label = {at = all}]
	\moves[numbers = false, white] {E6,G3,D6,F5,E7,E2,F8,E6,D8,D3,
		G7,C5,C6,D5,B5,G5,B4,B2,H6,C3,
		J3,J2,J5,H3}
	
	\moves[from = 26]{F6,E5,J4,G6}
	\moves{H2} % Es blanca y etiquetada con 30, sin necesidad de expresarlo.
\end{goban}
\end{codeexample}

|H2| va en otro |\moves| y, sin indicarlo, es blanca y lleva el 30: la jugada anterior, |G6|, fue la 29, de negras.

\subsection{Tablero parcial y marcas}
\label{tut:parcial}

Para un ejercicio del artículo, \karu\ quiere mostrar solo una parte del tablero, con dos señales: un círculo en la piedra que está en atari y un «!» en el punto clave.

\begin{figure}[H]
	\centering
	\begin{goban}[partial = B2:E4, background = {line = false}]
		\stones[b] {C2,B3,C4}
		\stones[w] {D4,E3,D2}
		\stone[w, mark = {circle, line = {color = red}}] {C3}
		\mark[custom = {!}] {D3}
	\end{goban}
\end{figure}

Empieza, como siempre, por las piedras.

\begin{codeexample}[width=5cm]
	\begin{goban}
		\stones[b]{C2,B3,C4}
		\stones[w]{D4,E3,D2,C3}
	\end{goban}
\end{codeexample}

Con la \emph{key} |partial|, se indica qué parte del tablero mostrar: dos esquinas opuestas del rectángulo, separadas por dos puntos, como |B2:E4| (pág.~\pageref{subsec:parcial}).

\begin{codeexample}[]
	\begin{goban}[partial = B2:E4]
		\stones[b]{C2,B3,C4}
		\stones[w]{D4,E3,D2,C3}
	\end{goban}
\end{codeexample}

Si \karu\ no quiere buscar las esquinas, con |partial = auto| el tablero elige la parte: la que usa lo que se dibuja, con una línea más alrededor. Un lado que queda a dos líneas o menos del borde llega hasta el borde, como aquí el de abajo y el de la izquierda.

\begin{codeexample}[]
	\begin{goban}[partial = auto]
		\stones[b]{C2,B3,C4}
		\stones[w]{D4,E3,D2,C3}
	\end{goban}
\end{codeexample}

\karu\ prefiere el recorte justo, así que sigue con |B2:E4|. En un tablero parcial, el borde del tablero sobra: lo quita con |line = false| en el |background|.

\begin{codeexample}[]
	\begin{goban}[partial = B2:E4, background = {line = false}]
		\stones[b]{C2,B3,C4}
		\stones[w]{D4,E3,D2,C3}
	\end{goban}
\end{codeexample}

Las marcas se ponen con |\mark|, o con |\marks| para varias intersecciones; hay muchas formas (pág.~\pageref{sec:marcas}). A \karu\ le alcanzan dos: |circle|, el círculo, y |custom|, que pone cualquier texto, como «!». Una marca también puede ir en las opciones de la piedra, con |mark|.

\begin{codeexample}[]
	\begin{goban}[partial = B2:E4, background = {line = false}]
		\stones[b] {C2,B3,C4}
		\stones[w] {D4,E3,D2}
		\mark[custom = {!}] {D3}
		\stone[w, mark = circle] {C3}
	\end{goban}
\end{codeexample}

Para que el círculo se distinga, \karu\ lo pinta de rojo, con |line/color| de la marca.

\begin{codeexample}[]
	\begin{goban}[partial = B2:E4, background = {line = false}]
		\stones[b] {C2,B3,C4}
		\stones[w] {D4,E3,D2}
		\mark[custom = {!}] {D3}
		\stone[w, mark = {circle, line = {color = red}}] {C3}
	\end{goban}
\end{codeexample}

\subsection{Explicar una pelea: flechas, letras y vistas}
\label{tut:pelea}

En la segunda parte del artículo, \karu\ quiere explicar una pelea. Para eso necesita señalar las jugadas posibles con letras, mostrar hacia dónde corre cada grupo y, en el texto, nombrar las piedras sin obligar al lector a buscarlas en el diagrama.

\begin{figure}[H]
	\centering
	\begin{goban}[partial = A1:G7, label = {at = {north, west}}]
		\stones[b]{C3, D3, E4, C5}
		\stones[w]{D4, D5, E5, F3}
		\marks[dim]{F3}
		\marks[triangle]{D5}
		\mark[custom = A]{E3}
		\mark[custom = B]{C4}
		\arrow{C5}{C7}
		\arrow[color = red!80!black, bend = 30, shorten = 4pt]{E4}{G5}
		\line[dash = dashed]{A6}{G6}
	\end{goban}
\end{figure}

Empieza, como siempre, por las piedras. Para no ocupar espacio, muestra solo la esquina con |partial|.

\begin{codeexample}[width=4.8cm]
	\begin{goban}[partial = A1:G7, label = {at = {north, west}}]
		\stones[b]{C3, D3, E4, C5}
		\stones[w]{D4, D5, E5, F3}
	\end{goban}
\end{codeexample}

Las jugadas candidatas se marcan con letras, mediante |custom|, que acepta cualquier texto. La marca |triangle| destaca la piedra de la que habla el texto, y |dim| apaga la piedra blanca de F3, que ya no importa.

\begin{codeexample}[width=4.8cm]
	\begin{goban}[partial = A1:G7, label = {at = {north, west}}]
		\stones[b]{C3, D3, E4, C5}
		\stones[w]{D4, D5, E5, F3}
		\marks[dim]{F3}
		\marks[triangle]{D5}
		\mark[custom = A]{E3}
		\mark[custom = B]{C4}
	\end{goban}
\end{codeexample}

Para mostrar hacia dónde corren los grupos, \karu\ usa |\arrow|, que une dos intersecciones (pág.~\pageref{subsec:flechas}). La flecha se dibuja con un halo del color del tablero, para que no se confunda con la cuadrícula. Con |bend|, la flecha es un arco; con |shorten|, no tapa las piedras de sus extremos.

\begin{codeexample}[width=4.8cm]
	\begin{goban}[partial = A1:G7, label = {at = {north, west}}]
		\stones[b]{C3, D3, E4, C5}
		\stones[w]{D4, D5, E5, F3}
		\arrow{C5}{C7}
		\arrow[color = red!80!black, bend = 30, shorten = 4pt]{E4}{G5}
	\end{goban}
\end{codeexample}

|\line| es como |\arrow|, pero sin punta, y acepta las mismas \emph{keys} de trazo, como |dash|.

\begin{codeexample}[width=4.8cm]
	\begin{goban}[partial = A1:G7, label = {at = {north, west}}]
		\stones[b]{C3, D3, E4, C5}
		\stones[w]{D4, D5, E5, F3}
		\line[dash = dashed]{A6}{G6}
	\end{goban}
\end{codeexample}

Ahora, el texto. \karu\ quiere escribir «la piedra negra de E4» y mostrarla ahí mismo. Para eso, le da un nombre al tablero con |name| y, después, |\pgfgoview| dibuja una parte de ese tablero en medio del párrafo, como quedó al final (pág.~\pageref{subsec:vistas}): un punto o un rectángulo.

\begin{codeexample}[width=4.8cm]
	\begin{goban}[partial = A1:G7, name = pelea]
		\stones[b]{C3, D3, E4, C5}
		\stones[w]{D4, D5, E5, F3}
		\marks[triangle]{D5}
		\mark[custom = A]{E3}
	\end{goban}
\end{codeexample}

\begin{codeexample}[code only]
	La negra \pgfgoview{pelea}{E4} está separada de sus compañeras.
	Si juega en A, conecta: \pgfgoview{pelea}{C3:E5}.
	Sin marcas: \pgfgoview[marks = false]{pelea}{C3:E5}.
\end{codeexample}

La negra \pgfgoview{pelea}{E4} está separada de sus compañeras.
Si juega en A, conecta: \pgfgoview{pelea}{C3:E5}.
Sin marcas: \pgfgoview[marks = false]{pelea}{C3:E5}.

La vista tiene el estilo del tablero original, pero cualquier opción de |goban| puede cambiarlo: \karu\ prefiere quitar el fondo y las coordenadas en todas las vistas del artículo.

\begin{codeexample}[code only]
	\pgfgoset{view = {background = false, label = false}}

	La negra \pgfgoview{pelea}{E4} está separada de sus compañeras.
	Si juega en A, conecta: \pgfgoview{pelea}{C3:E5}.
\end{codeexample}

\begingroup
\pgfgoset{view = {background = false, label = false}}

La negra \pgfgoview{pelea}{E4} está separada de sus compañeras.
Si juega en A, conecta: \pgfgoview{pelea}{C3:E5}.
\endgroup

A veces el texto habla de un diagrama que viene después, o de una figura que \LaTeX\ lleva a la página siguiente. Una vista también puede ir antes de su |goban|: como con |\ref|, el |goban| queda en el archivo |.aux| y hay que compilar dos veces. La primera vez, la vista es |??| (pág.~\pageref{subsec:vistas-antes}).

\begin{codeexample}[code only]
	La flecha roja del diagrama de abajo sale de la negra \pgfgoview{flechas}{E4}.

	\begin{goban}[partial = A1:G7, name = flechas]
		\stones[b]{C3, D3, E4, C5}
		\stones[w]{D4, D5, E5, F3}
		\arrow{C5}{C7}
		\arrow[color = red!80!black, bend = 30, shorten = 4pt]{E4}{G5}
	\end{goban}
\end{codeexample}

\begingroup
\pgfgoset{view = {background = false, label = false}}

La flecha roja del diagrama de abajo sale de la negra \pgfgoview{flechas}{E4}.

\begin{goban}[partial = A1:G7, name = flechas]
	\stones[b]{C3, D3, E4, C5}
	\stones[w]{D4, D5, E5, F3}
	\arrow{C5}{C7}
	\arrow[color = red!80!black, bend = 30, shorten = 4pt]{E4}{G5}
\end{goban}
\endgroup

\subsection{Un ko: capturas y notas}
\label{tut:ko}

En el primer tutorial, |captures| retiró una piedra capturada de una partida en SGF. Para explicar un ko, \karu\ necesita lo mismo con jugadas escritas a mano: ponerlas en orden y que \pgfgoname\ retire las piedras capturadas.

\begin{figure}[H]
	\centering
	\begin{goban}[7, captures, notes]
		\stones[b]{B3, C2, C4}
		\stones[w]{C3, D2, D4, E3}
		\moves{D3, G6, G5, C3, A6, A5, D3}
	\end{goban}

	\pgfgomovenotes
\end{figure}

Negras toman el ko en D3 (1), blancas amenazan en otro lado (2), negras responden (3) y blancas retoman (4). Sin más, las jugadas se dibujan unas sobre otras y ninguna piedra se va del tablero.

\begin{codeexample}[width=4.8cm]
	\begin{goban}[7]
		\stones[b]{B3, C2, C4}
		\stones[w]{C3, D2, D4, E3}
		\moves{D3, G6, G5, C3}
	\end{goban}
\end{codeexample}

Con la \emph{key} |captures| (pág.~\pageref{subsec:capturas}), el |goban| lleva la cuenta de las libertades y retira cada grupo que se queda sin ellas. La blanca de C3 desaparece con la jugada 1, y la negra 1, con la 4. |\pgfgoprisoners| dice cuántas piedras capturó cada jugador.

\begin{codeexample}[width=4.8cm]
	\begin{goban}[7, captures]
		\stones[b]{B3, C2, C4}
		\stones[w]{C3, D2, D4, E3}
		\moves{D3, G6, G5, C3}
	\end{goban}

	Negras capturaron \pgfgoprisoners{b},
	blancas \pgfgoprisoners{w}.
\end{codeexample}

Si las piedras capturadas deben seguir viéndose, |captures = fade| las deja como prisioneras.

\begin{codeexample}[width=4.8cm]
	\begin{goban}[7, captures = fade]
		\stones[b]{B3, C2, C4}
		\stones[w]{C3, D2, D4, E3}
		\moves{D3, G6, G5, C3}
	\end{goban}
\end{codeexample}

El ko sigue: negras amenazan (5), blancas responden (6) y negras retoman en D3 (7), donde ya se había jugado la 1. En los libros, esa jugada no se dibuja: debajo del diagrama se escribe «7 en 1». Eso hace la \emph{key} |notes| (pág.~\pageref{subsec:notas}), y |\pgfgomovenotes| escribe las notas.

\begin{codeexample}[width=4.8cm]
	\begin{goban}[7, captures, notes]
		\stones[b]{B3, C2, C4}
		\stones[w]{C3, D2, D4, E3}
		\moves{D3, G6, G5, C3, A6, A5, D3} % Y así hasta que alguien se canse.
	\end{goban}

	\pgfgomovenotes.
\end{codeexample}

\subsection{Una partida entera}
\label{tut:partida}

El artículo también comenta una partida larga: la primera entre Lee Sedol y AlphaGo, de 2016, guardada en |lee.sgf|, que llama |lee|. Como en el primer tutorial (pág.~\pageref{tut:primera}), no hace falta escribir sus 186 jugadas: \pgfgoname\ las lee del archivo (pág.~\pageref{subsec:sgf}).

Lo primero son los datos de la partida, que |\sgfinfo| toma de la raíz del archivo.

\begin{codeexample}[code only]
	\pgfgodefinesgf{lee}{lee.sgf}

	\sgfinfo{lee}{PB} (negras) contra \sgfinfo{lee}{PW} (blancas),
	\sgfinfo{lee}{EV}, \sgfinfo{lee}{RO}, \sgfinfo{lee}{DT}.
	Resultado: \sgfinfo{lee}{RE}.
\end{codeexample}

\sgfinfo{lee}{PB} (negras) contra \sgfinfo{lee}{PW} (blancas),
\sgfinfo{lee}{EV}, \sgfinfo{lee}{RO}, \sgfinfo{lee}{DT}.
Resultado: \sgfinfo{lee}{RE}.

|\sgfinfo| escribe el resultado tal como está en el archivo; más adelante, |\sgfdiagrams| lo escribe en palabras.

La partida es de $19\times19$: |sgf size| le da al tablero el tamaño que dice el archivo. Con |to|, \karu\ la muestra hasta la jugada 30, y con |scale|, achica el tablero.

\begin{codeexample}[]
	\begin{goban}[sgf size = lee, captures, scale = 0.6]
		\sgf[to = 30]{lee}
	\end{goban}
\end{codeexample}

Para comentar un tramo, \karu\ numera solo las jugadas desde la 21 con |from|: las anteriores se dibujan sin número.

\begin{codeexample}[]
	\begin{goban}[sgf size = lee, captures, scale = 0.6]
		\sgf[from = 21, to = 30]{lee}
	\end{goban}
\end{codeexample}

Después, \karu\ quiere la partida entera, como en los libros: un diagrama cada 50 jugadas, cada uno con su epígrafe y sus notas «N en M», y el resultado al final. Eso hace |\sgfdiagrams|, que se usa fuera de |goban|. Con |layout|, los diagramas van en una grilla de dos columnas, y con |fit| se escalan para llenarlas (pág.~\pageref{subsec:disposicion}). Aquí se muestran solo los dos últimos, con |from|.

\begin{codeexample}[code only]
	\sgfdiagrams[step = 50, from = 101, layout = {columns = 2, fit}]{lee}
\end{codeexample}

\sgfdiagrams[step = 50, from = 101, layout = {columns = 2, fit}]{lee}

El epígrafe y el resultado están en castellano porque el documento usa |babel| en castellano (pág.~\pageref{subsec:idiomas}). Algunos libros reinician la numeración en cada diagrama; con |restart|, cada diagrama, desde el segundo, numera desde 1.

\begin{codeexample}[code only]
	\sgfdiagrams[step = 50, from = 51, to = 150, restart,
		goban = {label = {at = {north, west}}}, layout = {columns = 2, fit}]{lee}
\end{codeexample}

\sgfdiagrams[step = 50, from = 51, to = 150, restart,
	goban = {label = {at = {north, west}}}, layout = {columns = 2, fit}]{lee}

\subsection{Variantes: recordar un tablero}
\label{tut:variantes}

En los comentarios, \karu\ suele mostrar qué habría pasado con otra jugada. Copiar la posición en cada diagrama es tedioso y propenso a errores: si cambia una piedra, hay que cambiarla en todos (hablo por experiencia). En la pelea, \karu\ le dio un nombre a un tablero con |name| para mostrarlo en el texto. Ese tablero guarda todo lo que dibuja, y con |resume| otro |goban| empieza con esa posición (pág.~\pageref{subsec:recordar}).

\begin{codeexample}[]
	\pgfgoset{goban = {5, partial = A1:E4}}

	\begin{goban}[name = base]
		\stones[b]{B2, C3, D2}
		\stones[w]{B3, C4, D3}
	\end{goban}

	\begin{goban}[resume = base]
		\moves[b]{C2, E3}
	\end{goban}

	\begin{goban}[resume = base]
		\moves[b]{E3, C2}
	\end{goban}
\end{codeexample}

Las variantes también pueden tener nombre, y seguir desde ellas: con |name| y |resume| en el mismo |goban|, el nuevo nombre guarda lo que trajo más lo suyo. Por defecto, |resume| trae solo las piedras; |numbers| trae también los números.

\begin{codeexample}[]
	\pgfgoset{goban = {5, partial = A1:E4}}

	\begin{goban}[name = base]
		\stones[b]{B2, C3, D2}
		\stones[w]{B3, C4, D3}
	\end{goban}

	\begin{goban}[name = variante, resume = base]
		\moves[b]{E3, C2}
	\end{goban}

	\begin{goban}[resume = {variante, numbers}]
		\moves{E4}
	\end{goban}
\end{codeexample}

Un tablero con nombre también sirve para comparar estilos. Antes de elegir el perfil de su artículo, \karu\ dibuja el mismo tablero con varios perfiles del catálogo (pág.~\pageref{sec:catalogo}). Para ponerlos uno al lado del otro, del mismo ancho, usa el entorno |pgfgolayout| (pág.~\pageref{subsec:disposicion}). Con |marks|, |resume| trae también las marcas.

\begin{codeexample}[code only]
	\begin{goban}[5, name = estilos]
		\stones[b]{B2, C3, D2}
		\stones[w]{B3, C4, D3}
		\marks[circle]{C2}
	\end{goban}

	\begin{pgfgolayout}[columns = 4, fit]
		\begin{goban}[resume = {estilos, marks}, profile = pgfgo] \end{goban}
		\begin{goban}[resume = {estilos, marks}, profile = pgfgosmooth] \end{goban}
		\begin{goban}[resume = {estilos, marks}, profile = wain] \end{goban}
		\begin{goban}[resume = {estilos, marks}, profile = simple] \end{goban}
	\end{pgfgolayout}
\end{codeexample}

\begin{goban}[5, name = estilos]
	\stones[b]{B2, C3, D2}
	\stones[w]{B3, C4, D3}
	\marks[circle]{C2}
\end{goban}

\begin{pgfgolayout}[columns = 4, fit]
	\begin{goban}[resume = {estilos, marks}, profile = pgfgo] \end{goban}
	\begin{goban}[resume = {estilos, marks}, profile = pgfgosmooth] \end{goban}
	\begin{goban}[resume = {estilos, marks}, profile = wain] \end{goban}
	\begin{goban}[resume = {estilos, marks}, profile = simple] \end{goban}
\end{pgfgolayout}

\subsection{El final del juego: territorio}
\label{tut:territorio}

\begin{figure}[H]
	\centering
	\pgfgoset{profile = wain}
	\begin{goban}[label = {at = {above, left}}]
		\stones[b] {A5,A6,B1,B2,B3,B4,B5,C1,C4,D1,D2,D5,D7,E2,E5,E7,E9,F4,F5,F6,F7,F8,F9,G5,G6,G7,H6,H8,J7}
		\stones[w] {A7,B6,B8,C2,C3,C5,C6,C7,C8,D3,D4,D6,D8,D9,E3,E4,E6,E8,F1,F2,F3,G4,H2,H3,H5,J5,J6}
		\prisoners[b]{H1,J1,J2}
		\prisoner[white] {A1}

		\marks[b, t-circle] {A1,A2,A3,A4,G8,G9,H9,J9,J8,H7}
		\marks[w, t-square] {A9,B9,C9,A8,B7,H4,J4,J3,J2,J1,H1,G1,G2,G3}
		\mark[neutral, t-square] {E1}
	\end{goban}
\end{figure}

Para terminar el artículo, \karu\ quiere mostrar cómo se cuenta una partida. Parte de la posición final. Como tiene muchas piedras y va a dibujarla varias veces, la guarda con |name|, como las variantes del tutorial anterior, y los tableros siguientes empiezan con |resume|, sin repetir las coordenadas.

\begin{codeexample}[]
	\pgfgoset{profile = wain}

	\begin{goban}[label = {at = {above, left}}, name = final]
		\stones[b] {A5,A6,B1,B2,B3,B4,B5,C1,C4,D1,D2,D5,D7,E2,E5,
			E7,E9,F4,F5,F6,F7,F8,F9,G5,G6,G7,H6,H8,J7}
		\stones[w] {A7,B6,B8,C2,C3,C5,C6,C7,C8,D3,D4,D6,D8,D9,E3,
			E4,E6,E8,F1,F2,F3,G4,H2,H3,H5,J5,J6}
	\end{goban}
\end{codeexample}

Lo siguiente son los prisioneros: las piedras muertas, que se retiran al contar. Se dibujan con |\prisoner| y |\prisoners|, o con la \emph{key} |prisoner| de |\stone| y |\stones| (pág.~\pageref{sec:piedras}). \karu\ guarda también este tablero, con otro nombre: con |name| y |resume| en el mismo |goban|, como en las variantes.

\begin{codeexample}[width=4.75cm]
	\pgfgoset{profile = wain}

	\begin{goban}[label = {at = {above, left}},
		resume = final, name = final2]
		\stones[b, prisoner]{H1, J1, J2} % Equivalente a \prisoners[b]{...}
		\prisoner[w] {A1} % Equivalente a \stone[w, prisoner]{...}
	\end{goban}
\end{codeexample}

Ahora, el territorio. Una forma es marcar cada punto con |t-square| o |t-circle|. Estas marcas toman siempre como relleno el color de relleno de la piedra a la que pertenecen y como color y grosor de borde los de la cuadrícula.

\begin{codeexample}[width=4.75cm]
	\pgfgoset{profile = wain}

	\begin{goban}[label = {at = {above, left}}, resume = final2]
		\marks[b, t-circle]{A1, A2, A3, A4, G8, G9, H9, J9, J8, H7}
		\marks[w, t-square]{A9, B9, C9, A8, B7, H4, J4, J3, J2, J1,
			H1, G1, G2, G3}
		\marks[t-square, fill = {color = red}] {E1} % No se rellena de rojo, porque toma el color de neutral
	\end{goban}
\end{codeexample}

Otra forma es |\territory| (pág.~\pageref{subsec:territorio}). Con un polígono, marca los puntos de esa región que no tienen una piedra viva; con |auto| y el kernel en Lua, busca solo el territorio de cada jugador. |\pgfgoscore| escribe el puntaje.

\begin{codeexample}[width=4.75cm]
	\pgfgoset{profile = wain}

	\begin{goban}[label = {at = {above, left}}, resume = final2]
		\territory[dame]{auto}
	\end{goban}

	Negras: \pgfgoscore{black}. Blancas: \pgfgoscore[komi = 6.5]{white}.
\end{codeexample}

Para el artículo, \karu\ prefiere que el territorio se vea como una mancha de color y no punto por punto. Con |draw = region|, el territorio de cada color es un solo polígono, cuyos lados van por las líneas del tablero.

\begin{codeexample}[width=4.75cm]
	\pgfgoset{profile = wain}

	\begin{goban}[label = {at = {above, left}}, resume = final2]
		\territory[draw = region]{auto}
	\end{goban}
\end{codeexample}

\subsection{Un cuaderno de problemas}
\label{tut:problemas}

Para sus estudiantes, \karu\ prepara un cuaderno de problemas de vida y muerte. Los problemas suelen guardarse también en SGF: la posición inicial, quién juega y un árbol con cada variante de la solución, correcta o no, con sus comentarios. |\sgfproblem| dibuja todo eso: primero el enunciado y después un diagrama por variante. Si el archivo no dice qué parte del tablero mostrar, se muestra solo la que usa el problema, como con |partial = auto|.

\begin{codeexample}[code only]
	\pgfgodefinesgf{problema}{problema.sgf}

	\sgfproblem[title = Problema 1]{problema}
\end{codeexample}

\sgfproblem[title = Problema 1]{problema}

En un cuaderno, las soluciones no van junto al enunciado. \karu\ pone los enunciados en la primera parte, con |solutions = false|, y las soluciones al final, con |statement = false|.

\begin{codeexample}[code only]
	\sgfproblem[title = Problema 1, solutions = false]{problema}
	% ... al final del cuaderno:
	\sgfproblem[statement = false, layout = {columns = 4, fit}]{problema}
\end{codeexample}

\sgfproblem[title = Problema 1, solutions = false]{problema}

\sgfproblem[statement = false, layout = {columns = 4, fit}]{problema}

Algunos problemas no están en ningún archivo, y \karu\ los dibuja a mano, con la misma parte del tablero que muestra |\sgfproblem|. Aquí escribe la parte, |A16:F19|: con |partial = auto|, cada tablero mostraría solo la suya, y no coincidirían. Para que tengan el mismo tamaño que los demás, los pone junto con los de |\sgfproblem| en un mismo |pgfgolayout|, que acomoda en una grilla todos los tableros que contiene.

\begin{codeexample}[code only]
	\begin{pgfgolayout}[columns = 3, fit]
		\sgfproblem[title = Problema 1, solutions = false]{problema}
		\begin{goban}[19, partial = A16:F19]
			\stones[b]{A16, B16, C16, D17, E18, E19}
			\stones[w]{A17, B17, C17, D18, D19}
		\end{goban}
		\begin{goban}[19, partial = A16:F19]
			\stones[b]{A17, B17, C17, D18, D19}
			\stones[w]{A18, B18, C18, C19}
		\end{goban}
	\end{pgfgolayout}
\end{codeexample}

\begin{pgfgolayout}[columns = 3, fit]
	\sgfproblem[title = Problema 1, solutions = false]{problema}
	\begin{goban}[19, partial = A16:F19]
		\stones[b]{A16, B16, C16, D17, E18, E19}
		\stones[w]{A17, B17, C17, D18, D19}
	\end{goban}
	\begin{goban}[19, partial = A16:F19]
		\stones[b]{A17, B17, C17, D18, D19}
		\stones[w]{A18, B18, C18, C19}
	\end{goban}
\end{pgfgolayout}

\subsection{Un tablero de madera: un perfil con imágenes}
\label{tut:madera}

Para la portada del cuaderno, \karu\ quiere un tablero que se parezca a uno de verdad, ya que el presupuesto no da para uno de kaya: madera, piedras con volumen y blancas con vetas. Los perfiles del catálogo dibujan las piedras como círculos; para esto hacen falta imágenes (PNG, JPG o PDF). \pgfgoname\ trae algunas, junto a sus archivos |.sty|, y \karu\ empieza con ellas: una foto de madera, |pgfgo-board.jpg|, una piedra negra de pizarra, |pgfgo-shell-black.png|, y cinco blancas de concha, |pgfgo-shell-white-1.png| a |pgfgo-shell-white-5.png|. Se usan por su nombre, desde cualquier documento; las imágenes propias se usan igual, como con |\includegraphics|.

\begin{figure}[H]
	\centering
	\pgfgodefineprofile[from = default]{madera}{
		goban = {background = {fill = {image = {file = pgfgo-board.jpg}}}},
		all = {stone = {shadow = {true, offset = 0.75pt, angle = -45}}},
		black = {stone = {image = {file = pgfgo-shell-black.png}}},
		white = {stone = {image = {file = {
			pgfgo-shell-white-1.png, pgfgo-shell-white-2.png, pgfgo-shell-white-3.png,
			pgfgo-shell-white-4.png, pgfgo-shell-white-5.png}}}}}
	\begin{goban}[profile = madera, label = {at = {north, west}}]
		\stones[b]{B2,E2,J2,C3,D3,G3,H3,C5,D5,E5,F5,G5,E6,G6}
		\stones[w]{D8,F8,E7,H7,C6,D6,F6,H6,B5,J5,B4,E4,J4,J3,H2}
	\end{goban}
\end{figure}

Parte del primer tablero del artículo. La madera va en el fondo del tablero, con |background/fill/image| (pág.~\pageref{sec:tablero}). Por defecto, la imagen cubre todo el tablero, recortada si hace falta; con |fit|, puede estirarse o repetirse en baldosas.

\begin{codeexample}[width=4.8cm]
	\pgfgoset{goban = {background = {fill = {image = {file = pgfgo-board.jpg}}}}}

	\begin{goban}[label = {at = {north, west}}]
		\stones[b]{B2,E2,J2,C3,D3,G3,H3,C5,D5,E5,F5,G5,E6,G6}
		\stones[w]{D8,F8,E7,H7,C6,D6,F6,H6,B5,J5,B4,E4,J4,J3,H2}
	\end{goban}
\end{codeexample}

Las piedras de un jugador se dibujan con una imagen mediante |stone/image| (pág.~\pageref{subsec:piedrasimagen}), en las opciones de ese jugador. La imagen ocupa el lugar del círculo: no se dibujan el relleno ni el borde.

\begin{codeexample}[width=4.8cm]
	\pgfgoset{
		goban = {background = {fill = {image = {file = pgfgo-board.jpg}}}},
		players = {
			black = {stone = {image = {file = pgfgo-shell-black.png}}},
			white = {stone = {image = {file = pgfgo-shell-white-1.png}}}}}

	\begin{goban}[label = {at = {north, west}}]
		\stones[b]{B2,E2,J2,C3,D3,G3,H3,C5,D5,E5,F5,G5,E6,G6}
		\stones[w]{D8,F8,E7,H7,C6,D6,F6,H6,B5,J5,B4,E4,J4,J3,H2}
	\end{goban}
\end{codeexample}

Las negras se ven bien, pero las blancas son todas iguales, y en un tablero de verdad cada piedra de concha tiene sus vetas. Con varios archivos en |file|, cada piedra usa uno de ellos, elegido por su posición: una piedra en un punto muestra siempre la misma imagen, en todos los diagramas.

\begin{codeexample}[width=4.8cm]
	\pgfgoset{
		goban = {background = {fill = {image = {file = pgfgo-board.jpg}}}},
		players = {
			black = {stone = {image = {file = pgfgo-shell-black.png}}},
			white = {stone = {image = {file = {
				pgfgo-shell-white-1.png, pgfgo-shell-white-2.png, pgfgo-shell-white-3.png,
				pgfgo-shell-white-4.png, pgfgo-shell-white-5.png}}}}}}

	\begin{goban}[label = {at = {north, west}}]
		\stones[b]{B2,E2,J2,C3,D3,G3,H3,C5,D5,E5,F5,G5,E6,G6}
		\stones[w]{D8,F8,E7,H7,C6,D6,F6,H6,B5,J5,B4,E4,J4,J3,H2}
	\end{goban}
\end{codeexample}

Para que las piedras parezcan apoyadas sobre la madera, falta una sombra. |shadow| la dibuja con la forma de la imagen (si la imagen tiene transparencia, como estos PNG). Como es la misma para los dos jugadores, va en |all|. Cada archivo se guarda una sola vez en el PDF, por más piedras que lo usen.

\begin{codeexample}[width=4.8cm]
	\pgfgoset{
		goban = {background = {fill = {image = {file = pgfgo-board.jpg}}}},
		players = {
			all = {stone = {shadow = {true, offset = 0.75pt, angle = -45}}},
			black = {stone = {image = {file = pgfgo-shell-black.png}}},
			white = {stone = {image = {file = {
				pgfgo-shell-white-1.png, pgfgo-shell-white-2.png, pgfgo-shell-white-3.png,
				pgfgo-shell-white-4.png, pgfgo-shell-white-5.png}}}}}}

	\begin{goban}[label = {at = {north, west}}]
		\stones[b]{B2,E2,J2,C3,D3,G3,H3,C5,D5,E5,F5,G5,E6,G6}
		\stones[w]{D8,F8,E7,H7,C6,D6,F6,H6,B5,J5,B4,E4,J4,J3,H2}
	\end{goban}
\end{codeexample}

Con cinco imágenes, todavía puede verse alguna repetida. |rotate| gira cada piedra un ángulo distinto, elegido también por su posición. Estas piedras tienen la luz arriba a la izquierda, así que \karu\ las gira poco, entre $-15$ y $15$ grados; la sombra gira con la imagen, pero sigue cayendo hacia el mismo lado.

\begin{codeexample}[width=4.8cm]
	\pgfgoset{
		goban = {background = {fill = {image = {file = pgfgo-board.jpg}}}},
		players = {
			all = {stone = {shadow = {true, offset = 0.75pt, angle = -45}}},
			black = {stone = {image = {file = pgfgo-shell-black.png}}},
			white = {stone = {image = {file = {
				pgfgo-shell-white-1.png, pgfgo-shell-white-2.png, pgfgo-shell-white-3.png,
				pgfgo-shell-white-4.png, pgfgo-shell-white-5.png},
				rotate = -15:15}}}}}

	\begin{goban}[label = {at = {north, west}}]
		\stones[b]{B2,E2,J2,C3,D3,G3,H3,C5,D5,E5,F5,G5,E6,G6}
		\stones[w]{D8,F8,E7,H7,C6,D6,F6,H6,B5,J5,B4,E4,J4,J3,H2}
	\end{goban}
\end{codeexample}

\karu\ va a usar este estilo en todo el cuaderno, y quizás en otros documentos. En lugar de repetir el |\pgfgoset|, lo guarda como un perfil con |\pgfgodefineprofile| (pág.~\pageref{subsec:perfiles}): |goban| son las opciones del tablero; |all|, las de los tres jugadores, y |black| y |white|, las de cada uno. Con |from = default|, cada parte empieza con el perfil |default|. Después, el perfil se usa como cualquier otro, con la \emph{key} |profile|.

\begin{codeexample}[code only]
	\pgfgodefineprofile[from = default]{madera}{
		goban = {background = {fill = {image = {file = pgfgo-board.jpg}}}},
		all = {stone = {shadow = {true, offset = 0.75pt, angle = -45}}},
		black = {stone = {image = {file = pgfgo-shell-black.png}}},
		white = {stone = {image = {file = {
			pgfgo-shell-white-1.png, pgfgo-shell-white-2.png, pgfgo-shell-white-3.png,
			pgfgo-shell-white-4.png, pgfgo-shell-white-5.png}}}}}

	\pgfgoset{profile = madera}
\end{codeexample}

\pgfgodefineprofile[from = default]{madera}{
	goban = {background = {fill = {image = {file = pgfgo-board.jpg}}}},
	all = {stone = {shadow = {true, offset = 0.75pt, angle = -45}}},
	black = {stone = {image = {file = pgfgo-shell-black.png}}},
	white = {stone = {image = {file = {
		pgfgo-shell-white-1.png, pgfgo-shell-white-2.png, pgfgo-shell-white-3.png,
		pgfgo-shell-white-4.png, pgfgo-shell-white-5.png}}}}}

Un perfil casi igual, con una sombra difusa y sin el borde del tablero, viene con \pgfgoname: |shell| (pág.~\pageref{subsec:perfilshell}); con |profile = shell| no hace falta definir nada. También está |yunzi|, con piedras mates. Sus definiciones están en el catálogo.

Un perfil vale para todo lo que dibuja \pgfgoname, también para los diagramas de un SGF. Así queda la partida de Lee Sedol y AlphaGo en madera, con un diagrama cada 50 jugadas y la numeración reiniciada en cada uno.

\begin{codeexample}[code only]
	\pgfgoset{profile = madera}

	\sgfdiagrams[step = 50, restart, layout = {columns = 2, fit}]{lee}
\end{codeexample}

\begingroup
\pgfgoset{profile = madera}

\sgfdiagrams[step = 50, restart, layout = {columns = 2, fit}]{lee}
\endgroup

\subsection{Diapositivas y animaciones}
\label{tut:diapositivas}

\karu\ también da clases con diapositivas. Un |goban| funciona dentro de un |frame| de |beamer| como en cualquier otro lugar, y |\sgf| con |to| permite mostrar la partida jugada por jugada, una diapositiva por jugada, con su comentario al lado.

\begin{codeexample}[code only]
	\documentclass{beamer}
	\usepackage{pgf-go}

	\begin{document}
		\begin{frame}[t]{Una partida corta}
			\begin{columns}[T]
				\column{0.5\textwidth}
					\foreach \k in {1, ..., 10} {
						\only<\k>{
							\begin{goban}[sgf size = partida, captures, scale = 1.5]
								\sgf[to = \k]{partida}
							\end{goban}}}
				\column{0.45\textwidth}
					\foreach \k in {1, ..., 10} {
						\only<\k>{\small\sgfcomment{partida}{\k}}}
			\end{columns}
		\end{frame}
	\end{document}
\end{codeexample}

Para repasar la partida entera, en cambio, le alcanzan los diagramas de Lee Sedol y AlphaGo del artículo (pág.~\pageref{tut:partida}), uno por diapositiva. Con |per frame|, |\sgfdiagrams| pone cada diagrama en su propia diapositiva del |frame|, con su epígrafe (pág.~\pageref{subsec:disposicion}). Fuera de |beamer|, |per page| hace lo mismo con las páginas.

\begin{codeexample}[code only]
	\begin{frame}{Lee Sedol contra AlphaGo}
		\centering
		\sgfdiagrams[step = 50, goban = {scale = 0.6}, layout = {per frame}]{lee}
	\end{frame}
\end{codeexample}

Las diapositivas son otro documento, pero \karu\ no quiere copiar los tableros del artículo. En el artículo, agrega |export| al |goban| de la pelea, que queda en el archivo |pelea.pgfgo|. En las diapositivas, lo trae con |\pgfgoimport|, con su nombre y su estilo, y lo usa como si lo hubiera dibujado ahí: con |\pgfgoview| o con |resume| (pág.~\pageref{subsec:exportar}).

\begin{codeexample}[code only]
	% En el artículo:
	\begin{goban}[partial = A1:G7, name = pelea, export = pelea]
		...
	\end{goban}

	% En las diapositivas:
	\pgfgoimport{pelea}
	...
	\begin{frame}{Si negras conectan}
		\begin{goban}[resume = {pelea, marks}, partial = auto]
			\moves{E3}
		\end{goban}
	\end{frame}
\end{codeexample}

En un PDF para leer en pantalla, en cambio, prefiere una animación: |\sgfanimation| dibuja un tablero por jugada con el paquete |animate|, y sus botones permiten avanzar, retroceder o reproducir la partida (pág.~\pageref{subsec:sgf}). El perfil también vale aquí. La animación se ve en Acrobat Reader y en algunos otros visores; en los demás, solo el primer cuadro.

\begin{codeexample}[code only]
	\pgfgoset{profile = madera}

	\sgfanimation[to = 40, fps = 4, goban = {captures, scale = 0.5},
		animate = {controls, loop}]{lee}
\end{codeexample}

\begin{center}
	\pgfgoset{profile = madera}
	\sgfanimation[to = 40, fps = 4, goban = {captures, scale = 0.5},
		animate = {controls, loop}]{lee}
\end{center}

\subsection{Un diccionario de josekis}
\label{tut:josekis}

Para otra clase, \karu\ quiere explicar las respuestas a la aproximación baja al punto 4-4. Las secuencias conocidas de la apertura en una esquina, los \emph{josekis}\footnote{Aprenderlos es fácil; olvidarlos, como pide el proverbio, cuesta más.}, suelen guardarse en un SGF con forma de diccionario: un árbol con una rama por cada jugada posible, comentarios en las que importan y letras en las jugadas que pueden seguir. \karu\ arma uno pequeño, |josekis.sgf|, y lo llama |josekis|.

Primero, la posición: las dos primeras jugadas, con las letras de las respuestas. El tablero es de $19\times19$, pero todo pasa en una esquina; con |partial = auto| (pág.~\pageref{tut:parcial}), el |goban| muestra solo la parte que usa lo que dibuja, también las letras.

\begin{codeexample}[]
	\pgfgodefinesgf{josekis}{josekis.sgf}

	\begin{goban}[sgf size = josekis, partial = auto]
		\sgf[to = 2]{josekis}
	\end{goban}
\end{codeexample}

Para las respuestas, \karu\ elige la línea con |variation|. En un diccionario hay tantas ramas que es más cómodo dar las jugadas de la línea, desde la primera, que buscar su número: |variation = {Q16, R14}| (pág.~\pageref{subsec:sgf}). Con |variations|, |\sgfdiagrams| dibuja además, cada una en su diagrama, las variantes con comentarios que se apartan de esa línea, y con |from = 3|, solo las que nacen en la jugada 3 o después: las respuestas, y no la invasión en el 3-3, que es otra jugada 2. Con |partial = auto|, todos los diagramas muestran la misma parte del tablero.

\begin{codeexample}[code only]
	\sgfdiagrams[variation = {Q16, R14}, variations = {commented, from = 3},
		partial = auto, comments, layout = {columns = 3, fit}]{josekis}
\end{codeexample}

\sgfdiagrams[variation = {Q16, R14}, variations = {commented, from = 3},
	partial = auto, comments, layout = {columns = 3, fit}]{josekis}

El primer diagrama es la línea elegida, hasta el final. Las jugadas de cada variante se numeran desde 1, y el epígrafe dice en qué jugada se aparta. En un diccionario grande, cada jugada trae letras en las que pueden seguir; con |marks = false|, los diagramas no las dibujan.

\subsection{Otros idiomas}
\label{tut:idiomas}

Un colega de \karu\ quiere traducir el cuaderno al francés. Los textos que escribe \pgfgoname\ (los epígrafes, el resultado, las notas «N en M» y las anotaciones de los comentarios) siguen el idioma de |babel| o |polyglossia|, también si cambia dentro del documento: con |\selectlanguage{french}| se traducen solos. También se puede elegir el idioma con la \emph{key} |lang| (pág.~\pageref{subsec:idiomas}).

\begin{codeexample}[code only]
	\pgfgoset{lang = fr}

	\sgfdiagrams[step = 5, layout = {columns = 2}]{partida}
	\sgfproblem[solutions = false]{problema}
\end{codeexample}

\begingroup
\pgfgoset{lang = fr}

\sgfdiagrams[step = 5, layout = {columns = 2}]{partida}
\sgfproblem[solutions = false]{problema}
\endgroup

El comentario del problema, «Negras matan.», no se traduce: está escrito así en el archivo. Cada texto del paquete puede cambiarse, además, con su propia \emph{key}, en cualquier idioma: por ejemplo, el epígrafe de los diagramas, con |title|.

\begin{codeexample}[code only]
	\pgfgoset{sgf diagrams = {title = Figura}}
\end{codeexample}

\subsection{Escribir menos: la sintaxis de coordenadas}
\label{tut:abreviar}

En todos los tutoriales, \karu\ escribió las piedras una por una. Cuando son muchas, \pgfgoname\ permite abreviar las listas de |\stones|, |\marks| y |\prisoners| (pág.~\pageref{subsec:sintaxis-de-coordenadas}). Las formas más útiles son los tramos de una fila o de una columna: |C:G5| es la fila 5, de C a G, y |J3:5| es la columna J, de 3 a 5. Con ellas, el primer tablero del artículo queda así:

\begin{codeexample}[width=4.8cm]
	\pgfgoset{profile = pgfgo}

	\begin{goban}[label = {at = {north, west}}]
		\stones[b]{B2, E2, J2, C3, D3, G3, H3, C:G5, E6, G6}
		\stones[w]{D8, F8, E7, H7, C6, D6, F6, H6, B5, B4, E4, H2, J3:5}
	\end{goban}
\end{codeexample}

Para grupos más grandes, hay caminos que doblan en ángulo recto. \verb+A6 |- B5+ va de A6 hacia abajo hasta la fila 5 y después hacia la derecha hasta B5: son las piedras A6, A5 y B5. \verb+E9 -| F4+ va primero hacia la derecha y después hacia abajo. \verb+D5 -- D7+ es una línea recta. Los caminos pueden encadenarse, y con ellos la posición final del tutorial de territorio (pág.~\pageref{tut:territorio}) cabe en dos líneas:

\begin{codeexample}[width=4.75cm]
	\pgfgoset{profile = wain}

	\begin{goban}[label = {at = {above, left}}]
		\stones[b] {A6|-B5|-D1|-E2, C4, E9-|F4, D5-|G7--D7, H6, H8, J7}
		\stones[w] {E8-|C5, B:E6, A7, B8, D9, J6|-H5, H2:3, F1|-C3, C2, D:E4, G4}
		\prisoners[b]{H1-|J2}
		\prisoner[w] {A1}
		\territory[dame]{auto}
	\end{goban}
\end{codeexample}

Entre dos intersecciones que no están en la misma fila ni en la misma columna, |--| es una línea inclinada: \verb+D5 -- F7+ son D5, E6 y F7. Esas piedras se tocan solo por las esquinas; para una cadena conectada, \verb+-|-+ rellena cada escalón, primero hacia el costado y después hacia arriba o hacia abajo: \verb+D4 -|- F6+ son D4, E4, E5, F5 y F6 (\verb+|-|+ empieza por la columna). Con eso, una escalera (\emph{shicho}) se escribe con sus extremos:

\begin{codeexample}[width=4.8cm]
	\begin{goban}[label = {at = {north, west}}]
		\stones[w]{D4 -|- F6 -- G6}
		\stones[b]{C4, D3, D5 -- F7, E3 -- G5}
		\arrow{G6}{J8} % Blancas corren. No por mucho.
	\end{goban}
\end{codeexample}

Hay más: rectángulos llenos (|A1:C3|) o solo su borde (|A1 r C3|), columnas y filas enteras (|D|, |4|), todo el tablero (|*|) y, después de un |;|, piedras que se quitan de la lista. Todo está en la referencia, en la pág.~\pageref{subsec:sintaxis-de-coordenadas}.


	\part{Referencia}\label{part:referencia}
	\section{Configuración}
\label{sec:configuracion}

|\pgfgoset{...}| fija opciones para todo lo que sigue: el aspecto del tablero, de las piedras, de las jugadas, etc. Es equivalente a |\keys_set:nn{pgfgo}{...}|. También sirve como una puerta al formato directo.\footnote{Úsese con discreción\dots no como los pies de página.}

\begin{command}{\pgfgoset\marg{options}}
\end{command}

\begin{key}{/pgfgo/goban = \meta{goban options}}
	Opciones del tablero (pág.~\pageref{sec:tablero}): el tamaño, la cuadrícula, el fondo, las coordenadas, etc.
\end{key}

\begin{key}{/pgfgo/stone = \meta{stone options}}
	Opciones de las piedras de los tres jugadores (pág.~\pageref{sec:piedras}), igual que |players = {all = {stone = {...}}}|.
\end{key}

\begin{key}{/pgfgo/move = \meta{move options}}
	Opciones de las jugadas, como el formato de sus números, igual que |players = {all = {move = {...}}}|, salvo |numbers| (pág.~\pageref{subsec:movimientos}), que vale para todas las jugadas.
\end{key}

\begin{key}{/pgfgo/mark = \meta{mark options}}
	Opciones de las marcas: la forma por defecto de |\marks| (|\pgfgoset{mark = {triangle}}|) y el estilo de las marcas de los tres jugadores (|line|, |fill|, |label|, |sequence| y |scale|): |mark = {scale = 2}| es lo mismo que |players = {all = {mark = {scale = 2}}}|.
\end{key}

\begin{key}{/pgfgo/sgf animation = \meta{sgf animation options}}
	Opciones de todas las |\sgfanimation| del documento (pág.~\pageref{subsec:sgfanimacion}).
\end{key}

\begin{key}{/pgfgo/sgf diagrams = \meta{sgf diagrams options}}
	Opciones de todas las |\sgfdiagrams| del documento (pág.~\pageref{subsec:sgfdiagramas}).
\end{key}

\begin{key}{/pgfgo/sgf problem = \meta{sgf problem options}}
	Opciones de todos los |\sgfproblem| del documento (pág.~\pageref{subsec:sgfproblemas}).
\end{key}

\begin{key}{/pgfgo/sgf comment = \meta{sgf comment options}}
	Opciones de todos los |\sgfcomment| del documento (pág.~\pageref{subsec:sgfdatos}).
\end{key}

\begin{key}{/pgfgo/sgf info = \meta{sgf info options}}
	Opciones de todos los |\sgfinfo| del documento (pág.~\pageref{subsec:sgfdatos}).
\end{key}

\begin{key}{/pgfgo/sgf annotations = \meta{texts}}
	Los textos de las anotaciones de los archivos SGF (pág.~\pageref{subsec:sgfdatos}).
\end{key}

\begin{key}{/pgfgo/arrow = \meta{arrow options}}
	Opciones de todas las flechas del documento (pág.~\pageref{subsec:flechas}).
\end{key}

\begin{key}{/pgfgo/line = \meta{line options}}
	Opciones de todas las líneas del documento (pág.~\pageref{subsec:flechas}).
\end{key}

\begin{key}{/pgfgo/territory = \meta{territory options}}
	Opciones de todos los |\territory| del documento, como |draw = fill| (pág.~\pageref{subsec:territorio}).
\end{key}

\begin{key}{/pgfgo/score = \meta{score options}}
	Opciones de todos los |\pgfgoscore| del documento, como |komi = 6.5| (pág.~\pageref{subsec:territorio}).
\end{key}

\begin{key}{/pgfgo/players = \meta{players options}}
	Opciones de cada jugador: el negro, el blanco y el neutral (pág.~\pageref{sec:jugadores}).
\end{key}

\begin{codeexample}[]
	\pgfgoset{
		goban = {
			background = {
				fill = {color = pgfgo},
				line = false},
			grid = {line = {color = white}}},
		players = {
			black = {
				stone = {
					line = {color = magenta},
					fill = {color = blue!50!black}}},
			white = {
				stone = {line = {color = blue}},
				mark = {line = {color = orange}}
			}				
		}
	}
	
	\begin{goban}[3]
		\stone[black] {A1}
		\stone[white, mark = {triangle}] {A2}
	\end{goban}
	
	\begin{goban}[4]
		\stone[black] {A2}
		\stones[white, mark = {circle}] {B2,A3}
	\end{goban}
\end{codeexample}

Los |goban| que siguen heredan estas opciones, y cada uno puede cambiarlas. Las opciones de un |goban| aceptan, además de las del tablero, las de |\pgfgoset|, que valen solo para ese tablero: |\begin{goban}[9, stone = {radius = 4pt}, mark = {line = {color = red}}]|.

\label{subsec:prender}Muchas \emph{keys} prenden o apagan algo: la cuadrícula, el fondo, las coordenadas, el borde y el relleno de las piedras, sus sombras y sus brillos, etc. Todas siguen la misma regla:
\begin{itemize}
	\item la \emph{key} sola, o con |true|, lo prende: |shine|, |shine = true|;
	\item con |false| lo apaga: |shine = false|;
	\item dar opciones no lo prende: |shine = {angle = 45}| cambia el ángulo, pero el brillo sigue como estaba. Para prenderlo y cambiarlo a la vez, |shine = {true, angle = 45}|.
\end{itemize}
Las excepciones son las \emph{keys} cuyo valor es el contenido mismo, como |partial| (la parte del tablero) y |sequence| (desde dónde se numera).
\subsection{El kernel de cálculo}
\label{subsec:kernel}

\pgfgoname\ separa el cálculo del dibujo. El cálculo (leer las coordenadas, saber qué parte del tablero es visible, numerar las jugadas y llevar el turno, recordar los tableros, etc.) lo hace un \emph{kernel}. El dibujo, las \emph{keys} y los perfiles son siempre los mismos. Hay dos kernels:
\begin{itemize}
	\item |lua|: escrito en Lua. Necesita Lua\LaTeX, y es el que se usa por defecto con Lua\LaTeX.
	\item |expl3|: escrito en L3. Funciona con cualquier motor, y es el que se usa por defecto con los demás.
\end{itemize}

Ambos producen exactamente los mismos diagramas\footnote{O eso dicen los tests.}. El kernel en Lua, además, sabe hacer tres cosas que el otro no: las capturas (pág.~\pageref{subsec:capturas}), el cálculo del territorio con |\territory{auto}| (pág.~\pageref{subsec:territorio}) y la lectura de partidas en SGF (pág.~\pageref{subsec:sgf}). Este manual está compilado con el kernel en Lua.

\begin{key}{/pgfgo/package/kernel = \meta{expl3 o lua} (initially lua con Lua\LaTeX, expl3 con los demás motores)}
	Opción del paquete que elige el kernel. Solo hace falta para usar |expl3| con Lua\LaTeX. Solo puede darse al cargar \pgfgoname; no se puede cambiar con |\pgfgoset|. Si se pide |kernel = lua| con pdf\LaTeX, el paquete avisa y usa |expl3|. Con |expl3|, |captures|, |\territory{auto}| y |\sgf| producen un error.

	\begin{codeexample}[code only]
		\usepackage[kernel = expl3]{pgf-go}
	\end{codeexample}
\end{key}

\subsection{Idiomas}
\label{subsec:idiomas}

Algunos textos los escribe el propio paquete: los epígrafes de |\sgfdiagrams| y |\sgfproblem|, el resultado de la partida, las anotaciones de los comentarios (|/pgfgo/sgf annotations|) y la palabra de las notas «N en M». Están traducidos al inglés (|en|), español (|es|), francés (|fr|), alemán (|de|), portugués (|pt|), italiano (|it|), esperanto (|eo|), japonés (|ja|), chino (|zh|) y coreano (|ko|). Para el japonés, el chino y el coreano hace falta una fuente con esos caracteres.

Por defecto, los textos están en el idioma que fijan babel o polyglossia, y lo siguen si cambia dentro del documento (|\selectlanguage|, |otherlanguage|). Sin babel ni polyglossia, o en un idioma sin traducción, están en inglés. Cada texto puede cambiarse con su \emph{key} (por ejemplo, |/pgfgo/sgf diagrams/title|); lo que se fija así vale en cualquier idioma.

\begin{key}{/pgfgo/lang = \meta{idioma} (initially auto)}
	Idioma de los textos: |auto| (el de babel o polyglossia), el código de un idioma (|fr|) o el nombre que le dan babel o polyglossia (|french|). Si el idioma no se conoce, el paquete avisa y usa el inglés.

	\begin{codeexample}[code only]
		\pgfgoset{lang = fr}
	\end{codeexample}
\end{key}

\begin{key}{/pgfgo/package/lang = \meta{idioma} (initially auto)}
	Lo mismo, como opción del paquete: |\usepackage[lang = fr]{pgf-go}|.
\end{key}

\begin{command}{\pgfgotext\marg{texto}}
	Escribe uno de los textos en el idioma actual. Los textos, con su valor en inglés y en español, son:

	\begin{center}
		\small
		\begin{tabular}{lll}
			|tesuji| & Tesuji & Tesuji \\
			|very good tesuji| & Very good tesuji & Muy buen tesuji \\
			|bad move| & Bad move & Mala jugada \\
			|very bad move| & Very bad move & Muy mala jugada \\
			|doubtful| & Doubtful move & Jugada dudosa \\
			|interesting| & Interesting move & Jugada interesante \\
			|good for black| & Good for Black & Bueno para Negras \\
			|very good for black| & Very good for Black & Muy bueno para Negras \\
			|good for white| & Good for White & Bueno para Blancas \\
			|very good for white| & Very good for White & Muy bueno para Blancas \\
			|even| & Even position & Posición pareja \\
			|unclear| & Unclear position & Posición incierta \\
			|diagram| & Diagram & Diagrama \\
			|variation| & Variation & Variante \\
			|black to play| & Black to play & Juegan Negras \\
			|white to play| & White to play & Juegan Blancas \\
			|at| & at & en \\
			|result| & Result & Resultado \\
			|black| & Black & Negras \\
			|white| & White & Blancas \\
			|win| & \meta{ganador} wins & Ganan \meta{ganador} \\
			|win by points| & \meta{ganador} wins by \meta{puntos} points & Ganan \meta{ganador} por \meta{puntos} puntos \\
			|win by resignation| & \meta{ganador} wins by resignation & Ganan \meta{ganador} por abandono \\
			|win on time| & \meta{ganador} wins on time & Ganan \meta{ganador} por tiempo \\
			|win by forfeit| & \meta{ganador} wins by forfeit & Ganan \meta{ganador} por incomparecencia \\
			|draw| & Draw (jigo) & Empate (jigo) \\
			|void| & No result & Sin resultado \\
			|unknown result| & Unknown result & Resultado desconocido \\
		\end{tabular}
	\end{center}

	\begin{codeexample}[]
		\pgfgotext{black to play}. \pgfgotext{tesuji}.

		\pgfgoset{lang = fr}
		\pgfgotext{black to play}.
	\end{codeexample}
\end{command}

\begin{command}{\pgfgodefinelanguage\marg{idioma}\marg{textos}}
	Agrega un idioma, o cambia textos de uno que ya existe. Los textos que falten se toman del inglés. En los textos del resultado, |\pgfgowinner| es el nombre del ganador y |\pgfgopoints| la diferencia (pág.~\pageref{subsec:sgfdiagramas}).

	\begin{codeexample}[code only]
		\pgfgodefinelanguage{es}{black = Negro, white = Blanco,
			win by points = Gana \pgfgowinner{} por \pgfgopoints{} puntos}
	\end{codeexample}
\end{command}

\begin{command}{\pgfgolanguagealias\marg{nombre}\marg{idioma}}
	Otro nombre para un idioma, por ejemplo el que le dan babel o polyglossia, si \pgfgoname\ no lo conoce.

	\begin{codeexample}[code only]
		\pgfgolanguagealias{castellano}{es}
	\end{codeexample}
\end{command}

	\section{Sintaxis de coordenadas}
\label{subsec:sintaxis-de-coordenadas}

Algo que distingue a \pgfgoname\ de otros paquetes es la posibilidad de introducir coordenadas con diversas sintaxis, ofreciendo flexibilidad para representar posiciones en el tablero de manera intuitiva. 
\begin{itemize}
	\item Posibilidad de tener múltiples entradas separando con comas.
	\item Posibilidad de eliminar piedras agregando el grupo sustractivo despu\'es de un <<;>>.
\end{itemize}

Estas formas solo funcionan en los comandos plurales (|\stones|, |\marks|, |\prisoners|); |\stone|, |\mark| y |\prisoner| reciben una sola intersección. |\moves| acepta solo las que tienen un orden claro (pág.~\pageref{subsec:movimientos}). Sin embargo, por motivos de rendimiento y consistencia, se han restringido algunas libertades que en un comienzo no estaban. Por ejemplo, para referirse a una intersección, la coordenada alfabética debe escribirse en mayúsculas y sin espacios. Así, la intersección B5 debe ingresarse como |B5|, y no como |B 5|, |B  5| o |b5|, ya que los espacios entre letra y número no son válidos en este contexto. Las restricciones se explicarán en cada caso.


Cuando \pgfgoname\ tiene que leer el argumento de las coordenadas se puede encontrar con algo como esto
\begin{center}
	|\stones{|{\color{red}|<item>, <item>, |\dots;} {\color{blue} |<item>, <item>, |\dots; }{\color{red} \dots}|}|
\end{center}

Lo primero que intentar\'a hacer es dividir el argumento entre los |;| y construir dos grupos: el grupo aditivo y el grupo sustractivo. El grupo aditivo es el encargado de a\~nadir piedras mediante nuestra sintaxis. El grupo sustractivo es el encargado de eliminarlas. Estos grupos alternan automáticamente al colocar un |;|. Esto se ve con mas detalle en la subsecci\'on correspondiente. Si no se coloca un |;|, \pgfgoname\ a\~nadir\'a todas las piedras al grupo aditivo. Se pueden a\~nadir varias piedras o varios grupos separados por |,|. En conjunto, todo esto provee una interfaz para crear grupos complejos con pocos \emph{tokens};

\begin{codeexample}[width=0pt]
	\pgfgoset{profile = pgfgosmooth}
	
	\begin{goban}[17, partial = A1:Q5, label = {at = {north, west}}]
		\stones[b]{*;B,D,H,L,Q2,P,F2:4,J2rK5,N[1:2,4];B3}
	\end{goban}
\end{codeexample}


\subsection{Grupos de piedras}

\def\sssseccion#1{\vskip0.5em\noindent\textbf{#1}\par}

\sssseccion{Todas las piedras del tablero (\emph{?`Por qu\'e no?})}

A\~nade al correspondiente grupo todas las intersecciones \textbf{visibles} del tablero. Es decir, si  es un tablero parcial, solo se a\~nadir\'an las correspondientes a la regi\'on visible. Es importante aclarar esto por si se usa en combinaci\'on con |resume|. Su sintaxis es |*|.

\begin{codeexample}[]
	\pgfgoset{goban = {5, partial = B2:D4}}
	
	\begin{goban}[label = {at = {south, west}}]
		\stones[b]{*}
	\end{goban}
\end{codeexample}


\sssseccion{Intersección}

Las intersecciones se especifican indicando la letra seguida del número, sin espacios entre el componente alfabético y el numérico: |AN|

\begin{codeexample}[]
	\pgfgoset{goban = {5, partial = B2:D4}}
	
	\begin{goban}[label = {at = {south, west}}]
		\stones[b]{C2, D3}
		\stone[w]{B4}
		\stone[w]{B 3} % No se reconoce porque tiene un espacio.
	\end{goban}
\end{codeexample}

En el ejemplo, |B 3| no se reconoce debido al espacio, mientras que |B4| se dibuja correctamente como una piedra blanca.

\sssseccion{Rectas verticales u horizontales}
Las líneas son grupos de piedras verticales u horizontales delimitadas impl\'icitamente por el tama\~no del tablero: una columna o una fila entera.

\begin{itemize}
	\item Para una columna (una línea vertical), la sintaxis es |L|: |D| es la columna D.
	\item Para una fila (una línea horizontal), la sintaxis es |N|: |4| es la fila 4.
\end{itemize}

\begin{codeexample}[width=6cm]
	\pgfgoset{goban = {5, partial = B2:D4}}
	
	\begin{goban}
		\stones[b]{B2,B3,B4}
		\stones[fill = {color = red}]{D}
	\end{goban}
	
	\begin{goban}
		\stones[b]{B2,C2,D2}
		\stones[fill = {color = blue}]{4}
	\end{goban}
	
	\begin{goban}
		%Agrego las líneas C y 3. Luego elimino C3.
		\stones[fill = {color = orange}]{C, 3; C3} %OJO C,3 != C3.
	\end{goban}
\end{codeexample}

\sssseccion{Segmentos}

Los segmentos son grupos de piedras verticales u horizontales explícitamente delimitados. Como son unidireccionales (solo verticales o solo horizontales), una de sus coordenadas permanece constante y no necesita repetirse; solo se especifica el rango de la otra variable.

\begin{itemize}
	\item Para segmentos horizontales, la sintaxis es |L:LN| o |L:L N|.
	\item Para segmentos verticales, la sintaxis es |LN:N| o |L N:N|.
\end{itemize}

\begin{codeexample}[width=6cm]
	\pgfgoset{goban = {5, partial = B2:D4}}
	
	\begin{goban}
		\stones[b]{B2,B3,B4}
		\stones[w]{C2,D2}
	\end{goban}
	
	\begin{goban}
		\stones[b]{B2:4}
		\stones[w]{C:D2}
	\end{goban}
	
	\begin{goban}
		\stones[b]{B 2:4} % Los espacios funcionan entre elementos distintos.
		\stones[w]{C:D 2}
	\end{goban}
\end{codeexample}

Aquí, |B2:4| representa un segmento vertical de B2 a B4, mientras que |C:D2| indica un segmento horizontal de C2 a D2.

\begin{codeexample}[width=6cm]
	\begin{goban}[label = {at = all}]
		\stones[fill = {color = red}]{A:E5}
		\stones[fill = {color = blue}]{H9:2}
		\stones[fill = {color = orange}]{E:A2}
		\stones[fill = {color = magenta}]{A:D7, A:C9} % Pueden combinarse.
	\end{goban}
\end{codeexample}

\sssseccion{Rectángulo: una consecuencia artificial}

La sintaxis de segmentos permite crear rangos entre coordenadas para insertar piedras. Me pareció útil extenderla para incluir dos rangos y así ofrecer una primera aproximación a los rectángulos (aunque no es mi método favorito, aquí está). En este caso, la sintaxis es |L:L N:N| o |L:LN:N|. 

\begin{codeexample}[width=6cm]
	\begin{goban}[label = {at = all}]
		\stones[fill = {color = red}]{A:E 7:8} % El espacio no es obligatorio.
		\stones[fill = {color = blue}]{I:G 9:2}
		\stones[fill = {color = orange}]{A:B 2:4, F:D 1:3}
	\end{goban}
\end{codeexample}

Por ejemplo, |A:E 7:8| crea un rectángulo que abarca las columnas de A a E en las filas 7 y 8, mientras que |I:G 9:2| desciende de I9 a G2. Aunque funcional, esta sintaxis puede resultar menos intuitiva, lo que lleva a los operadores más específicos que veremos más adelante.


\sssseccion{Productos Cartesianos}
	En este contexto, el producto cartesiano es una extensi\'on m\'as poderosa de los segmentos, pero a su vez, una generalizaci\'on a todas las notaciones antes presentadas. La idea detr\'as de esto es la siguiente: Supongamos que queremos ingresar un grupo de piedras sobre la columna A: |A1:3, A5, A7:8|. Hasta el momento no hab\'ia forma mas sencilla de escribir esto. Pero todas repiten A al comienzo: ?`qu\'e tal si en lugar de escribirla tres veces la escribimos solo una? Algo como
	\begin{center}
		|A[1:3, 5, 7:8]|
	\end{center}
	En este caso, el interprete va a distribuir el A de la izquierda para cada uno de los elementos de la derecha.
	
	Los elementos de la lista de la derecha no necesariamente tienen que estar formados por n\'umeros. Si se tiene algo como esto |A:D 5:7, A1:2, A3, A:G9| tambi\'en se puede reescribir as\'i
	\begin{center}
		|A[:D 5:7, 1:2, 3, :G9]|
	\end{center}
	Claro est\'a que nadie abstrae tanto al momento de poner coordenadas, por lo que esto es solo mero capricho. La \'unica restricci\'on para el contenido entre corchetes es que no pueden anidarse corchetes. 
	
	Del mismo modo que hay una sintaxis para letras a la izquierda, hay una para n\'umeros a la derecha, y funciona igual. La expresi\'on |A12, C12, E:G12| puede escribirse como |[A,C,E:G]12|.
	
	Existe un tercer operador, que s\'i, es m\'as un capricho que otra cosa y es el producto cartesiano sin restricciones. No me detendr\'e mucho en \'el, porque este requiere mayor grado de abstracci\'on y realmente no tiene utilidad humana. Este operador es una conjunci\'on de los dos anteriores: La expresi\'on |[A,D:F][2,5,7:9]| se interpreta como
	\begin{center}
		|A2, A5, A7:9, D:F2, D:F5, D:F 7:9|.
	\end{center}
 
	\begin{codeexample}[width=4cm]
		\begin{goban}
			\stones[b]{[A, D:F][2, 5, 7:9]}
			\marks[b]{A2, A5, A7:9, D:F2, D:F5, D:F 7:9} %<- lo que lee el interprete.
		\end{goban}
	\end{codeexample}

	La sintaxis de producto cartesiano entonces se puede invocar de tres formas distintas:
\begin{itemize}
	\item |L[itm1, itm2,...]|.
	\item |[itm1, itm2,...]N|.
	\item |[itm1,itm2,...][itm3,itm4,...]|.
\end{itemize}
donde los corchetes no pueden anidarse. 

Esto puede combinarse tambi\'en con algunos operadores que veremos a continuaci\'on para crear nodos en los caminos.

Hasta aquí hemos visto comandos para intersecciones que no pueden combinarse con otros en una misma instrucción y que, por lo tanto, \textbf{deben} separarse por comas. Ahora exploraremos comandos que sí permiten encadenarse entre sí, ofreciendo mayor flexibilidad para patrones complejos.

\sssseccion{Horizontal--Vertical y Vertical--Horizontal}

Dadas dos intersecciones no alineadas vertical ni horizontalmente, estos comandos generan las intersecciones que conectan la primera con la segunda mediante un recorrido «horizontal--vertical» o «vertical--horizontal». Esto puede recordar al operador \verb+-|+ de TikZ, de donde tomé prestada la idea. La sintaxis es \verb+LN -| LN+ (para el primer caso) o \verb+LN |- LN+ (para el segundo), con o sin espacios. Estos comandos, junto con el siguiente, pueden encadenarse para formar escaleras o recorridos propios de la m\'etrica de Manhattan.

\begin{codeexample}[width=6cm]
	\begin{goban}[label = {at = all}]
		\stones[fill = {color = red}]{A1 -| D3}
		\stones[fill = {color = blue}]{D4 -| A2}
		\stones[fill = {color = orange}] {A9 -| C7 -| G5} % Encadenados.
		\stones[fill = {color = magenta}] {F1 |- I3 -| E4} % Intente adivinar.
	\end{goban}
\end{codeexample}

En este ejemplo, \verb+A1 -| D3+ va horizontalmente de A1 a D1 y luego verticalmente a D3.

\sssseccion{Línea horizontal o vertical}

Dadas dos intersecciones alineadas vertical u horizontalmente, este operador las conecta con una línea recta de piedras. Su sintaxis es |LN -- LN|. Puede encadenarse con los dos comandos anteriores para formar patrones más elaborados. Si las intersecciones no están alineadas, el operador no dibuja nada: para eso están las líneas en cualquier dirección y las escaleras, más abajo.

\begin{codeexample}[width=6cm]
	\begin{goban}[label = {at = all}]
		\stones[fill = {color = red}]{A1 -- A5}
		\stones[fill = {color = blue}] {A1 -- D1}
		\stones[fill = {color = magenta}] {A6 -- B9} % No dibuja porque no están alineados.
		\stones[fill = {color = orange}] {I9 -| F6 -- I6} % Encadenado.
	\end{goban}
\end{codeexample}

Aquí, |A1 -- A5| dibuja una línea vertical, mientras que \verb+I9 -| F6 -- I6+ combina un recorrido «horizontal--vertical» con una línea horizontal adicional.

¡Ojo! \verb+A1 -- D1 -| F3+ es un encadenamiento permitido, mientras que \verb+A:D1 -| F3+ no lo es, porque el |:| de los rangos no es un operador de cadena.

\sssseccion{Líneas en cualquier dirección y escaleras}

Para dos intersecciones cualesquiera, |LN :--: LN| dibuja la línea de piedras más parecida a la recta que las une: una piedra por fila (o por columna, si la línea avanza más de lo que sube o baja), la más cercana a la recta. Esas piedras se tocan solo por las esquinas en cada escalón, como en una diagonal, así que en go no forman un grupo. Para que lo formen, \verb+LN :-|: LN+ y \verb+LN :|-: LN+ rellenan cada escalón con la esquina de \verb+-|+ (primero horizontal) o de \verb+|-+ (primero vertical): el resultado es una escalera conectada. Si las dos intersecciones están en una misma fila o columna, los tres operadores dan la misma línea recta.

\begin{codeexample}[width=6cm]
	\begin{goban}[label = {at = all}]
		\stones[fill = {color = red}]{A1 :--: C7}
		\stones[fill = {color = blue}]{D1 :-|: F7}
		\stones[fill = {color = orange}]{G1 :|-: I7}
	\end{goban}
\end{codeexample}

Cuando la recta pasa justo entre dos intersecciones, la piedra va del lado del extremo más cercano, y en el medio exacto, a la columna (o la fila) menor. Así, |A1 :--: C5| y |C5 :--: A1| dan las mismas piedras, y la línea queda simétrica, como la recta. En las escaleras el orden sí importa, como en \verb+-|+ y \verb+|-+.

Las intersecciones van en el orden del camino, de la primera a la segunda, con la esquina de cada escalón antes del escalón: |\moves| y |sequence| las numeran a lo largo de la línea. Los tres operadores se encadenan con los demás, como en \verb+A1 :--: E3 -| G6+; los dos puntos son parte del operador, así que no se confunden con un rango.

\begin{codeexample}[width=6cm]
	\begin{goban}[7, label = {at = {south, west}}]
		\stones[w]{A1 :-|: F4}
		\marks[sequence]{A1 :-|: F4}
	\end{goban}
\end{codeexample}

\sssseccion{Rectángulos rellenos y vacíos}

La sintaxis previa para rectángulos |L:L N:N| es poco intuitiva y más bien una herencia de los segmentos acotados. Por ello, existen dos operadores más directos para crear rectángulos: |R| y |r|. La diferencia radica en que |R| genera un rectángulo relleno, mientras que |r| solo dibuja su contorno. Dados dos vértices diametralmente opuestos del rectángulo, la sintaxis es |LN R LN| o |LN r LN|. Aunque los espacios no son obligatorios por ahora, recomiendo usarlos para mayor legibilidad.

\begin{codeexample}[width=6cm]
	\begin{goban}[label = {at = all}]
		\stones[fill = {color = red}]{A1 R D3}
		\stones[fill = {color = blue}] {A5 r D8}
		\stones[fill = {color = orange}] {I8 r G4, H6} % ¡Está vivo!
	\end{goban}
\end{codeexample}

Por ejemplo, |A1 R D3| rellena un rectángulo de A1 a D3, mientras que |A5 r D8| dibuja solo el borde de un rectángulo mayor.

Los rectángulos también son operadores de cadena y pueden usarse como tales para crear patrones más complejos.

\begin{codeexample}[width=6cm]
	\begin{goban}[label = {at = all}]
		\stones[fill = {color = orange}] {I8 r G6 r I4} % Ud. adivine.
		\marks[sequence]{I8,G6,I4}
		\stones[fill = {color = blue!20}] {A1 R E3 r A6 |- E9 -- E8} % ¡Wow!
		\marks[sequence]{A1,E3,A6,E9,E8}
	\end{goban}
\end{codeexample}

En este caso, |I8 r G6 r I4| encadena dos rectángulos de contorno, mientras que \verb+A1 R E3 r A6 |- E9 -- E8+ combina un rectángulo relleno, un contorno, un recorrido «vertical--horizontal» y una línea, mostrando la potencia del encadenamiento.

\subsection{Grupos aditivos y sustractivos}

Supongamos que queremos dibujar el siguiente diagrama:

\begin{figure}[H]
	\centering
	\begin{goban}
		\stones[b]{A1 R E6; C6 |- E3, A1:2;}
	\end{goban}
\end{figure}

Esto puede lograrse con los operadores ya vistos, dividiendo el patrón en varias partes:

\begin{codeexample}[width=6cm]
	\begin{goban}[label = {at = all}]
		\stones[fill = {color = orange}] {A6 R B3, B2 R E1, D4 R E6}
	\end{goban}
\end{codeexample}

Pero también podemos imaginar que partimos de un rectángulo completo al que le hemos quitado las piedras marcadas, ofreciendo una perspectiva más conceptual:

\begin{figure}[H]
	\centering
	\begin{goban}
		\stones[b]{A1 R E6}
		\marks[cross, line = {color = white}]{C6 |- E3, A1:2}
	\end{goban}
\end{figure}

De esto tratan los grupos aditivos y sustractivos. Un grupo aditivo añade intersecciones a la lista, mientras que uno sustractivo las elimina. Por defecto, los comandos siempre comienzan procesando grupos aditivos, lo que explica que dibujen piedras directamente. En \pgfgoname, podemos alternar entre ambos tipos usando |;|.

\begin{codeexample}[width=6cm]
	\begin{goban}[label = {at = all}, partial = A1:E6]
		\stones[b]{A1 R E6} % Grupo aditivo.
	\end{goban}
	
	\begin{goban}[label = {at = all}, partial = A1:E6]
		\stones[b]{A1 R E6; A 1:2} % Grupo aditivo; sustractivo.
	\end{goban}
	
	\begin{goban}[label = {at = all}, partial = A1:E6]
		\stones[b]{A1 R E6; A 1:2, C6 |- E3} % Grupo aditivo; sustractivo.
	\end{goban}
\end{codeexample}

En esta secuencia, primero se dibuja un rectángulo completo (|A1 R E6|), luego se elimina el segmento horizontal |A1:2|, y finalmente se quita el recorrido \verb+C6 |- E3+. Esto resulta en el mismo diagrama que el ejemplo inicial, pero con un enfoque más compacto.

La ventaja de este sistema es que permite realizar ajustes cuando sea necesario. A veces, un patrón aditivo es ideal salvo por una o dos piedras; en esos casos, podemos eliminarlas con un grupo sustractivo sin rehacer todo el diseño.

\begin{codeexample}[width=6cm]
	\begin{goban}[label = {at = all}]
		\stones[b]{A1 R E6; B:D 5, C2}
	\end{goban}
\end{codeexample}

Aquí, |B:D 5| elimina un segmento horizontal en la fila 5, y |C2| quita una piedra específica, ajustando el rectángulo inicial.

Recordemos que los grupos aditivos y sustractivos pueden encadenarse, separándolos con |;|, lo que permite alternar entre añadir y quitar elementos en una sola línea tantas veces como se desee.

\begin{codeexample}[width=6cm]
	\begin{goban}[label = {at = all}, partial = A1:E6]
		\stones[b]{A1 R E6} % AD
	\end{goban}
	
	\begin{goban}[label = {at = all}, partial = A1:E6]
		\stones[b]{A1 R E6; A1 R C4, E6} % AD, SU
	\end{goban}
	
	\begin{goban}[label = {at = all}, partial = A1:E6]
		\stones[b]{A1 R E6; A1 R C4, E6; A 1:2} % AD, SU, AD
	\end{goban}
\end{codeexample}

En este último ejemplo, se crea un rectángulo, se eliminan dos regiones (|A1 R C4| y |E6|), y luego se añade un segmento (|A 1:2|).

	\section{El tablero}
\label{sec:tablero}

Cuando hablamos de «tablero», nos referimos a un entorno o \emph{scope} en el que se renderizan tanto el goban (la cuadrícula del juego) como los elementos que contiene: piedras, etiquetas de coordenadas y otra información relevante. Este entorno está diseñado para ser flexible y personalizable, permitiendo a los usuarios adaptar su apariencia y funcionalidad según las necesidades del diagrama, ya sea para representar partidas completas, secciones específicas o incluso añadir anotaciones visuales.

\begin{environment}{{goban}\opt{\oarg{options}}}
	
	Inicializa el entorno del goban. Por defecto, dibuja un tablero cuadrado de 9×9 intersecciones, con un borde visible, fondo blanco y sin etiquetas en los laterales.
	
	\begin{codeexample}[width=5cm]
		\begin{goban}
		\end{goban}
	\end{codeexample}
	
	Dado que el tablero es la única macro que utiliza la capa de TikZ, cada intersección puede ser accedida mediante su coordenada alfanumérica (como en \pgfname) o mediante valores numéricos en el sistema de coordenadas de \tikzname. Esto permite combinar la sintaxis intuitiva del Go con las capacidades gráficas de TikZ, lo que resulta especialmente útil para añadir elementos personalizados, como flechas o marcas, directamente sobre el tablero.
	
	\begin{codeexample}[width=5.85cm]
		\begin{goban}[19, partial = {from = J1 to S9},
			label = {at = {north, east}}]
			\stones {K8,R8,Q5,R5,O4,Q4,L3,O3,M2,N2}
			\mark[black, circle] {K8}
			\stones[white] {K6,R4,R3,P3,N1,O2,Q2}
			
			\tikzset{flecha/.style = {-{latex[]}, line width = 2pt, blue}}
			\draw[flecha] (12,8)--(17,8);
		\end{goban}
	\end{codeexample}
	
	\begin{key}{/pgfgo/goban/use i=\meta{boolean} (default true, initially true)}
		
		Permite activar o desactivar el uso de la letra «I» en las coordenadas y etiquetas del tablero. Cuando está desactivada («false»), cualquier referencia a «I» se interpreta como «J».
		
		\begin{codeexample}[]
			\pgfgoset{goban = {13, label = {at = north}}}
			
			\begin{goban}[label = {at = west}, partial = {from = H1 to M3}]
				\stone[b]{J2}
			\end{goban}
			
			\begin{goban}[use i = false, partial = {from = H1 to N3}]
				\stone[b]{J2}
			\end{goban}
			
			\begin{goban}[use i = false, label = {at = east}, partial = {from = H1 to N3}]
				\stone[white, mark = {custom = !}] {I2} % Si es I, entonces pasa a J.
			\end{goban}
		\end{codeexample}
	\end{key}
	
	\begin{key}{/pgfgo/goban/size=\meta{nlines or nlines |x| nlines} (initially 9)}
		
		Define el número de líneas verticales y horizontales del tablero. Puede omitirse la palabra «size» y especificarse directamente el valor. Si se indica un único entero (como |9|), el tablero será cuadrado (9×9). Si se usa una expresión como |5x3|, se creará un tablero rectangular de 5 columnas por 3 filas.
		
		\begin{codeexample}[width=6.5cm]
			\begin{goban}[5]
			\end{goban}
			
			\begin{goban}[5x3]
			\end{goban}
		\end{codeexample}
	\end{key}
	
	\begin{key}{/pgfgo/goban/scale=\meta{float} (initially 1)}
		
		Controla la escala global del tablero, afectando tanto las dimensiones generales como el grosor de ciertos elementos (como el «line width» de las líneas).
		
		\begin{codeexample}[width=5.5cm]
			\begin{goban}[2, scale = 2]
				\stone[b, mark]{A1}
			\end{goban}
			
			\begin{goban}[2, scale = 1]
				\stone[b, mark]{A1}
			\end{goban}
			
			\begin{goban}[2, scale = 0.5]
				\stone[b, mark]{A1}
			\end{goban}
		\end{codeexample}
	\end{key}
	
	\subsection{Cuadrícula}
	
	La cuadrícula se refiere a las líneas centrales que forman las intersecciones del tablero.
	
	\begin{key}{/pgfgo/goban/grid=\meta{options}}
		Controla el comportamiento y la apariencia de las líneas centrales del tablero. Este conjunto de opciones permite modificar su visibilidad, color, grosor y separación, ofreciendo un alto grado de personalización.
	\end{key}
	
	\begin{stylekey}{/pgfgo/goban/grid/true}

		Vuelve a activar las líneas centrales después de un |grid = false| (por ejemplo, uno global hecho con |\pgfgoset|).
	\end{stylekey}

	\begin{stylekey}{/pgfgo/goban/grid/false}
		
		Desactiva las líneas centrales, dejando visibles solo el borde y las piedras.
		
		\begin{codeexample}[]
			\begin{goban}[3]
				\stone[b]{A1}
				\stone[w]{B2}
			\end{goban}
			
			\begin{goban}[3, grid = false]
				\stone[b]{A1}
				\stone[w]{B2}
			\end{goban}
		\end{codeexample}
		
		El primer tablero muestra la cuadrícula por defecto.
	\end{stylekey}
	
	\begin{key}{/pgfgo/goban/grid/line=\meta{options}}
		Permite controlar las líneas de la cuadrícula.
	\end{key}
	
	\begin{key}{/pgfgo/goban/grid/line/color =\meta{color} (initially black)}
		
		Permite cambiar el color de las líneas de la cuadrícula.
		
		\begin{codeexample}[]
			\pgfgoset{goban = {5, partial = {from = B2 to D4}}}
			
			\begin{goban}[grid = {line = {color = blue}}, new remember]
				\stones[b]{C2,D3}
				\stone[w]{B4}
			\end{goban}
			
			\begin{goban}[grid = {line = {color = red}}, resume]
			\end{goban}
		\end{codeexample}
	\end{key}
	
	\begin{key}{/pgfgo/goban/grid/line/opacity =\meta{float} (initially 1)}
		
		Establece la opacidad de las líneas de la cuadrícula, desde 0 (transparente) hasta 1 (opaco).
		
		\begin{codeexample}[]
			\pgfgoset{goban = {5, partial = {from = B2 to D4}}}
			
			\begin{goban}[new remember]
				\stones[b]{C2,D3}
				\stone[w]{B4}
			\end{goban}
			
			\begin{goban}[grid = {line = {opacity = 0.2}}, resume]
			\end{goban}
		\end{codeexample}
	\end{key}
	
	\begin{key}{/pgfgo/goban/grid/line/width =\meta{dimension} (initially 0.55pt)}
		
		Define el grosor de las líneas de la cuadrícula.
		
		\begin{codeexample}[]
			\pgfgoset{goban = {5, partial = {from = B2 to D4}}}
			
			\begin{goban}[new remember]
				\stones[b]{C2,D3}
				\stone[w]{B4}
			\end{goban}
			
			\begin{goban}[grid = {line = {width = 1.5pt}}, resume]
			\end{goban}
		\end{codeexample}
	\end{key}
	
	\begin{key}{/pgfgo/goban/grid/x sep =\meta{dimension} (initially 1.25em)}
		Establece la separación entre las líneas verticales, ajustando el ancho de las columnas del tablero.
	\end{key}
	
	\begin{key}{/pgfgo/goban/grid/y sep =\meta{dimension} (initially 1.25em)}
		Establece la separación entre las líneas horizontales, ajustando la altura de las filas.
	\end{key}
	
	\begin{key}{/pgfgo/goban/grid/sep =\meta{dimension} (initially 1.25em)}
		Define simultáneamente la separación vertical y horizontal de las líneas, asegurando un espaciado uniforme en ambas direcciones. Es una opción conveniente para mantener proporciones equilibradas.
	\end{key}
	
	\subsection{Borde}
	
	El borde del tablero es el contorno que enmarca la cuadrícula y las piedras, proporcionando un límite visual claro. Por defecto, está activo, tiene un fondo blanco y un contorno negro.
	
	\begin{key}{/pgfgo/goban/background=\meta{options} (default true)}
		Controla las propiedades del borde del tablero, como su visibilidad, color, grosor y forma, permitiendo una personalización detallada.
	\end{key}
	
	\begin{stylekey}{/pgfgo/goban/background/true}

		Vuelve a activar el fondo y el borde del tablero después de un |background = false|.
	\end{stylekey}

	\begin{stylekey}{/pgfgo/goban/background/false}
		
		Desactiva completamente el borde del tablero, eliminando tanto el contorno como el relleno. Esto es útil para diagramas minimalistas o cuando el fondo del documento ya proporciona el contraste necesario.
		
		\begin{codeexample}[]
			\pgfgoset{goban = {5, partial = {from = B2 to D4}}}
			
			\begin{goban}[new remember]
				\stones[b]{C2,D3}
				\stone[w]{B4}
			\end{goban}
			
			\begin{goban}[background = false, resume]
			\end{goban}
		\end{codeexample}
	\end{stylekey}
	
	\begin{key}{/pgfgo/goban/background/fill=\meta{options} (initially true)}
		Controla las opciones relacionadas con el relleno del fondo del tablero.
	\end{key}
	
	\begin{key}{/pgfgo/goban/background/fill=\meta{boolean} (initially true)}
		
		Determina si el relleno del fondo del tablero está activo. Si se desactiva, el área dentro del borde será transparente, mostrando solo el contorno (si está habilitado).
		
		\begin{codeexample}[]
			\pgfgoset{goban = {5, partial = {from = B2 to D4}}}
			
			\begin{goban}[background = {line = {color = red}}, new remember]
				\stones[b]{C2,D3}
				\stone[w]{B4}
			\end{goban}
			
			\begin{goban}[background = {line = false}, resume]
			\end{goban}
		\end{codeexample}
	\end{key}
	
	\begin{key}{/pgfgo/goban/background/fill/color =\meta{color} (default white)}
		
		Define el color de relleno del tablero.
		
		\begin{codeexample}[]
			\pgfgoset{goban = {5, partial = {from = B2 to D4}}}
			
			\begin{goban}[background = {fill = {color = orange!40}}, new remember]
				\stones[b]{C2,D3}
				\stone[w]{B4}
			\end{goban}
			
			\begin{goban}[background = {fill = {color = olive!30}}, resume]
			\end{goban}
		\end{codeexample}
	\end{key}
	
	\begin{key}{/pgfgo/goban/background/fill/opacity =\meta{float} (initially 1)}
		
		Controla la opacidad del relleno, desde 0 (transparente) hasta 1 (opaco).
		
		\begin{codeexample}[]
			\pgfgoset{goban = {5, partial = {from = B2 to D4}}}
			
			\begin{goban}[background = {fill = {color = orange!40}}, new remember]
				\stones[b]{C2,D3}
				\stone[w]{B4}
			\end{goban}
			
			\begin{goban}[background = {fill = {color = orange!40}, fill = {opacity = 0.5}}, resume]
			\end{goban}
		\end{codeexample}
		
	\end{key}
	
	\begin{key}{/pgfgo/goban/background/line=\meta{options} (default true)}
		Controla las opciones de la línea de contorno del fondo.
	\end{key}
	
	\begin{key}{/pgfgo/goban/background/line=\meta{boolean}}
		Activa o desactiva la línea del borde.
		
		\begin{codeexample}[]
			\pgfgoset{goban = {5, partial = {from = B2 to D4}}}
			
			\begin{goban}[background = {line = {color = red}}, new remember]
				\stones[b]{C2,D3}
				\stone[w]{B4}
			\end{goban}
			
			\begin{goban}[background = {line = false}, resume]
			\end{goban}
		\end{codeexample}
	\end{key}
	
	\begin{key}{/pgfgo/goban/background/line/color =\meta{color} (default black)}
		
		Define el color de la línea del borde, permitiendo distinguirlo del fondo o de las líneas de la cuadrícula.
		
		\begin{codeexample}[]
			\pgfgoset{goban = {5, partial = {from = B2 to D4}}}
			
			\begin{goban}[background = {line = {color = red}}, new remember]
				\stones[b]{C2,D3}
				\stone[w]{B4}
			\end{goban}
			
			\begin{goban}[background = {line = {color = blue}}, resume]
			\end{goban}
		\end{codeexample}		
	\end{key}
	
	\begin{key}{/pgfgo/goban/background/line/opacity =\meta{float} (default 1)}
		
		Ajusta la opacidad de la línea del borde.
		
		\begin{codeexample}[]
			\pgfgoset{goban = {5, partial = {from = B2 to D4}}}
			
			\begin{goban}[background = {line = {color = red}}, new remember]
				\stones[b]{C2,D3}
				\stone[w]{B4}
			\end{goban}
			
			\begin{goban}[background = {line = {color = red}, line = {opacity = 0.5}}, resume]
			\end{goban}
		\end{codeexample}
	\end{key}
	
	\begin{key}{/pgfgo/goban/background/line/width =\meta{dimension} (default 0.75pt)}
		
		Controla el grosor de la línea del borde.
		
		\begin{codeexample}[]
			\pgfgoset{goban = {5, partial = {from = B2 to D4}}}
			
			\begin{goban}[new remember]
				\stones[b]{C2,D3}
				\stone[w]{B4}
			\end{goban}
			
			\begin{goban}[background = {line = {width = 1.5pt}}, resume]
			\end{goban}
		\end{codeexample}
	\end{key}
	
	\begin{key}{/pgfgo/goban/background/rounded corners =\meta{dimension} (default 1pt)}
		
		Ajusta el redondeo de las esquinas del borde, desde 0pt (esquinas rectas) hasta valores mayores para un efecto más suave y estilizado. 
		
		\begin{codeexample}[]
			\pgfgoset{goban = {5, partial = {from = A1 to C3}}}
			
			\begin{goban}[new remember]
				\stones[b]{B1,C2}
				\stone[w]{A3}
			\end{goban}
			
			\begin{goban}[background = {rounded corners = 10pt}, resume]
			\end{goban}
		\end{codeexample}
	\end{key}
	
	
	\begin{key}{/pgfgo/goban/background/rounded corners/outer =\meta{dimension} (default 1pt)}
		
		Ajusta solamente el redondeo de las esquinas exteriores del borde. 
		
		\begin{codeexample}[]
			\pgfgoset{goban = {5, partial = {from = A1 to C3}}}
			
			\begin{goban}[background = {rounded corners = 0pt}, new remember]
				\stones[b]{B1,C2}
				\stone[w]{A3}
			\end{goban}
			
			\begin{goban}[background = {rounded corners = {outer = 10pt}}, resume]
			\end{goban}
		\end{codeexample}
	\end{key}
	
	
	\begin{key}{/pgfgo/goban/background/rounded corners/inner =\meta{dimension} (default 1pt)}
		
		Ajusta solamente el redondeo de las esquinas interiores del tablero (solo visibles en tableros parciales). 
		
		\begin{codeexample}[]
			\pgfgoset{goban = {5, partial = {from = A1 to C3}}}
			
			\begin{goban}[background = {rounded corners = 0pt}, new remember]
				\stones[b]{B1,C2}
				\stone[w]{A3}
			\end{goban}
			
			\begin{goban}[background = {rounded corners = {outer = 8pt, inner = 2pt}}, resume]
			\end{goban}
		\end{codeexample}
	\end{key}
	
	\begin{key}{/pgfgo/goban/background/sep =\meta{dimension} (initially 0.825em)}
		
		Define la separación entre la cuadrícula y el borde, basada por defecto en el radio de las piedras multiplicado por 1.5. Este ajuste asegura que las piedras cercanas al borde no se superpongan con él.
		
		\begin{codeexample}[]
			\pgfgoset{goban = {5, partial = {from = B2 to D4}}}
			
			\begin{goban}[background = {sep = 0.625em}, new remember]
				\stones[b]{C2,D3}
				\stone[w]{B4}
			\end{goban}
			
			\begin{goban}[background = {sep = 1.5em}, resume]
			\end{goban}
		\end{codeexample}
	\end{key}
	
	\subsection{Tablero parcial}
	
	En muchas ocasiones, no es necesario representar un tablero completo, sino solo una sección rectangular, ya sea por limitaciones de espacio, para enfocar la atención en una zona específica o para simplificar la creación de diagramas. El entorno |goban| incluye la \emph{key} |partial|, que permite definir esta región mediante la sub-\emph{key} |from|. Esta indica los vértices opuestos de un rectángulo, sin importar el orden en que se especifiquen, lo que facilita su uso. Por ejemplo, |from = A1 to D3| es equivalente a |from = D3 to A1|.
	
	\begin{figure}[H]
		\centering
		
		\begin{tikzpicture}
			\node[anchor=north] (X) {
				\begin{goban}[label = {at = {north, west}}]
					\coordinate (A) at (4,6);
					\coordinate (B) at (9,9);
					\coordinate (C) at (A|-B);
					\coordinate (D) at (A-|B);
					
					\tikzset{punto/.style = {circle, fill = #1, draw, inner sep=0.25em, fill opacity=0.5, anchor=center}}
					
					\node[punto = red] at (A){};
					\node[punto = red] at (B){};
					\node[punto = blue] at (C){};
					\node[punto = blue] at (D){};
					
					\draw[red, {latex}-{latex}, shorten <=5pt, shorten >= 5pt] (A)--(B);
					\draw[blue, {latex}-{latex}, shorten <=5pt, shorten >= 5pt] (C)--(D);
					
					\node[fit=(A)(B), draw, inner sep=0.6em, rounded corners, line width=1pt, fill opacity=0.1, fill=orange] {};
				\end{goban}
			};
			
			\node[anchor=north, xshift=15em] {
				\tikz{
					\node[inner sep=0] (Y) {
						\begin{goban}[label = {at = {north, west}}, partial = {from = I9 to D6}]
						\end{goban}
					};	
					\node[align = left, anchor=north, inner sep=0, yshift=-1em] at (Y.south){
						{\color{blue}|partial = {from = D9 to I6}|}\\
						{\color{blue}|partial = {from = I6 to D9}|}\\
						{\color{red} |partial = {from = D6 to I9}|}\\
						{\color{red} |partial = {from = I9 to D6}|}\\
					};
				}
			};
		\end{tikzpicture}
	\end{figure}
	
	\begin{key}{/pgfgo/goban/partial=\meta{options}}
		
		Controla las opciones para representar tableros parciales, permitiendo delimitar una región específica del goban.
	\end{key}
	
	\begin{key}{/pgfgo/goban/partial/from=\meta{intersection |to| intersection}}
		
		Especifica los vértices opuestos del rectángulo que define el tablero parcial. Estas intersecciones pueden ser cualquier par de puntos que formen una diagonal del área deseada.
		
	\end{key}
	
	\begin{codeexample}[]
		\tikzset{flecha/.style = {-{latex[]}, line width = 1pt, red}}
		
		\begin{goban}[partial = {from = A1 to D3}]
			\draw[flecha] (1,1)--(4,3); % Desde A1 hasta D3
		\end{goban}
		
		\begin{goban}[partial = {from = A9 to I6}]
			\draw[flecha] (1,9)--(9,6); % Desde A9 hasta I6
		\end{goban}
		
		\begin{goban}[partial = {from = E2 to B4}]
			\draw[flecha] (5,2)--(2,4); % Desde E2 hasta B4
		\end{goban}
	\end{codeexample}	
	
	\subsection{Etiquetas del tablero}
	
	Las etiquetas son los indicadores alfanuméricos (letras y números) que aparecen en los lados del tablero para identificar las intersecciones.
	
	\begin{key}{/pgfgo/goban/label=\meta{options}}
		Controla las propiedades de las etiquetas, como su ubicación y formato.
	\end{key}
	
	\begin{key}{/pgfgo/goban/label/false}
		Desactiva todas las etiquetas.
	\end{key}
	
	\begin{key}{/pgfgo/goban/label/at =\meta{options}}
		
		Define en qué lados del tablero aparecen las etiquetas. Las opciones incluyen combinaciones de posiciones cardinales o valores predefinidos como «all» o «none».
		
		\begin{stylekey}{/pgfgo/goban/label/at/all}
			
			Muestra etiquetas en todos los lados del tablero (norte, sur, este y oeste).
		\end{stylekey}
		
		\begin{stylekey}{/pgfgo/goban/label/at/none}
			
			Equivale a |label=false|, eliminando todas las etiquetas.
		\end{stylekey}
		
		\begin{stylekey}{/pgfgo/goban/label/at/east}
			Coloca etiquetas solo en el lado derecho (números verticales).
		\end{stylekey}
		\begin{stylekey}{/pgfgo/goban/label/at/right}
			Equivalente a |east|.
		\end{stylekey}
		
		\begin{stylekey}{/pgfgo/goban/label/at/north}
			Coloca etiquetas solo en la parte superior (letras horizontales).
		\end{stylekey}
		\begin{stylekey}{/pgfgo/goban/label/at/above}
			Equivalente a |north|.
		\end{stylekey}
		
		\begin{stylekey}{/pgfgo/goban/label/at/south}
			Coloca etiquetas solo en la parte inferior (letras horizontales).
		\end{stylekey}
		\begin{stylekey}{/pgfgo/goban/label/at/below}
			Equivalente a |south|.
		\end{stylekey}
		
		\begin{stylekey}{/pgfgo/goban/label/at/west}
			Coloca etiquetas solo en el lado izquierdo (números verticales).
		\end{stylekey}
		\begin{stylekey}{/pgfgo/goban/label/at/left}
			Equivalente a |west|.
		\end{stylekey}
		
		\begin{codeexample}[]
			\begin{goban}[12, label = {at = {all}}]
			\end{goban}
			
			\begin{goban}[partial = {from = A6 to I9}, label = {at = north}]
			\end{goban}
			
			\begin{goban}[partial = {from = B2 to E4}, label = {at = {north, east}}]
			\end{goban}
		\end{codeexample}
	\end{key}
	
	\begin{key}{/pgfgo/goban/label/font size = \meta{dimension} (initially el tamaño de la fuente actual)}

		Ajusta el tamaño de la fuente de las etiquetas.

		\begin{codeexample}[]
			\begin{goban}[5, label = {at = {south, west}}]
			\end{goban}

			\begin{goban}[5, label = {at = {south, west}, font size = 6pt}]
			\end{goban}
		\end{codeexample}
	\end{key}

	\begin{key}{/pgfgo/goban/label/font = \meta{font commands} (initially \texttt{\string\sffamily})}

		Establece la fuente de las etiquetas.

		\begin{codeexample}[]
			\begin{goban}[5, label = {at = {south, west}, font = \bfseries}]
			\end{goban}
		\end{codeexample}
	\end{key}
	
	\begin{key}{/pgfgo/goban/label/sep = \meta{dimension} (initially 0.35cm)}
		
		Controla la distancia entre las etiquetas y el borde del tablero.
	\end{key}
	
	\begin{key}{/pgfgo/goban/label/text color = \meta{color} (initially black)}
		
		Define el color del texto de las etiquetas.
	\end{key}
	
	\subsection{Capturas}
	\label{subsec:capturas}

	Con el kernel en Lua (pág.~\pageref{subsec:kernel}), \pgfgoname\ puede llevar las piedras que hay sobre el tablero y retirar las que quedan sin libertades, como en una partida.

	\begin{key}{/pgfgo/goban/captures=\meta{remove, fade o false} (default remove, initially false)}

		Activa las capturas. Cuando una jugada deja sin libertades a un grupo rival, ese grupo se retira del tablero. Si deja sin libertades a su propio grupo (un suicidio), se retira ese grupo. Las piedras colocadas con |\stone| y |\stones| también cuentan; las neutrales y los prisioneros, no. Una jugada fuera de un tablero parcial también captura. Con |kernel = expl3|, esta \emph{key} produce un error.

		\begin{codeexample}[]
			\pgfgoset{goban = {5, partial = {from = A1 to C3}}}

			\begin{goban}
				\moves*{B1,A1,A2}
			\end{goban}

			\begin{goban}[captures]
				\moves*{B1,A1,A2} % 3 captura a 2.
			\end{goban}
		\end{codeexample}

		Como \pgfname\ no puede borrar lo que ya dibujó, con |captures| las piedras se dibujan al cerrar el |goban|. Por eso, si se cambia el estilo con |\pgfgoset| dentro de un |goban| con capturas, el cambio no llega a las piedras; sí llegan las opciones de cada comando.
	\end{key}

	\begin{stylekey}{/pgfgo/goban/captures/remove}

		Las piedras capturadas desaparecen. Es lo mismo que |captures| o |captures = true|.
	\end{stylekey}

	\begin{stylekey}{/pgfgo/goban/captures/fade}

		Las piedras capturadas quedan en el tablero, dibujadas como prisioneros (con el estilo de |\prisoner| y de |/pgfgo/prisoner|).

		\begin{codeexample}[]
			\pgfgoset{goban = {5, partial = {from = A1 to E3}}}

			\begin{goban}[captures = fade]
				\stones[w]{C1,D1}
				\stones[b]{B1,C2,D2}
				\moves*{E1}
			\end{goban}
		\end{codeexample}
	\end{stylekey}

	\begin{stylekey}{/pgfgo/goban/captures/false}

		Desactiva las capturas. Es el valor inicial.
	\end{stylekey}

	\subsection{Recordar un tablero}
	El sistema de «recordar» permite almacenar y reutilizar el estado de un tablero entre diferentes entornos |goban|. Esta funcionalidad es extremadamente experimental y, en rigor, no debería estar documentada aún.
	
	\begin{stylekey}{/pgfgo/goban/new remember}
		
		Inicia un nuevo estado de memoria, guardando las piedras y configuraciones del tablero actual para usarlas más adelante.
	\end{stylekey}
	
	\begin{stylekey}{/pgfgo/goban/add to remember}
		
		Añade elementos al estado memorizado sin borrar lo anterior, permitiendo construir diagramas incrementalmente.
	\end{stylekey}
	
	\begin{stylekey}{/pgfgo/goban/forget}
		
		Borra el estado memorizado, reiniciando el sistema de memoria.
	\end{stylekey}
	
	\begin{stylekey}{/pgfgo/goban/resume}
		
		Recupera el estado memorizado y lo renderiza, útil para continuar desde un punto previo sin reescribir el código.
		
	\end{stylekey}
	
\end{environment}
	\section{Piedras, jugadas y prisioneros}
\label{sec:piedras}

Las piedras son los elementos fundamentales. En \pgfgoname, se proporcionan comandos y opciones para dibujarlas de manera flexible, permitiendo personalizar su apariencia y comportamiento según las necesidades del diagrama.

\subsection{Piedras}

\begin{command}{\stone\opt{\oarg{options}}\marg{intersection}}
	Este comando dibuja una piedra individual en la intersección especificada. El argumento obligatorio es la coordenada (ej. «C2»), mientras que las opciones opcionales permiten ajustar su estilo, como color, tamaño o marca. Si no se indica el jugador (|b|, |w| o |n|), la piedra es \emph{neutral}.
\end{command}


\begin{command}{\stones\opt{\oarg{options}}\marg{intersection \opt{, intersection, ...}}}
	Similar a |\stone|, pero permite colocar múltiples piedras con la sintaxis correspondiente (pág.~\pageref{subsec:sintaxis-de-coordenadas}). Si no se indica el jugador, las piedras también son \emph{neutrales}: para piedras negras hay que escribir |\stones[b]{...}|.
\end{command}

Las piedras, marcas y etiquetas que caen fuera de la región visible de un tablero parcial no se dibujan. Una coordenada que no existe en el tablero (por ejemplo, |E9| en un tablero de 5$\times$5) se ignora y produce una advertencia.

Es posible cambiar una opci\'on en todos los jugadores mediante |\pgfgoset{stone = {...}}|. Esto es equivalente a |\pgfgoset{players = {all = {stone = {...}}}}|.


\begin{key}{/pgfgo/stone/fill = \meta{options}}
	Controla las opciones de relleno de una piedra.
\end{key}


\begin{key}{/pgfgo/stone/fill/color = \meta{color}}	
	Define el color de relleno de una piedra.
	
	\begin{codeexample}[]
		\pgfgoset{goban = {5, partial = B2:D4}}
		
		\begin{goban}
			\stones[b]{C2,D3}
			\stone[w]{B4}
		\end{goban}
		
		\begin{goban}
			\stones[fill = {color = cyan!80!black}]{C2,D3}
			\stone[fill = {color = magenta!80!black}]{B4}
		\end{goban}
	\end{codeexample}
\end{key}

\begin{key}{/pgfgo/stone/fill/opacity = \meta{float} (initially 1)}
	
	Controla la opacidad del relleno de la piedra, desde 0 (transparente) hasta 1 (opaco).
	
	\begin{codeexample}[]
		\pgfgoset{goban = {5, partial = B2:D4}}
		
		\begin{goban}
			\stones[b]{C2,D3}
			\stone[w]{B4}
		\end{goban}
		
		\begin{goban}
			\stones[b, fill = {opacity = 0.8}]{C2,D3}
			\stone[w, fill = {opacity = 0.5}]{B4}
		\end{goban}
	\end{codeexample}
\end{key}


\begin{key}{/pgfgo/stone/line = \meta{option} (initially true)}
	Controla las opciones de la frontera de una piedra.
\end{key}

\begin{key}{/pgfgo/stone/line/color = \meta{color} (initially black)}
	
	Establece el color del borde de la piedra, que por defecto es negro para todas las piedras.
	
	\begin{codeexample}[]
		\pgfgoset{goban = {5, partial = B2:D4}}
		
		\begin{goban}
			\stones[b]{C2,D3}
			\stone[w]{B4}
		\end{goban}
		
		\begin{goban}
			\stones[b, line = {color = magenta}]{C2,D3}
			\stone[w, line = {color = cyan}]{B4}
		\end{goban}
	\end{codeexample}
\end{key}

\begin{key}{/pgfgo/stone/line/opacity = \meta{float} (initially 1)}
	
	Define la opacidad del borde de la piedra, permitiendo hacerlo más sutil o transparente.
	
	\begin{codeexample}[]
		\pgfgoset{goban = {5, partial = B2:D4}}
		
		\begin{goban}
			\stones[b]{C2,D3}
			\stone[w]{B4}
		\end{goban}
		
		\begin{goban}
			\stones[b, line = {opacity = 0.5}]{C2,D3}
			\stone[w, line = {opacity = 0.8}]{B4}
		\end{goban}
	\end{codeexample}
\end{key}

\begin{key}{/pgfgo/stone/line/width = \meta{dimension} (initially 0.7pt)}
	
	Controla el grosor del borde de la piedra.
	
	\begin{codeexample}[]
		\pgfgoset{goban = {5, partial = B2:D4}}
		
		\begin{goban}
			\stones[b]{C2,D3}
			\stone[w]{B4}
		\end{goban}
		
		\begin{goban}
			\stones[b]{C2,D3}
			\stone[w, line = {width = 2pt}]{B4}
		\end{goban}
	\end{codeexample}	
\end{key}

\begin{key}{/pgfgo/stone/radius = \meta{dimension} (initially 5.5pt)}
	
	Establece el tamaño de la piedra mediante su radio.
	
	\begin{codeexample}[]
		\pgfgoset{goban = {5, partial = B2:D4}}
		
		\begin{goban}
			\stones[b]{C2,D3}
			\stone[w]{B4}
		\end{goban}
		
		\begin{goban}
			\stones[b, radius = 0.15cm]{C2,D3}
			\stone[w, radius = 0.25cm]{B4}
		\end{goban}
	\end{codeexample}
\end{key}

\begin{key}{/pgfgo/stone/scale = \meta{float} (initially 1)}
	
	Define un factor de escala que afecta todos los elementos de la piedra (radio, borde, marcas), manteniendo sus proporciones.

	\begin{codeexample}[]
		\begin{goban}[5, partial = B2:D4]
			\stones[b, mark = circle]{C2}
			\stones[b, scale = 0.7, mark = circle]{D3}
			\stone[w, scale = 1.3, mark = triangle]{B4}
		\end{goban}
	\end{codeexample}
\end{key}

\begin{key}{/pgfgo/stone/black}
	
	Asigna a la piedra las características del primer jugador (negro por defecto), incluyendo color de relleno y borde. Sin |black| (o |b|), |\stone| y |\stones| dibujan piedras neutrales.
	
	\begin{codeexample}[]
		\pgfgoset{goban = {5, partial = B2:D4}}
		
		\begin{goban}
			\stones[b]{C2,D3}
			\stone[w]{B4}
		\end{goban}
		
		\begin{goban}
			\stones[b]{C2,D3}
			\stone[black]{B4}
		\end{goban}
	\end{codeexample}
\end{key}

\begin{key}{/pgfgo/stone/b}
	
	Abreviatura de «black», con idéntica funcionalidad.
\end{key}

\begin{key}{/pgfgo/stone/white}
	
	Asigna a la piedra las características del segundo jugador (blanco por defecto), con relleno blanco y borde negro.
	
	\begin{codeexample}[]
		\pgfgoset{goban = {5, partial = B2:D4}}
		
		\begin{goban}
			\stones[b]{C2,D3}
			\stone[white]{B4}
		\end{goban}
		
		\begin{goban}
			\stones[white]{C2,D3}
			\stone[white]{B4}
		\end{goban}
	\end{codeexample}
\end{key}

\begin{key}{/pgfgo/stone/w}
	
	Abreviatura de «white», con la misma funcionalidad.
\end{key}

\begin{key}{/pgfgo/stone/neutral}
	La piedra es del jugador neutral. Es el jugador por defecto de |\stone| y de |\stones|.

	\begin{codeexample}[]
		\begin{goban}[5, partial = B2:D4]
			\stones[b]{C2}
			\stones[w]{D3}
			\stone{B4}
			\stone[n]{C3}
		\end{goban}
	\end{codeexample}
\end{key}

\begin{key}{/pgfgo/stone/n}
	Abreviatura de «neutral».
\end{key}

\begin{key}{/pgfgo/stone/shadow = \meta{options} (default true)}
	Sombra de la piedra, dibujada debajo: |true|, |false|, |angle|, |offset|, |opacity|, |color|, |scale| e |image|. Sin |image| tiene la silueta de la piedra; con |image|, es una imagen propia (tal cual, o solo su forma con |silhouette|). Las opciones se describen en la referencia de jugadores (pág.~\pageref{subsec:jugadores}).

	\begin{codeexample}[]
		\pgfgoset{goban = {5, partial = B2:D4}}

		\begin{goban}
			\stones[b, shadow = {true, angle = -45, offset = 1.5pt}]{C2,D3}
			\stone[w, shadow = {true, angle = -45, offset = 1.5pt}]{B4}
		\end{goban}
	\end{codeexample}
\end{key}

\begin{key}{/pgfgo/stone/shine = \meta{options} (default true)}
	Brillo de la piedra, dibujado encima: |true|, |false|, |angle|, |opacity|, |color| e |image|. Las opciones se describen en la referencia de jugadores (pág.~\pageref{subsec:jugadores}).

	\begin{codeexample}[]
		\pgfgoset{goban = {5, partial = B2:D4}}

		\begin{goban}
			\stones[b, shine]{C2,D3}
			\stone[w, shine = {true, color = gray}]{B4}
		\end{goban}
	\end{codeexample}

	Con |image|, el brillo es una imagen con transparencia (PNG o PDF), centrada en la piedra (tal cual, o solo su forma, del color del brillo, con |silhouette|). Si la piedra también es una imagen, el brillo se dibuja encima y no gira con ella: las piedras pueden tener vetas giradas al azar (|rotate = random|) y la luz siempre del mismo lado.

	\begin{codeexample}[]
		\begin{goban}[5, partial = B2:D4]
			\stones[w, image = {file = pgfgo-yunzi-white-1.png, rotate = random},
				shine = {true, image = {file = shine.png}}]{B2, C2, D2, B3, C3, D3}
		\end{goban}
	\end{codeexample}
\end{key}


\label{subsec:piedrasimagen}\begin{key}{/pgfgo/stone/image = \meta{options}}
	Dibuja la piedra con una imagen (PNG, JPG o PDF) en lugar del círculo: no se dibujan el relleno ni el borde, y la imagen no tiene por qué ser redonda. Si el brillo (|shine|) está prendido, se dibuja encima de la imagen. La sombra (|shadow|) se dibuja con la forma de la imagen; para eso la imagen tiene que tener transparencia (un PNG con canal alfa o un PDF): la sombra de un JPG es un rectángulo. Un prisionero con imagen se ve translúcido, como cualquier prisionero.

	Los archivos se buscan como con |\includegraphics| (por ejemplo, en las carpetas de |\graphicspath|) y, si no están, entre las imágenes que trae \pgfgoname, en su carpeta |img| (las de los perfiles |shell| y |yunzi|, pág.~\pageref{subsec:perfilshell}). Cada archivo, con cada tamaño, se guarda una sola vez en el PDF, por más piedras que lo usen (pág.~\pageref{subsubsec:imagenespdf}). La sombra con la forma de la imagen necesita pdf\LaTeX\ o Lua\LaTeX\ (con salida PDF); con otros motores la piedra se dibuja sin sombra y con un aviso.

	Un nombre de archivo con coma va entre llaves dobles, |file = {{a,b.png}}|; con llaves simples, la coma separa dos archivos. Como con |\includegraphics|, el nombre no puede tener |&|, |#|, |%|, |$|, |{| ni |}|, ni dos espacios seguidos.

	Las imágenes SVG no se pueden usar directamente: hay que convertirlas antes a PDF, por ejemplo con |inkscape --export-type=pdf piedra.svg| o con |rsvg-convert -f pdf -o piedra.pdf piedra.svg|. Lo mismo vale para el fondo, la cuadrícula, los puntos de estrella y las marcas.

	\begin{codeexample}[]
		\pgfgoset{goban = {5, partial = B2:D4}}

		\begin{goban}
			\stones[b, image = {file = pgfgo-yunzi-black.png}, shadow = {true, angle = -45, offset = 1.5pt}]{C2,D3}
			\stone[w, image = {file = pgfgo-yunzi-white-1.png}, shadow = {true, angle = -45, offset = 1.5pt}]{B4}
		\end{goban}
	\end{codeexample}

	Para todas las piedras de un jugador, la imagen se da en sus opciones:

	\begin{codeexample}[]
		\pgfgoset{players = {
			black = {stone = {image = {file = pgfgo-yunzi-black.png}}},
			white = {stone = {image = {file = pgfgo-yunzi-white-1.png}}}}}

		\begin{goban}[5, partial = B2:D4]
			\moves[numbers = false]{C3, D3, C2, B3}
		\end{goban}
	\end{codeexample}
\end{key}

\begin{key}{/pgfgo/stone/image/file = \meta{file \opt{, file, ...}}}
	El archivo de la imagen. Con varios archivos, cada piedra usa uno, elegido por su posición: una piedra en un punto muestra siempre la misma imagen, en todos los diagramas. Sirve para variar el aspecto de las piedras, por ejemplo las vetas de las blancas.

	\begin{codeexample}[]
		\begin{goban}[5, partial = B2:D4]
			\stones[w, image = {file = {pgfgo-yunzi-white-1.png, pgfgo-yunzi-white-2.png}}]{B2, C2, D2, B3, C3, D3}
		\end{goban}
	\end{codeexample}
\end{key}

\begin{key}{/pgfgo/stone/image/width = \meta{dimension} (initially el diámetro de la piedra)}
	El ancho de la imagen; el alto mantiene la proporción. Por defecto es el diámetro de la piedra: $2 \times$ |radius| $\times$ |scale|.
\end{key}

\begin{key}{/pgfgo/stone/image/rotate = \meta{ángulo}, |random| o \meta{mínimo}|:|\meta{máximo} (initially 0)}
	Gira la imagen \meta{ángulo} grados, en sentido antihorario. Con |random|, cada piedra se gira un ángulo distinto, elegido por su posición: una piedra en un punto queda siempre igual, en todos los diagramas y en todos los cuadros de una animación. Sirve para que las vetas de una misma imagen no se vean repetidas. La sombra gira con la imagen, pero sigue cayendo hacia el mismo lado.

	\begin{codeexample}[]
		\begin{goban}[5, partial = B2:D4]
			\stones[w, image = {file = pgfgo-yunzi-white-1.png,
				rotate = random}]{B2, C2, D2, B3, C3, D3}
		\end{goban}
	\end{codeexample}

	Con un rango \meta{mínimo}|:|\meta{máximo}, el ángulo se elige igual que con |random|, pero entre \meta{mínimo} y \meta{máximo} grados (|random| es lo mismo que |0:360|). Sirve para imágenes que no se pueden girar del todo, como una piedra con luz de un lado.

	\begin{codeexample}[]
		\begin{goban}[5, partial = B2:D4]
			\stones[w, image = {file = pgfgo-yunzi-white-1.png,
				rotate = -20:20}]{B2, C2, D2, B3, C3, D3}
		\end{goban}
	\end{codeexample}
\end{key}

\begin{key}{/pgfgo/stone/image/opacity = \meta{float} (initially 1)}
	La opacidad de la imagen.

	\begin{codeexample}[]
		\begin{goban}[5, partial = B2:D4]
			\stones[b, image = {file = pgfgo-yunzi-black.png}]{C2}
			\stones[b, image = {file = pgfgo-yunzi-black.png, width = 8pt}]{D3}
			\stone[w, image = {file = pgfgo-yunzi-white-1.png, opacity = 0.5}]{B4}
		\end{goban}
	\end{codeexample}
\end{key}

\begin{stylekey}{/pgfgo/stone/image/false}
	Quita la imagen: la piedra vuelve a ser un círculo.
\end{stylekey}

\begin{key}{/pgfgo/stone/mark = \meta{mark keys}}
	
	Coloca una marca sobre la piedra (ej. triángulo, círculo). Las marcas pueden personalizarse ulteriormente con otras claves, como color o tamaño.
	
	\begin{codeexample}[]
		\pgfgoset{goban = {5, partial = B2:D4}}
		
		\begin{goban}
			\stones[b]{C2,D3}
			\stone[w]{B4}
		\end{goban}
		
		\begin{goban}
			\stones[b, mark = triangle]{C2,D3}
			\stone[w, mark = t-circle]{B4}
		\end{goban}
	\end{codeexample}
\end{key}

\begin{key}{/pgfgo/stone/forget}

	|resume| no trae esta piedra (p.~\pageref{subsec:recordar}). Queda en la memoria del |goban| con nombre: |\pgfgoview| la muestra con |forgotten|.

	\begin{codeexample}[]
		\pgfgoset{goban = {5, partial = B2:D4}}

		\begin{goban}[name = a]
			\stones[b]{C2,D3}
			\stone[w, forget]{B4}
		\end{goban}

		\begin{goban}[resume = a]
		\end{goban}
	\end{codeexample}
\end{key}

\begin{key}{/pgfgo/stone/prisoner=\meta{prisoner options} (default fill opacity = 0.5)}
	
	Caracteriza a una piedra como prisionera, típicamente con una opacidad reducida (0.5 por defecto) para indicar que ha sido capturada. Pueden añadirse otras opciones para personalizar su apariencia: se aplican después del estilo de los prisioneros (|/pgfgo/prisoner| y |/pgfgo/players/white/stone/prisoner|).
	
	\begin{codeexample}[]
		\pgfgoset{goban = {5, partial = B2:D4}}
		
		\begin{goban}
			\stones[b]{C2,D3}
			\stone[w]{B4}
		\end{goban}
		
		\begin{goban}
			\stones[b, prisoner]{C2,D3}
			\stone[white]{B4}
		\end{goban}
	\end{codeexample}
\end{key}

\subsubsection{Imágenes y tamaño del PDF}
\label{subsubsec:imagenespdf}

Con pdf\LaTeX\ o Lua\LaTeX\ (con salida PDF), cada imagen (de las piedras, el fondo, la cuadrícula, los hoshi, las marcas, el brillo o la sombra) se guarda en el PDF una sola vez, la primera vez que se usa; después, en todo el documento, cada piedra o baldosa que la usa solo la nombra. Una partida dibujada con |profile = shell| guarda la madera, las seis piedras y la sombra una vez, aunque tenga cientos de piedras repartidas en muchos diagramas. Girar una imagen (|rotate|) no agrega otra copia, y la sombra con la forma de la imagen (también con |silhouette|) usa la misma copia. Lo que sí agrega una copia es otro tamaño: la misma imagen con otro |width|, o en tableros con otro |scale| (por ejemplo, con |fit|), se guarda una vez por cada tamaño. Con otros motores (Xe\LaTeX, o salida DVI), la imagen se incluye cada vez que se usa, y el PDF crece con cada piedra.

Las piedras dibujadas, sin imagen, no se guardan así: cada una es un círculo y unas pocas instrucciones de dibujo, casi lo mismo que nombrar una copia guardada. Para piedras vectoriales más elaboradas (degradados, vetas dibujadas), conviene hacerlas en un PDF y usarlo con |stone/image|: un PDF es una imagen vectorial, se ve nítido con cualquier escala y se guarda una sola vez, como las demás imágenes.





\subsection{Movimientos}
\label{subsec:movimientos}

|\moves| representa secuencias de jugadas alternando automáticamente entre negro y blanco. Sus opciones deciden si las jugadas llevan número y personalizan las piedras asociadas.

\begin{command}{\moves\opt{\oarg{options}}\marg{intersection \opt{, intersection, ...}}}

	Dibuja piedras como parte de una secuencia de jugadas, alternando colores automáticamente (negro, blanco, negro, etc.) y con su número (|numbers|). Las jugadas son una secuencia: se juegan en el orden en que se escriben y una intersección puede repetirse (por ejemplo, para retomar un ko con |captures|). Acepta la sintaxis de |\stones| (pág.~\pageref{subsec:sintaxis-de-coordenadas}) y juega cada elemento en su orden (pág.~\pageref{subsec:orden-coordenadas}): un camino (|A1 -- A3 -- C3|) va de su primer punto al último, y un segmento, del primer extremo al segundo (|C:A 4| va de C4 a A4; |C5:1|, de C5 a C1). Lo que no tiene un recorrido fijo (rectángulos, |r|, columnas y filas enteras, |*|) se juega igual, con un aviso. Los grupos sustractivos (|;|) dan error: |\moves| solo agrega jugadas.

	\begin{codeexample}[]
		\begin{goban}[5]
			\moves{B1 -| D3, C:A 4, E5}
		\end{goban}
	\end{codeexample}
\end{command}

Dentro de un mismo |goban|, la numeración y el turno continúan de un |\moves| al siguiente. Cada |goban| nuevo empieza en la jugada 1 con negras; para continuar una partida en otro diagrama, se usan |from| y |b| o |w|, o bien |resume|.

Las opciones de |\moves| también aceptan las del estilo de la piedra (|fill|, |line|, |radius|, |shine|, |image|, etc.) y |mark|, como |\stones|. Por ejemplo, para marcar la última jugada:

\begin{codeexample}[]
	\begin{goban}[5, partial = B2:D4]
		\moves[numbers = false]{C3, D3, C2}
		\moves[numbers = false, mark = circle]{B3}
	\end{goban}
\end{codeexample}

Es posible cambiar las opciones de todos los jugadores mediante |\pgfgoset{move = {...}}| que es equivalente a |\pgfgoset{players = {all = {move = {...}}}}|, salvo |numbers|, que no es de un jugador: vale para todas las jugadas.

\begin{key}{/pgfgo/move/numbers=\meta{boolean} (default true, initially true)}

	Si las jugadas llevan su número. Sin números, las jugadas igual cuentan: la jugada siguiente sigue la numeración y el turno. Como opción de |\moves| vale para esas jugadas; con |\pgfgoset{move = {numbers = false}}| o |\begin{goban}[move = {numbers = false}]|, para todas las que siguen.

	\begin{codeexample}[]
		\begin{goban}[5, partial = B2:D4]
			\moves[numbers = false]{C3, D3}
			\moves{C2, B3}
		\end{goban}

		\begin{goban}[5, partial = B2:D4, move = {numbers = false}]
			\moves{C3, D3}
			\moves[numbers]{C2, B3}
		\end{goban}
	\end{codeexample}
\end{key}

\begin{key}{/pgfgo/move/black}
	
	Establece que la secuencia de movimientos comienza con el jugador negro, útil para forzar el inicio de una jugada en diagramas específicos.
	
	\begin{codeexample}[]
		\pgfgoset{goban = {5, partial = B2:D4}}
		
		\begin{goban}
			\stones[b]{C2,D3}
			\stone[w]{B4}
		\end{goban}
		
		\begin{goban}
			\moves[numbers = false, b]{C2,B4}
			\moves[numbers = false]{D3}
		\end{goban}
	\end{codeexample}
\end{key}

\begin{key}{/pgfgo/move/white}
	
	Establece que la secuencia comienza con el jugador blanco, alterando el orden natural de las jugadas.
	
	\begin{codeexample}[]
		\pgfgoset{goban = {5, partial = B2:D4}}
		
		\begin{goban}
			\stones[b]{C2,D3}
			\stone[w]{B4}
		\end{goban}
		
		\begin{goban}
			\moves[numbers = false, white]{C2,B4}
			\moves[numbers = false]{D3}
		\end{goban}
	\end{codeexample}
\end{key}

\begin{key}{/pgfgo/move/b}
	Abreviatura de |black|.
\end{key}

\begin{key}{/pgfgo/move/w}
	Abreviatura de |white|.
\end{key}

\begin{key}{/pgfgo/move/from = \meta{number}}
	
	Define el número inicial de una secuencia numerada de movimientos, permitiendo continuar desde un punto específico en lugar de empezar en 1.
	
	\begin{codeexample}[]
		\pgfgoset{goban = {5, partial = B2:D4}}
		
		\begin{goban}
			\stones[b]{C2,D3}
			\stone[w]{B4}
		\end{goban}
		
		\begin{goban}
			\moves[from = 7]{C2,B4}
			\moves{D3}
		\end{goban}
	\end{codeexample}
\end{key}

\begin{key}{/pgfgo/move/label = \meta{options}}
	
	Controla las propiedades de las etiquetas numéricas en los movimientos estrellados, como tamaño, color o formato.
\end{key}

\begin{key}{/pgfgo/move/label/font size=\meta{dimension} (initially automático)}

	Ajusta el tamaño de la fuente de las etiquetas, permitiendo adaptarlas al tamaño de las piedras o al diseño general. Si no se indica, el tamaño se ajusta solo según la cantidad de cifras del número.
	
	\begin{codeexample}[]
		\pgfgoset{goban = {5, partial = B2:D4}}
		
		\begin{goban}
			\stones[b]{C2,D3}
			\stone[w]{B4}
		\end{goban}
		
		\begin{goban}
			\moves[from = 7]{C2,B4}
			\moves[label = {font size = 0.5em}]{D3}
		\end{goban}
	\end{codeexample}
\end{key}

\begin{key}{/pgfgo/move/label/format = \meta{options} (initially arabic)}
	
	Define el estilo de numeración de las etiquetas.
\end{key}

\begin{key}{/pgfgo/move/label/format/arabic}
	
	Usa números arábigos, el formato por defecto para las etiquetas de movimientos.
	
	\begin{codeexample}[]
		\pgfgoset{goban = {5, partial = B2:D4}}
		
		\begin{goban}
			\stones[b]{C2,D3}
			\stone[w]{B4}
		\end{goban}
		
		\begin{goban}
			\moves[from = 7]{C2,B4}
			\moves{D3}
		\end{goban}
	\end{codeexample}
\end{key}

\begin{key}{/pgfgo/move/label/format/alph}
	
	Usa letras minúsculas para las etiquetas.
	
	\begin{codeexample}[]
		\pgfgoset{goban = {5, partial = B2:D4}}
		
		\begin{goban}
			\stones[b]{C2,D3}
			\stone[w]{B4}
		\end{goban}
		
		\pgfgoset{move={label = {format=alph}}}
		\begin{goban}
			\moves[from = 7]{C2,B4}
			\moves{D3}
		\end{goban}
	\end{codeexample}
\end{key}

\begin{key}{/pgfgo/move/label/format/Alph}
	
	Usa letras mayúsculas.
	
	\begin{codeexample}[]
		\pgfgoset{goban = {5, partial = B2:D4}}
		
		\begin{goban}
			\stones[b]{C2,D3}
			\stone[w]{B4}
		\end{goban}
		
		\pgfgoset{move = {label = {format = Alph}}}
		\begin{goban}
			\moves[from = 7]{C2,B4}
			\moves{D3}
		\end{goban}
	\end{codeexample}
\end{key}

\begin{key}{/pgfgo/move/label/format/roman}
	
	Usa números romanos minúsculos.
	
	\begin{codeexample}[]
		\pgfgoset{goban = {5, partial = B2:D4}}
		
		\begin{goban}
			\stones[b]{C2,D3}
			\stone[w]{B4}
		\end{goban}
		
		\pgfgoset{move = {label = {format = roman}}}
		\begin{goban}
			\moves{C2,B4}
			\moves{D3}
		\end{goban}
	\end{codeexample}
\end{key}

\begin{key}{/pgfgo/move/label/format/Roman}
	
	Usa números romanos mayúsculos.
	
	\begin{codeexample}[]
		\pgfgoset{goban = {5, partial = B2:D4}}
		
		\begin{goban}
			\stones[b]{C2,D3}
			\stone[w]{B4}
		\end{goban}
		
		\pgfgoset{move = {label = {format = Roman}}}
		\begin{goban}
			\moves{C2,B4}
			\moves{D3}
		\end{goban}
	\end{codeexample}
\end{key}

\begin{key}{/pgfgo/move/label/color = \meta{color}}
	
	Define el color de las etiquetas.
	
	\begin{codeexample}[]
		\pgfgoset{goban = {5, partial = B2:D4}}
		
		\begin{goban}
			\moves{C2,B4}
			\moves{D3}
		\end{goban}
		
		\begin{goban}
			\moves[label = {color = red}]{C2,B4}
			\moves[label = {color = blue}]{D3}
		\end{goban}
	\end{codeexample}
\end{key}

\begin{key}{/pgfgo/move/label/forget}

	Con |resume = {..., numbers}|, esta jugada viene sin su número (p.~\pageref{subsec:recordar}). Con |\pgfgoview|, el número solo se ve con |forgotten|.

	\begin{codeexample}[]
		\pgfgoset{goban = {5, partial = B2:D4}}

		\begin{goban}[name = a]
			\moves{C2}
			\moves[label = forget]{B4}
			\moves{D3}
		\end{goban}

		\begin{goban}[resume = {a, numbers}]
			\moves{C:D4}
		\end{goban}
	\end{codeexample}
\end{key}

\begin{key}{/pgfgo/move/forget}

	|resume| no trae estas jugadas, ni la piedra ni el número. La numeración sigue desde la última jugada que trae.
\end{key}


\subsection{Prisioneros}

\pgfgoname\ ofrece comandos específicos para representar prisioneros, generalmente con una opacidad reducida para distinguirlas de las piedras activas.

\begin{command}{\prisoner\opt{\oarg{options}}\marg{intersection}}
	
	Dibuja una piedra prisionera en una intersección específica, con opciones para personalizar su estilo.
\end{command}

\begin{command}{\prisoners\opt{\oarg{options}}\marg{intersection \opt{, intersection, ...}}}
	
	Coloca múltiples piedras prisioneras en una sola línea, eficiente para mostrar grupos capturados.
\end{command}

\begin{key}{/pgfgo/prisoner=\meta{stone options} (default fill opacity = 0.5)}
	
	Define las características de las piedras prisioneras, con una opacidad de relleno de 0.5 por defecto para indicar su estado. Pueden aplicarse otras opciones de |\stone|, como color o tamaño. Vale para los prisioneros de los tres jugadores; el estilo de cada uno (|/pgfgo/players/white/stone/prisoner|) va después.

	\begin{codeexample}[]
		\pgfgoset{goban = {5, partial = B2:D4}}

		\begin{goban}
			\stones[b]{C2}
			\prisoners[b]{D3}
			\prisoner[w]{B4}
		\end{goban}

		\pgfgoset{prisoner = {fill = {opacity = 0.2}, line = {color = red}}}

		\begin{goban}
			\stones[b]{C2}
			\prisoners[b]{D3}
			\prisoner[w]{B4}
		\end{goban}
	\end{codeexample}
\end{key}

	\section{Marcas}
\label{sec:marcas}

Las marcas son elementos gráficos que se colocan en las intersecciones del tablero para resaltar posiciones específicas, indicar puntos de interés o numerar secuencias, sin necesidad de añadir piedras. \pgfgoname\ ofrece comandos y opciones versátiles para personalizar estas marcas, desde formas predefinidas hasta etiquetas personalizadas.

\begin{command}{\mark\opt{\oarg{options}}\marg{intersection}}
	
	Dibuja una marca individual en la intersección especificada (ej. «C3»). Las opciones permiten definir su forma, color y otros atributos.
\end{command}

\begin{command}{\marks\opt{\oarg{options}}\marg{intersection \opt{, intersection, ...}}}
	
	Similar a |\mark|, pero aplica la misma marca a múltiples intersecciones mediante la sintaxis de coordenadas.
\end{command}

\begin{key}{/pgfgo/mark = \meta{options} (default cross)}
	
	Define el tipo de marca a colocar, con una cruz («cross») como valor predeterminado si no se especifica otra opción.
	
	\begin{codeexample}[]
		\pgfgoset{goban = {5, partial = {from = B2 to D4}}}
		
		\begin{goban}
			\stones[b]{C2,D3}
			\stone[w]{B4}
		\end{goban}
		
		\begin{goban}
			\stones[b]{C2,D3}
			\stone[w]{B4}
			\mark{C3}
		\end{goban}
	\end{codeexample}	
\end{key}

\begin{key}{/pgfgo/mark/circle}
	
	Coloca un círculo hueco como marca.
	
	\begin{codeexample}[]
		\pgfgoset{goban = {5, partial = {from = B2 to D4}}}
		
		\begin{goban}
			\stones[b]{C2,D3}
			\stone[w]{B4}
		\end{goban}
		
		\begin{goban}
			\stones[b]{C2,D3}
			\stone[w]{B4}
			\mark[circle]{C3}
		\end{goban}
	\end{codeexample}	
\end{key}

\begin{key}{/pgfgo/mark/triangle}
	
	Dibuja un triángulo hueco.
	
	\begin{codeexample}[]
		\pgfgoset{goban = {5, partial = {from = B2 to D4}}}
		
		\begin{goban}
			\stones[b]{C2,D3}
			\stone[w]{B4}
		\end{goban}
		
		\begin{goban}
			\stones[b]{C2,D3}
			\stone[w]{B4}
			\mark[triangle]{C3}
		\end{goban}
	\end{codeexample}	
\end{key}

\begin{key}{/pgfgo/mark/square}
	
	Coloca un cuadrado hueco.
	
	\begin{codeexample}[]
		\pgfgoset{goban = {5, partial = {from = B2 to D4}}}
		
		\begin{goban}
			\stones[b]{C2,D3}
			\stone[w]{B4}
		\end{goban}
		
		\begin{goban}
			\stones[b]{C2,D3}
			\stone[w]{B4}
			\mark[square]{C3}
		\end{goban}
	\end{codeexample}
\end{key}

\begin{key}{/pgfgo/mark/cross}
	
	Dibuja una cruz.
	
	\begin{codeexample}[]
		\pgfgoset{goban = {5, partial = {from = B2 to D4}}}
		
		\begin{goban}
			\stones[b]{C2,D3}
			\stone[w]{B4}
		\end{goban}
		
		\begin{goban}
			\stones[b]{C2,D3}
			\stone[w]{B4}
			\mark[cross]{C3}
		\end{goban}
	\end{codeexample}
\end{key}

\begin{key}{/pgfgo/mark/filled square}
	
	Coloca un cuadrado relleno.
	
	\begin{codeexample}[]
		\pgfgoset{goban = {5, partial = {from = B2 to D4}}}
		
		\begin{goban}
			\stones[b]{C2,D3}
			\stone[w]{B4}
		\end{goban}
		
		\begin{goban}
			\stones[b]{C2,D3}
			\stone[w]{B4}
			\mark[filled square]{C3}
		\end{goban}
	\end{codeexample}
\end{key}

\begin{key}{/pgfgo/mark/filled circle}
	
	Dibuja un círculo relleno.
	
	\begin{codeexample}[]
		\pgfgoset{goban = {5, partial = {from = B2 to D4}}}
		
		\begin{goban}
			\stones[b]{C2,D3}
			\stone[w]{B4}
		\end{goban}
		
		\begin{goban}
			\stones[b]{C2,D3}
			\stone[w]{B4}
			\mark[filled circle]{C3}
		\end{goban}
	\end{codeexample}
\end{key}

\begin{key}{/pgfgo/mark/t-circle}
	
	Coloca un círculo pequeño con un borde más grueso, diseñado específicamente para marcar territorios. También puede escribirse |tcircle|.
	
	\begin{codeexample}[]
		\pgfgoset{goban = {5, partial = {from = B2 to D4}}}
		
		\begin{goban}
			\stones[b]{C2,D3}
			\stone[w]{B4}
		\end{goban}
		
		\begin{goban}
			\stones[b]{C2,D3}
			\stone[w]{B4}
			\mark[t-circle]{C3}
		\end{goban}
	\end{codeexample}
\end{key}

\begin{key}{/pgfgo/mark/t-square}
	
	Dibuja un cuadrado pequeño con borde grueso, similar a «t-circle», usado para señalar territorios. También puede escribirse |tsquare|.
	
	\begin{codeexample}[]
		\pgfgoset{goban = {5, partial = {from = B2 to D4}}}
		
		\begin{goban}
			\stones[b]{C2,D3}
			\stone[w]{B4}
		\end{goban}
		
		\begin{goban}
			\stones[b]{C2,D3}
			\stone[w]{B4}
			\mark[t-square]{C3}
		\end{goban}
	\end{codeexample}
\end{key}

\begin{key}{/pgfgo/mark/hatched lines}
	
	Coloca un patrón de líneas cruzadas (tramado).
	
	\begin{codeexample}[]
		\pgfgoset{goban = {5, partial = {from = B2 to D4}}}
		
		\begin{goban}
			\stones[b]{C2,D3}
			\stone[w]{B4}
		\end{goban}
		
		\begin{goban}
			\stones[b]{C2,D3}
			\stone[w]{B4}
			\mark[hatched lines]{C3}
		\end{goban}
	\end{codeexample}
\end{key}

\begin{key}{/pgfgo/mark/sequence=\meta{options} (default from = 1)}
	Permite numerar una serie de marcas en orden.
\end{key}

\begin{key}{/pgfgo/mark/sequence/use}
	Muestra la marca |sequence| sobre el tablero. Es lo mismo que |sequence = true|.
\end{key}

\begin{key}{/pgfgo/mark/sequence/false}
	Desactiva la secuencia (por ejemplo, una activada con |\pgfgoset|).
\end{key}

\begin{key}{/pgfgo/mark/sequence/format = \meta{format}}
	Activa la secuencia y elige el formato de los números. Es lo mismo que |sequence, label = {format = ...}| (ver |label/format|).
\end{key}

\begin{key}{/pgfgo/mark/sequence/from=\meta{number}}
	
	Establece el número inicial de la secuencia de marcas numeradas, permitiendo continuar desde un punto específico en lugar de 1.
	
	\begin{codeexample}[]
		\pgfgoset{goban = {5, partial = {from = B2 to D4}}}
		
		\begin{goban}
			\stones[b]{C2,D3}
			\stone[w]{B4}
		\end{goban}
		
		\begin{goban}
			\stones[b]{C2,D3}
			\stone[w]{B4}
			\marks[sequence = {use, from = 3}]{C3,B3,C4}
		\end{goban}
	\end{codeexample}
\end{key}




\begin{key}{/pgfgo/mark/dim}
	Atenúa el punto: lo cubre con el color del fondo del tablero, translúcido, sobre la piedra y la cuadrícula. Es lo que dibuja |\sgf| para la propiedad |DD|.

	\begin{codeexample}[]
		\begin{goban}[5]
			\stones{B2, C3}\stone[w]{D4}
			\marks[dim]{A1:5, B2}
		\end{goban}
	\end{codeexample}
\end{key}

\begin{key}{/pgfgo/mark/custom=\meta{token list}}
	
	Permite definir una marca personalizada usando texto o símbolos (como «!» o «?») e incluso im\'agenes.
	
	\begin{codeexample}[]
		\pgfgoset{goban = {5, partial = {from = B2 to D4}}}
		
		\begin{goban}
			\stones[b]{C2,D3}
			\stone[w]{B4}
		\end{goban}
		
		\begin{goban}
			\stones[b]{C2,D3}
			\stone[w]{B4}
			\mark[custom = {!}]{C3}
		\end{goban}
		
		\begin{goban}
			\stones[b]{C2,D3}
			\stone[w]{B4}
			\mark[custom = {\tikz{\node[rounded corners, draw, inner sep = 1em, red]{\dots};}}]{C3}
		\end{goban}
	\end{codeexample}	
\end{key}


\begin{key}{/pgfgo/mark/line = \meta{options}}
	Controla el contorno de la marca.
\end{key}


\begin{key}{/pgfgo/mark/line/color = \meta{color}}
	Define el color del contorno de la marca.
	
	\begin{codeexample}[]
		\pgfgoset{goban = {5, partial = {from = B2 to D4}}}
		
		\begin{goban}
			\stones[b]{C2,D3}
			\stone[w]{B4}
		\end{goban}
		
		\begin{goban}
			\stones[b]{C2,D3}
			\stone[w]{B4}
			\mark[triangle, line = {color = red}]{C3}
		\end{goban}
	\end{codeexample}
\end{key}


\begin{key}{/pgfgo/mark/line/width = \meta{dimension}}
	
	Ajusta el grosor del contorno de las marcas.
	
	\begin{codeexample}[]
		\pgfgoset{goban = {5, partial = {from = B2 to D4}}}
		
		\begin{goban}
			\stones[b]{C2,D3}
			\stone[w]{B4}
			\mark[triangle]{C3}
		\end{goban}
		
		\begin{goban}
			\stones[b]{C2,D3}
			\stone[w]{B4}
			\mark[triangle, line = {width = 0.1pt}]{C3}
		\end{goban}
	\end{codeexample}
\end{key}



\begin{key}{/pgfgo/mark/fill = \meta{options}}
	Controla el relleno de las marcas rellenas (|filled square|, |filled circle| y |hatched lines|).
\end{key}

\begin{key}{/pgfgo/mark/fill/true}
	Las marcas rellenas usan |fill/color|. Sin |fill = true|, se rellenan de gris.
\end{key}

\begin{key}{/pgfgo/mark/fill/false}
	Las marcas rellenas vuelven a rellenarse de gris. Es el valor inicial.
\end{key}

\begin{key}{/pgfgo/mark/fill/color = \meta{color}}
	
	Establece el color de relleno de una marca. Solo se usa con |fill = true|.
	
	\begin{codeexample}[]
		\pgfgoset{goban = {5, partial = {from = B2 to D4}}}
		
		\begin{goban}
			\stones[b]{C2,D3}
			\stone[w]{B4}
		\end{goban}
		
		\begin{goban}
			\stones[b]{C2,D3}
			\stone[w]{B4}
			\mark[filled circle, fill = {true, color = red}]{C3}
		\end{goban}
	\end{codeexample}
\end{key}


\begin{key}{/pgfgo/mark/label = \meta{option}}
	Controla la apariencia de las etiquetas de las marcas que las tengan.
\end{key}


\begin{key}{/pgfgo/mark/label/color = \meta{color}}
	
	Controla el color del texto o símbolo en marcas personalizadas (como con |custom|).
	
	\begin{codeexample}[]
		\pgfgoset{goban = {5, partial = {from = B2 to D4}}}
		
		\begin{goban}
			\stones[b]{C2,D3}
			\stone[w]{B4}
		\end{goban}
		
		\begin{goban}
			\stones[b]{C2,D3}
			\stone[w]{B4}
			\mark[custom = {!}, label = {color = red}]{C3}
		\end{goban}
	\end{codeexample}
\end{key}

\begin{key}{/pgfgo/mark/label/font = \meta{font commands}}
	Fuente de las etiquetas de las marcas (|sequence| y |custom|).

	\begin{codeexample}[]
		\pgfgoset{goban = {5, partial = {from = B2 to D4}}}

		\begin{goban}
			\stones[b]{C2,D3}
			\stone[w]{B4}
			\marks[sequence, label = {font = \itshape}]{C3,B3,C4}
		\end{goban}
	\end{codeexample}
\end{key}

\begin{key}{/pgfgo/mark/label/format=\meta{options} (initially arabic)}
	
	Controla el estilo de numeración de las etiquetas de enumeraci\'on (como |sequence|), con opciones como números arábigos, letras o romanos.
\end{key}

\begin{key}{/pgfgo/mark/label/format/arabic}
	
	Usa números arábigos para secuencias.
	
	\begin{codeexample}[]
		\pgfgoset{goban = {5, partial = {from = B2 to D4}}}
		
		\begin{goban}
			\stones[b]{C2,D3}
			\stone[w]{B4}
		\end{goban}
		
		\begin{goban}
			\stones[b]{C2,D3}
			\stone[w]{B4}
			\marks[label = {format = arabic}, sequence]{C3,B3,C4}
		\end{goban}
	\end{codeexample}
\end{key}

\begin{key}{/pgfgo/mark/label/format/alph}
	
	Usa letras minúsculas para la secuencia.
	
	\begin{codeexample}[]
		\pgfgoset{goban = {5, partial = {from = B2 to D4}}}
		
		\begin{goban}
			\stones[b]{C2,D3}
			\stone[w]{B4}
		\end{goban}
		
		\begin{goban}
			\stones[b]{C2,D3}
			\stone[w]{B4}
			\marks[label = {format = alph}, sequence]{C3,B3,C4}
		\end{goban}
	\end{codeexample}
\end{key}

\begin{key}{/pgfgo/mark/label/format/Alph}
	
	Usa letras mayúsculas.
	
	\begin{codeexample}[]
		\pgfgoset{goban = {5, partial = {from = B2 to D4}}}
		
		\begin{goban}
			\stones[b]{C2,D3}
			\stone[w]{B4}
		\end{goban}
		
		\begin{goban}
			\stones[b]{C2,D3}
			\stone[w]{B4}
			\marks[label = {format = Alph}, sequence]{C3,B3,C4}
		\end{goban}
	\end{codeexample}
\end{key}

\begin{key}{/pgfgo/mark/label/format/roman}
	
	Usa números romanos minúsculos.
	
	\begin{codeexample}[]
		\pgfgoset{goban = {5, partial = {from = B2 to D4}}}
		
		\begin{goban}
			\stones[b]{C2,D3}
			\stone[w]{B4}
		\end{goban}
		
		\begin{goban}
			\stones[b]{C2,D3}
			\stone[w]{B4}
			\marks[label = {format = roman}, sequence]{C3,B3,C4}
		\end{goban}
	\end{codeexample}
\end{key}

\begin{key}{/pgfgo/mark/label/format/Roman}
	
	Usa números romanos mayúsculos.
	
	\begin{codeexample}[]
		\pgfgoset{goban = {5, partial = {from = B2 to D4}}}
		
		\begin{goban}
			\stones[b]{C2,D3}
			\stone[w]{B4}
		\end{goban}
		
		\begin{goban}
			\stones[b]{C2,D3}
			\stone[w]{B4}
			\marks[label = {format = Roman}, sequence]{C3,B3,C4}
		\end{goban}
	\end{codeexample}
\end{key}



\begin{key}{/pgfgo/mark/black}
	
	Aplica las propiedades del primer jugador (negro) a la marca.
	
	\begin{codeexample}[]
		\pgfgoset{goban = {5, partial = {from = B2 to D4}}}
		
		\begin{goban}
			\stones[b]{C2,D3}
			\stone[w]{B4}
			\mark[triangle]{C3}
		\end{goban}
		
		\begin{goban}
			\stones[b]{C2,D3}
			\stone[w]{B4}
			\mark[black, t-square]{C3}
		\end{goban}
	\end{codeexample}
\end{key}

\begin{key}{/pgfgo/mark/white}
	
	Aplica las propiedades del segundo jugador (blanco) a la marca.
\end{key}

\begin{key}{/pgfgo/mark/neutral}
	
	Mantiene la marca sin asociarla a ningún jugador.
\end{key}

\begin{key}{/pgfgo/mark/b}
	Abreviatura de |black|.
\end{key}

\begin{key}{/pgfgo/mark/w}
	Abreviatura de |white|.
\end{key}

\begin{key}{/pgfgo/mark/n}
	Abreviatura de |neutral|.
\end{key}

\begin{key}{/pgfgo/mark/forget}
	Esta marca no se guarda en la memoria del tablero, aunque el |goban| use |remember|.
\end{key}

\begin{key}{/pgfgo/mark/remember}
	Esta marca se guarda en la memoria del tablero, aunque el |goban| use |forget|.
\end{key}


\subsection{Flechas y líneas}
\label{subsec:flechas}

Las flechas y las líneas unen dos intersecciones. Se dibujan encima de las piedras, con un halo del color del fondo del tablero que las separa de las piedras y de la cuadrícula. Como |\mark|, |\arrow| y |\line| solo existen dentro de |goban|: afuera, |\line| sigue siendo el comando de \LaTeX. |\sgf| dibuja las flechas (|AR|) y las líneas (|LN|) del archivo.

\begin{command}{\arrow\opt{\oarg{options}}\marg{from}\marg{to}}
	Dibuja una flecha de una intersección a otra.

	\begin{codeexample}[]
		\begin{goban}[9]
			\stones{C3, D3}\stone[w]{E5}
			\arrow{B3}{F3}
			\arrow[color = red, halo = false]{C7}{G9}
		\end{goban}
	\end{codeexample}
\end{command}

\begin{command}{\line\opt{\oarg{options}}\marg{from}\marg{to}}
	Dibuja una línea de una intersección a otra.

	\begin{codeexample}[]
		\begin{goban}[9]
			\stones{C3, D4}\stone[w]{E5}
			\line{C3}{G7}
			\line[color = blue, width = 2pt]{A9}{J9}
		\end{goban}
	\end{codeexample}
\end{command}

\begin{key}{/pgfgo/arrow/color=\meta{color} (initially black)}
	Color de la flecha.
\end{key}

\begin{key}{/pgfgo/arrow/width=\meta{dimension} (initially 1.2pt)}
	Grosor de la flecha. La punta crece con él.
\end{key}

\begin{key}{/pgfgo/arrow/opacity=\meta{number} (initially 1)}
	Opacidad de la flecha.
\end{key}

\begin{key}{/pgfgo/arrow/tip=\meta{arrow tip} (initially Stealth)}
	La punta de la flecha, una de las de la biblioteca |arrows.meta| de \tikzname, como |Latex| o |Triangle|.
\end{key}

\begin{key}{/pgfgo/arrow/halo=\meta{halo options} (default true)}
	El halo alrededor de la flecha.
\end{key}

\begin{stylekey}{/pgfgo/arrow/halo/true}
	Dibuja el halo. Es el valor inicial.
\end{stylekey}

\begin{stylekey}{/pgfgo/arrow/halo/false}
	No dibuja el halo.
\end{stylekey}

\begin{key}{/pgfgo/arrow/halo/color=\meta{color}}
	Color del halo. Sin esta \emph{key}, es el del fondo del tablero.
\end{key}

\begin{key}{/pgfgo/arrow/halo/width=\meta{dimension} (initially 0.8pt)}
	Ancho del halo a cada lado de la flecha.
\end{key}

\begin{key}{/pgfgo/arrow/remember}
	Esta flecha se guarda en la memoria del tablero, aunque el |goban| use |forget|.
\end{key}

\begin{key}{/pgfgo/arrow/forget}
	Esta flecha no se guarda en la memoria del tablero, aunque el |goban| use |remember|.
\end{key}

Las líneas tienen las mismas opciones, salvo |tip|:

\begin{key}{/pgfgo/line/color=\meta{color} (initially black)}
	Como |/pgfgo/arrow/color|.
\end{key}

\begin{key}{/pgfgo/line/width=\meta{dimension} (initially 1.2pt)}
	Como |/pgfgo/arrow/width|.
\end{key}

\begin{key}{/pgfgo/line/opacity=\meta{number} (initially 1)}
	Como |/pgfgo/arrow/opacity|.
\end{key}

\begin{key}{/pgfgo/line/halo=\meta{halo options} (default true)}
	Como |/pgfgo/arrow/halo|.
\end{key}

\begin{stylekey}{/pgfgo/line/halo/true}
	Como |/pgfgo/arrow/halo/true|.
\end{stylekey}

\begin{stylekey}{/pgfgo/line/halo/false}
	Como |/pgfgo/arrow/halo/false|.
\end{stylekey}

\begin{key}{/pgfgo/line/halo/color=\meta{color}}
	Como |/pgfgo/arrow/halo/color|.
\end{key}

\begin{key}{/pgfgo/line/halo/width=\meta{dimension} (initially 0.8pt)}
	Como |/pgfgo/arrow/halo/width|.
\end{key}

\begin{key}{/pgfgo/line/remember}
	Como |/pgfgo/arrow/remember|.
\end{key}

\begin{key}{/pgfgo/line/forget}
	Como |/pgfgo/arrow/forget|.
\end{key}

	\section{El juego: capturas, notas y territorio}
\label{sec:juego}

Con el kernel en Lua, un |goban| puede seguir las reglas del juego: retirar las piedras capturadas, anotar las jugadas que caen en puntos ya jugados y contar el territorio.

	\subsection{Capturas}
	\label{subsec:capturas}

	Con el kernel en Lua (pág.~\pageref{subsec:kernel}), \pgfgoname\ puede llevar las piedras que hay sobre el tablero y retirar las que quedan sin libertades, como en una partida.

	\begin{key}{/pgfgo/goban/captures=\meta{remove, fade o false} (default remove, initially false)}

		Activa las capturas. Cuando una jugada deja sin libertades a un grupo rival, ese grupo se retira del tablero. Si deja sin libertades a su propio grupo (un suicidio), se retira ese grupo. Las piedras colocadas con |\stone| y |\stones| también cuentan; las neutrales y los prisioneros, no. Una jugada fuera de un tablero parcial también captura. Una jugada en un punto ocupado produce una advertencia, porque en una partida no es posible. Con |kernel = expl3|, esta \emph{key} produce un error.

		\begin{codeexample}[]
			\pgfgoset{goban = {5, partial = A1:C3}}

			\begin{goban}
				\moves{B1,A1,A2}
			\end{goban}

			\begin{goban}[captures]
				\moves{B1,A1,A2} % 3 captura a 2.
			\end{goban}
		\end{codeexample}

		Como \pgfname\ no puede borrar lo que ya dibujó, con |captures| las piedras se dibujan al cerrar el |goban|. Por eso, si se cambia el estilo con |\pgfgoset| dentro de un |goban| con capturas, el cambio no llega a las piedras; sí llegan las opciones de cada comando.
	\end{key}

	\begin{stylekey}{/pgfgo/goban/captures/remove}

		Las piedras capturadas desaparecen. Es lo mismo que |captures| o |captures = true|.
	\end{stylekey}

	\begin{stylekey}{/pgfgo/goban/captures/fade}

		Las piedras capturadas quedan en el tablero, dibujadas como prisioneros (con el estilo de |\prisoner| y de |/pgfgo/prisoner|).

		\begin{codeexample}[]
			\pgfgoset{goban = {5, partial = A1:E3}}

			\begin{goban}[captures = fade]
				\stones[w]{C1,D1}
				\stones[b]{B1,C2,D2}
				\moves{E1}
			\end{goban}
		\end{codeexample}
	\end{stylekey}

	\begin{stylekey}{/pgfgo/goban/captures/false}

		Desactiva las capturas. Es el valor inicial.
	\end{stylekey}

	\begin{key}{/pgfgo/goban/captured marks=\meta{keep o remove} (initially keep)}

		Decide qué pasa con las marcas que están sobre una piedra cuando la capturan. Solo tiene efecto con |captures|. Las marcas sobre puntos vacíos no cambian.

		\begin{codeexample}[]
			\pgfgoset{goban = {5, partial = A1:C3, captures}}

			\begin{goban}
				\stones[w]{A1}\stones[b]{B1}
				\marks[triangle]{A1}
				\moves{A2}
			\end{goban}

			\begin{goban}[captured marks = remove]
				\stones[w]{A1}\stones[b]{B1}
				\marks[triangle]{A1}
				\moves{A2} % el triángulo se va con la piedra blanca.
			\end{goban}
		\end{codeexample}
	\end{key}

	\begin{stylekey}{/pgfgo/goban/captured marks/keep}

		Las marcas quedan sobre el punto vacío. Es el valor inicial.
	\end{stylekey}

	\begin{stylekey}{/pgfgo/goban/captured marks/remove}

		Las marcas se van con la piedra, también con |captures = fade|. Las secuencias conservan su numeración.
	\end{stylekey}

	\begin{command}{\pgfgoprisoners\marg{player}}
		Escribe cuántos prisioneros tiene el jugador (|black| o |b|, |white| o |w|) en el último |goban| con |captures|: las piedras del rival que capturaron sus jugadas. Con otro jugador, da un error y escribe 0.

		\begin{codeexample}[]
			\pgfgoset{goban = {5, partial = A1:C3}}

			\begin{goban}[captures]
				\moves{B1,A1,A2}
			\end{goban}

			Negras: \pgfgoprisoners{black}. Blancas: \pgfgoprisoners{white}.
		\end{codeexample}
	\end{command}

	\subsection{Jugadas en puntos ya jugados}
	\label{subsec:notas}

	En los diagramas numerados, cuando una jugada cae en un punto donde ya se jugó (por ejemplo, al retomar un ko), los libros no la dibujan: el punto sigue mostrando la primera piedra y debajo del diagrama se anota, por ejemplo, «3 en 1».

	\begin{key}{/pgfgo/goban/notes=\meta{boolean} (default true, initially false)}

		Una jugada numerada en un punto donde ya hubo una jugada numerada del mismo |goban| no se dibuja: el punto sigue mostrando la primera, aunque haya sido capturada, y se anota «N at M». Solo cuentan las jugadas numeradas (|\moves| sin |numbers = false|) que se ven en el tablero. Funciona con los dos kernels.

		\begin{codeexample}[]
			\pgfgoset{goban = {5, partial = A2:E4}}

			\begin{goban}[notes]
				\stones[b]{B3,C2,C4}
				\stones[w]{D2,D4,E3}
				\moves[w]{C3}
				\moves{D3}
				\moves[from = 3, w]{C3}
			\end{goban}

			\pgfgomovenotes[at = en]
		\end{codeexample}
	\end{key}

	\begin{command}{\pgfgomovenotes\opt{\oarg{options}}}
		Escribe las notas del último |goban|, separadas por comas. Si no hay notas, no escribe nada. Las opciones son las de |/pgfgo/notes|.
	\end{command}

	\begin{key}{/pgfgo/notes = \meta{options}}
		Opciones de |\pgfgomovenotes| para todo el documento, por ejemplo |\pgfgoset{notes = {at = en}}|.
	\end{key}

	\begin{key}{/pgfgo/notes/at = \meta{text} (initially at o su traducción)}
		Texto entre los dos números.
	\end{key}

	\begin{key}{/pgfgo/notes/pass = \meta{text} (initially pass o su traducción)}
		Un pase de un archivo SGF (pág.~\pageref{subsec:sgf}) cuenta, pero no se dibuja. Con |notes|, si se numeraría, se anota «N: pass», en orden con las demás notas. Los pases con los que termina la partida no se anotan. Esta clave es el texto después del número.
	\end{key}

	\begin{key}{/pgfgo/notes/separator = \meta{text} (initially una coma y un espacio)}
		Texto entre las notas.

		\begin{codeexample}[]
			\begin{goban}[5, partial = A2:E4, notes]
				\stones[b]{B3,C2,C4}
				\stones[w]{D2,D4,E3}
				\moves[w]{C3}
				\moves{D3}
				\moves[from = 3, w]{C3}
				\moves{D3}
			\end{goban}

			\pgfgomovenotes[at = en, separator = {; }]
		\end{codeexample}
	\end{key}

	\subsection{Territorio}
	\label{subsec:territorio}

	Al final de la partida, el territorio puede marcarse a mano con |\marks| (pág.~\pageref{sec:marcas}) o con |\territory|, que además lo cuenta.

	\begin{command}{\territory\opt{\oarg{options}}\marg{polygon}}
		Marca como territorio los puntos del polígono que tiene esos vértices (en orden), incluido su borde. No marca los puntos que tienen una piedra viva: los prisioneros cuentan como territorio. Los vértices se escriben como en |\stones|; uno solo marca un punto, y dos, un segmento. Funciona con los dos kernels.

		El territorio se dibuja al terminar el |goban|. Cada punto queda con el último |\territory| que lo marcó, con el jugador y las opciones de ese |\territory| (así, |\territory[n]{A1}| después de otro |\territory| vuelve neutral a A1), y no se dibuja ni se cuenta si al final tiene una piedra viva.

		\begin{codeexample}[width=4.75cm]
			\begin{goban}[label = {at = {above, left}}]
		\stones[b] {A6|-B5|-D1|-E2, C4, E9-|F4, D5-|G7--D7, H6, H8, I7}
		\stones[w] {E8-|C5, B:E6, A7, B8, D9, I6|- H5, H2:3, F1|-C3, C2, D:E4, G4}
		\prisoners[b]{H1-|I2}
		\prisoner[w] {A1}

				\territory[b]{A1,A5,B5,B1}
				\territory[b]{F9,F7,G7,G6,I6,I9}
				\territory[w]{F1,F4,G4,G5,I5,I1}
				\territory[w]{A9,A6,C6,C8,D8,D9}
			\end{goban}
		\end{codeexample}
	\end{command}

	\begin{command}{\territory\opt{\oarg{options}}\marg{auto}}
		Con |auto| en lugar de un polígono, busca solo el territorio: cada región de puntos sin piedras vivas es del color que la rodea. Si la tocan los dos colores (dame, seki) o ninguno, es neutral. Los prisioneros cuentan como piedras muertas. Sin jugador, marca el territorio de los dos; con |b| o |w|, solo el de ese jugador. Solo con el kernel en Lua (pág.~\pageref{subsec:kernel}).

		\begin{codeexample}[width=4.75cm]
			\begin{goban}[label = {at = {above, left}}]
		\stones[b] {A6|-B5|-D1|-E2, C4, E9-|F4, D5-|G7--D7, H6, H8, I7}
		\stones[w] {E8-|C5, B:E6, A7, B8, D9, I6|- H5, H2:3, F1|-C3, C2, D:E4, G4}
		\prisoners[b]{H1-|I2}
		\prisoner[w] {A1}

				\territory[dame]{auto}
			\end{goban}
		\end{codeexample}

		|auto| solo mira qué color toca cada región; no sabe leer la posición. Por eso:
		\begin{itemize}
			\item Es para partidas terminadas. En una posición abierta, casi todo el tablero toca los dos colores y es neutral; con piedras de un solo color, todo el tablero es suyo. Si un borde tiene un hueco, el territorio de adentro se junta con el dame y también queda neutral.
			\item No sabe qué piedras están muertas: hay que ponerlas con |\prisoners|. Una piedra muerta puesta con |\stones| cuenta como viva, y la región que toca queda neutral.
			\item Los ojos de un grupo en seki cuentan como territorio (en las reglas japonesas no se cuentan). Pueden volverse neutrales después, con un |\territory[n]| para cada uno.
			\item Las piedras neutrales (sin jugador) ocupan su punto, pero no son de nadie: una región que solo toca piedras neutrales es neutral.
			\item Busca las regiones con las piedras que hay cuando aparece |\territory|, así que conviene ponerlo al final del |goban|.
			\item Con |partial = auto|, el tablero crece hasta mostrar todo el territorio (y el dame, con |dame|).
		\end{itemize}
	\end{command}

	\begin{key}{/pgfgo/territory/black}
		El territorio es de las negras. También puede escribirse |b|.
	\end{key}

	\begin{key}{/pgfgo/territory/white}
		El territorio es de las blancas. También puede escribirse |w|.
	\end{key}

	\begin{key}{/pgfgo/territory/neutral}
		Puntos neutrales. Es lo que marca un polígono sin jugador. También puede escribirse |n|.
	\end{key}

	\begin{key}{/pgfgo/territory/b}
		Abreviatura de |black|.
	\end{key}

	\begin{key}{/pgfgo/territory/w}
		Abreviatura de |white|.
	\end{key}

	\begin{key}{/pgfgo/territory/n}
		Abreviatura de |neutral|.
	\end{key}

	\begin{key}{/pgfgo/territory/draw=\meta{marks, fill o region} (initially marks)}
		Con |marks|, cada punto lleva la marca de territorio del jugador. Con |fill|, la región se rellena debajo de las piedras.

		\begin{codeexample}[width=4.75cm]
			\begin{goban}[label = {at = {above, left}}]
		\stones[b] {A6|-B5|-D1|-E2, C4, E9-|F4, D5-|G7--D7, H6, H8, I7}
		\stones[w] {E8-|C5, B:E6, A7, B8, D9, I6|- H5, H2:3, F1|-C3, C2, D:E4, G4}
		\prisoners[b]{H1-|I2}
		\prisoner[w] {A1}

				\territory[draw = fill]{auto}
			\end{goban}
		\end{codeexample}

		Con |region|, el territorio de cada color es un solo polígono: se pinta cada casilla de la cuadrícula que tiene alguna esquina en el territorio (las piedras muertas cuentan como territorio). Sus lados van por las líneas del tablero, sus vértices están siempre en intersecciones, y llega hasta el borde del tablero. Dos zonas que se tocan por una sola esquina quedan separadas. Todos los |\territory| de un color en el mismo |goban| (también el territorio |TB| y |TW| de |\sgf|) forman una sola región, que se dibuja al terminar el |goban| con el estilo del último. Los puntos neutrales se rellenan uno por uno, como con |fill|. El contorno y sus vértices pueden dibujarse encima, con |line| y |vertices|, y |\pgfgoregion| da sus propiedades.

		\begin{codeexample}[width=4.75cm]
			\begin{goban}[label = {at = {above, left}}]
		\stones[b] {A6|-B5|-D1|-E2, C4, E9-|F4, D5-|G7--D7, H6, H8, I7}
		\stones[w] {E8-|C5, B:E6, A7, B8, D9, I6|- H5, H2:3, F1|-C3, C2, D:E4, G4}
		\prisoners[b]{H1-|I2}
		\prisoner[w] {A1}

				\territory[draw = region]{auto}
			\end{goban}
		\end{codeexample}
	\end{key}

	\begin{key}{/pgfgo/territory/dame=\meta{boolean} (default true, initially false)}
		Con |auto|, marca también los puntos neutrales.
	\end{key}

	\begin{key}{/pgfgo/territory/forget}
		|resume| no trae este territorio, aunque se pida |territory| (p.~\pageref{subsec:recordar}).
	\end{key}

	Las demás opciones de |\territory| cambian el estilo del jugador solo para ese comando, como |mark| y |fill| en |\territory[w, mark = {cross}]{auto}|. Sin jugador, cambian el de los tres. El estilo de cada jugador se fija con |/pgfgo/players|:

	\begin{key}{/pgfgo/players/black/territory=\meta{territory options}}
		Estilo del territorio de cada jugador (|black|, |white| o |neutral|).
	\end{key}

	\begin{key}{/pgfgo/players/black/territory/mark=\meta{mark options} (initially t-circle)}
		Marca de cada punto, con |draw = marks|. Es |t-circle| para las negras y |t-square| para las blancas y los puntos neutrales.

		\begin{codeexample}[]
			\begin{goban}[5]
				\stones[b]{C1:5}
				\stones[w]{D1:5}
				\territory[b, mark = {t-square}]{A1, A5, B5, B1}
				\territory[w]{E1, E5}
			\end{goban}
		\end{codeexample}
	\end{key}

	\begin{key}{/pgfgo/players/black/territory/fill=\meta{fill options} (default true)}
		Relleno, con |draw = fill| y |draw = region|.
	\end{key}

	\begin{key}{/pgfgo/players/black/territory/fill/true}
		Dibuja el relleno. Es el valor inicial.
	\end{key}

	\begin{key}{/pgfgo/players/black/territory/fill/false}
		No dibuja el relleno: con |draw = region| quedan solo el contorno o los vértices.
	\end{key}

	\begin{key}{/pgfgo/players/black/territory/fill/color=\meta{color} (initially black)}
		Color del relleno: |black| para las negras, |white!80!black| para las blancas y |gray| para los puntos neutrales.
	\end{key}

	\begin{key}{/pgfgo/players/black/territory/fill/opacity=\meta{number} (initially 0.35)}
		Opacidad del relleno: |0.35| para las negras y los puntos neutrales, |0.7| para las blancas.

		\begin{codeexample}[]
			\begin{goban}[5]
				\stones[b]{C1:5}
				\stones[w]{D1:5}
				\territory[b, draw = fill,
					fill = {color = blue, opacity = 0.5}]{A1, A5, B5, B1}
				\territory[w, draw = fill]{E1, E5}
			\end{goban}
		\end{codeexample}
	\end{key}

	\begin{key}{/pgfgo/players/black/territory/line=\meta{line options} (default true)}
		Contorno de la región, con |draw = region|. Sola o con |true|, dibuja el contorno; las opciones solas no lo prenden.

		\begin{codeexample}[width=4.75cm]
			\begin{goban}[label = {at = {above, left}}]
		\stones[b] {A6|-B5|-D1|-E2, C4, E9-|F4, D5-|G7--D7, H6, H8, I7}
		\stones[w] {E8-|C5, B:E6, A7, B8, D9, I6|- H5, H2:3, F1|-C3, C2, D:E4, G4}
		\prisoners[b]{H1-|I2}
		\prisoner[w] {A1}

				\territory[b, draw = region, line,
					vertices]{auto}
				\territory[w, draw = region,
					line = {true, color = red}]{auto}
			\end{goban}
		\end{codeexample}
	\end{key}

	\begin{key}{/pgfgo/players/black/territory/line/true}
		Dibuja el contorno.
	\end{key}

	\begin{key}{/pgfgo/players/black/territory/line/false}
		No dibuja el contorno. Es el valor inicial.
	\end{key}

	\begin{key}{/pgfgo/players/black/territory/line/color=\meta{color} (initially black)}
		Color del contorno: |black| para las negras y |white!40!black| para las blancas.
	\end{key}

	\begin{key}{/pgfgo/players/black/territory/line/width=\meta{dimension} (initially 1.2pt)}
		Grosor del contorno.
	\end{key}

	\begin{key}{/pgfgo/players/black/territory/line/opacity=\meta{number} (initially 1)}
		Opacidad del contorno y de los vértices.
	\end{key}

	\begin{key}{/pgfgo/players/black/territory/line/above=\meta{boolean} (default true, initially true)}
		Con |true|, el contorno y los vértices se dibujan encima de las piedras y de las marcas. Con |false|, debajo de las piedras, junto con el relleno.

		\begin{codeexample}[]
			\begin{goban}[5]
				\stones[b]{C1:5}
				\stones[w]{D1:5}
				\territory[b, draw = region, fill = false,
					line = {true, color = red, width = 2pt, opacity = 0.7}]{A1, A5, B5, B1}
			\end{goban}

			\begin{goban}[5]
				\stones[b]{C1:5}
				\stones[w]{D1:5}
				\territory[b, draw = region, fill = false,
					line = {true, color = red, width = 2pt, opacity = 0.7, above = false}]{A1, A5, B5, B1}
			\end{goban}
		\end{codeexample}
	\end{key}

	\begin{key}{/pgfgo/players/black/territory/vertices=\meta{options} (default true)}
		Puntos en los vértices del contorno, con |draw = region|, como los puntos de estrella. Se dibujan en la misma capa que el contorno. Sola o con |true|, los dibuja; las opciones solas no los prenden.
	\end{key}

	\begin{key}{/pgfgo/players/black/territory/vertices/true}
		Dibuja los vértices.
	\end{key}

	\begin{key}{/pgfgo/players/black/territory/vertices/false}
		No dibuja los vértices. Es el valor inicial.
	\end{key}

	\begin{key}{/pgfgo/players/black/territory/vertices/radius=\meta{dimension} (initially 2pt)}
		Radio de los puntos.
	\end{key}

	\begin{key}{/pgfgo/players/black/territory/vertices/color=\meta{color}}
		Color de los puntos. Sin esta \emph{key}, toman el color del contorno.

		\begin{codeexample}[]
			\begin{goban}[5]
				\stones[b]{C1:5}
				\stones[w]{D1:5}
				\territory[b, draw = region,
					vertices = {true, radius = 3pt, color = red}]{A1, A5, B5, B1}
			\end{goban}
		\end{codeexample}
	\end{key}

	\begin{command}{\pgfgoscore\opt{\oarg{options}}\marg{player}}
		Escribe el puntaje del jugador (|black| o |b|, |white| o |w|) en el último |goban|, con las reglas japonesas: los puntos de su territorio, más las piedras muertas del rival que hay en él, más las piedras que capturó (con |captures|), más el |komi| para las blancas. Un punto marcado dos veces cuenta una sola vez. Con |neutral| o |n|, escribe la cantidad de puntos neutrales marcados. Con otro jugador, da un error y escribe 0.

		\begin{codeexample}[code only]
			Negras: \pgfgoscore{black}. Blancas: \pgfgoscore[komi = 6.5]{white}.
		\end{codeexample}

		Con el tablero de arriba, escribe «Negras: 11. Blancas: 23.5».
	\end{command}

	\begin{key}{/pgfgo/score/komi=\meta{number} (initially 0)}
		Komi que se suma al puntaje de las blancas. Puede fijarse para todo el documento con |\pgfgoset{score = {komi = 6.5}}|.

		\begin{codeexample}[]
			\begin{goban}[5]
				\stones[b]{C1:5}
				\stones[w]{D1:5}
				\territory{auto}
			\end{goban}

			Negras: \pgfgoscore{b}. Blancas: \pgfgoscore[komi = 6.5]{w}.
		\end{codeexample}
	\end{key}

	\begin{command}{\pgfgoregion\marg{player}\marg{property}}
		Escribe una propiedad de la región (|draw = region|) del jugador (|black| o |b|, |white| o |w|) en el |goban| actual o en el último. Es expandible, así que puede usarse en una cuenta o dentro de las coordenadas de |\marks|. Las propiedades son del contorno completo, aunque el tablero sea parcial. Sin región, los números son 0. Las propiedades son:
		\begin{itemize}
			\item |polygons|: la cantidad de contornos, exteriores y huecos.
			\item |outlines|: la cantidad de contornos exteriores, uno por cada zona separada.
			\item |holes|: la cantidad de huecos.
			\item |vertices|: la cantidad de vértices de todos los contornos.
			\item |area|: la cantidad de casillas pintadas.
			\item |perimeter|: el largo de todos los contornos, en casillas.
			\item |polygon| \meta{i}: los vértices del contorno \meta{i}, como coordenadas separadas por comas (por ejemplo, |A1,D1,D4,A4|). Vacío si no hay tantos contornos.
		\end{itemize}
		Cada contorno empieza en su vértice más bajo y, entre esos, en el de más a la izquierda. Los exteriores van en sentido antihorario y los huecos, en sentido horario. Los contornos están ordenados por ese primer vértice, de abajo hacia arriba.

		\begin{codeexample}[width=4.75cm]
			\begin{goban}[label = {at = {above, left}}]
				\stones[w]{D4:F6}
				\territory[b, draw = region, line]{B2,B8,H8,H2}
				\marks[triangle]{\pgfgoregion{b}{polygon 2}}
			\end{goban}

			Contornos: \pgfgoregion{b}{polygons}
			(\pgfgoregion{b}{holes} hueco).
			Casillas: \pgfgoregion{b}{area}.
		\end{codeexample}
	\end{command}

	\section{Recordar tableros y vistas}
\label{sec:memoria}

	\subsection{Recordar un tablero}
	\label{subsec:recordar}
	Un |goban| con nombre guarda todo lo que dibuja (piedras, jugadas, marcas, flechas, líneas y territorio, también lo que viene de |\sgf|), junto con sus opciones. Otro |goban| puede empezar con ese tablero tal como quedó, con |resume|, y |\pgfgoview| lo muestra en el texto (p.~\pageref{subsec:vistas}).

	\begin{codeexample}[]
		\pgfgoset{goban = {5, partial = B2:D4}}

		\begin{goban}[name = tablero1]
			\stones[b]{C2,D3}
			\moves*[w]{B4}
		\end{goban}

		\begin{goban}[resume = tablero1]
			\moves*{C3} % Sigue con la 2 (negras).
		\end{goban}

		\begin{goban}[resume = {tablero1, numbers}]
			\moves*{C3}
		\end{goban}
	\end{codeexample}

	\begin{key}{/pgfgo/goban/name=\meta{name}}
		Guarda con ese nombre todo lo que se dibuje en este |goban|, junto con sus opciones. El nombre es libre (por ejemplo, |gob:tablero1|). Otro |goban| con el mismo nombre lo reemplaza.
	\end{key}

	\begin{key}{/pgfgo/goban/resume=\meta{name, options}}
		Empieza el |goban| con el tablero que tiene ese nombre, como quedó al terminar (con las capturas resueltas: si ese |goban| usaba |captures|, este también). La numeración y el turno siguen desde ahí. Por defecto trae solo las piedras, sin números, marcas ni territorio; las opciones piden el resto. Las opciones del tablero (tamaño, |partial|, coordenadas) se dan en cada |goban|. Si no hay ningún |goban| con ese nombre, se produce un error.

		Con |name| y |resume| en el mismo |goban| (en cualquier orden), el nuevo nombre guarda lo que trajo, tal como se dibujó, más lo suyo, y el tablero de |resume| no cambia: así se arman variantes. Con el mismo nombre en las dos, el |goban| sigue agregando a ese tablero.

		\begin{codeexample}[]
			\pgfgoset{goban = {5, partial = B2:D4}}

			\begin{goban}[name = base]
				\stones[b]{C2,D3}
				\stones[w]{C3}
			\end{goban}

			\begin{goban}[name = variante1, resume = base]
				\moves*[w]{D2}
			\end{goban}

			\begin{goban}[name = variante2, resume = base]
				\moves*[w]{B3}
			\end{goban}

			\begin{goban}[resume = {variante1, numbers}]
				\moves*{B4}
			\end{goban}
		\end{codeexample}
	\end{key}

	\begin{key}{/pgfgo/goban/resume/numbers=\meta{boolean} (default true, initially false)}
		Trae también los números de las jugadas.
	\end{key}

	\begin{key}{/pgfgo/goban/resume/marks=\meta{boolean} (default true, initially false)}
		Trae también las marcas, las flechas y las líneas.
	\end{key}

	\begin{key}{/pgfgo/goban/resume/territory=\meta{boolean} (default true, initially false)}
		Trae también el territorio.

		\begin{codeexample}[]
			\pgfgoset{goban = {5, partial = B2:D4}}

			\begin{goban}[name = guardado]
				\stones[b]{C2:4}
				\moves*[w]{D3}
				\marks[triangle]{B3}
				\territory[b]{B2, B4}
			\end{goban}

			\begin{goban}[resume = guardado]
			\end{goban}

			\begin{goban}[resume = {guardado, numbers, marks, territory}]
			\end{goban}
		\end{codeexample}
	\end{key}

	Para que |resume| no traiga una piedra, una jugada, una marca, una flecha o un territorio, se dibuja con la opción |forget| (|\stone[forget]{C3}|, |\marks[forget, circle]{A1}|, \ldots); con |label = forget|, una jugada viene sin su número aunque se pida |numbers|. Lo olvidado queda guardado: |\pgfgoview| lo muestra con |forgotten|.

	\subsection{Vistas en el texto}
	\label{subsec:vistas}
	Un |goban| con nombre guarda todo lo que dibuja. Después, |\pgfgoview| muestra una parte de ese tablero, tal como quedó al terminar (con las capturas resueltas), en medio de un párrafo: una sola piedra o un recuadro.

	\begin{codeexample}[]
		\begin{goban}[5, name = tablero1, captures]
			\stones[b]{B3, C2, C4}
			\stones[w]{D2, D4, E3}
			\moves*[w]{C3, D3}
		\end{goban}

		La negra en D3 \pgfgoview{tablero1}{D3} captura la blanca de C3 \pgfgoview{tablero1}{C3}.
		Así quedó el ko: \pgfgoview{tablero1}{B2:E4}.
	\end{codeexample}

	\begin{command}{\pgfgoview\opt{\oarg{options}}\marg{name}\marg{part}}
		Dibuja la intersección o el rectángulo (escrito como en |partial|) del |goban| con ese nombre, como quedó al terminar. Por defecto se ve todo lo que dibujó el |goban| en esa parte (piedras, números de las jugadas, marcas, flechas y territorio) con el fondo y las coordenadas de ese |goban|, al tamaño del texto y centrado a la altura de las minúsculas. Las opciones son las de |/pgfgo/view|, que quitan las marcas, los números o el territorio, o agregan lo olvidado, y, además, cualquier opción de |goban|, que se aplica después de las del |goban| con nombre: así se quitan el fondo, las coordenadas, la cuadrícula o los puntos de estrella, o se cambia el estilo de las piedras, las marcas y el tablero (|players|, |stone|, |mark|, |grid|, |background|, \ldots).

		La vista usa el estilo que tenía ese |goban| al dibujarse (el perfil, lo fijado con |\pgfgoset| y sus opciones), aunque después cambie. Encima van el perfil de las vistas (|view/profile|), |\pgfgoset{view = ...}| y las opciones de la vista.

		Una vista no cambia nada de lo que guarda el último |goban|: |\pgfgoscore|, |\pgfgoprisoners| y |\pgfgomovenotes| siguen hablando de él, y lo que guarda cada |goban| con nombre queda igual.

		\begin{codeexample}[]
			\begin{goban}[5, name = tablero2]
				\stones[b]{B2, C3}
				\moves*[w]{C2, D3}
				\marks[b, triangle]{C3}
			\end{goban}

			\pgfgoview{tablero2}{B2:D3}
			\pgfgoview[marks = false, numbers = false]{tablero2}{B2:D3}
			\pgfgoview[label = false, grid = false]{tablero2}{B2:D3}
			\pgfgoview[background = false]{tablero2}{B2:D3}
		\end{codeexample}
	\end{command}

	\begin{key}{/pgfgo/view=\meta{options}}
		Opciones para todas las vistas, con |\pgfgoset|. Las opciones de |goban| que no son de |/pgfgo/view| se agregan a la que traen las vistas: |scale = 0.8|.

		\begin{codeexample}[]
			\pgfgoset{view = {background = false, scale = 1}}
			Sin fondo y más grande: \pgfgoview{tablero2}{C3}.
		\end{codeexample}
	\end{key}

	\begin{key}{/pgfgo/view/marks=\meta{boolean} (default true, initially true)}
		Con |false|, no dibuja las marcas, las flechas ni las líneas.
	\end{key}

	\begin{key}{/pgfgo/view/numbers=\meta{boolean} (default true, initially true)}
		Con |false|, no dibuja los números de las jugadas.
	\end{key}

	\begin{key}{/pgfgo/view/territory=\meta{boolean} (default true, initially true)}
		Con |false|, no dibuja el territorio.
	\end{key}

	\begin{key}{/pgfgo/view/forgotten=\meta{boolean} (default true, initially false)}
		Dibuja también lo que se dibujó con |forget| (p.~\pageref{subsec:recordar}).

		\begin{codeexample}[]
			\begin{goban}[5, name = tablero3]
				\stones[b]{B2, C3}
				\stone[w, forget]{D3}
			\end{goban}

			\pgfgoview{tablero3}{B2:D3}
			\pgfgoview[forgotten]{tablero3}{B2:D3}
		\end{codeexample}
	\end{key}

	\begin{key}{/pgfgo/view/raise=\meta{dimension} (initially 0.5ex)}
		A qué altura sobre la línea de base del texto queda el centro de la vista.

		\begin{codeexample}[]
			Texto \pgfgoview[raise = 0pt]{tablero2}{C3}
			texto \pgfgoview{tablero2}{C3}
			texto \pgfgoview[raise = 1.5ex]{tablero2}{C3}.
		\end{codeexample}
	\end{key}

	\begin{key}{/pgfgo/view/profile=\meta{name} (initially false)}
		Un perfil para las vistas (p.~\pageref{subsec:perfiles}), por ejemplo |simple|. Va después del estilo del |goban| con nombre y antes de las demás opciones de la vista. Con |false|, la vista usa el estilo del |goban| tal cual.

		\begin{codeexample}[]
			\pgfgoview{tablero1}{B2:E4}
			\pgfgoview[profile = simple]{tablero1}{B2:E4}
			\pgfgoview[profile = wain]{tablero1}{B2:E4}
		\end{codeexample}
	\end{key}

	\section{Partidas en SGF}
\label{subsec:sgf}

SGF es el formato que usan casi todos los programas de Go para guardar partidas. Con el kernel en Lua (pág.~\pageref{subsec:kernel}), \pgfgoname\ puede leerlas.

En todos los comandos que leen archivos SGF, el nombre puede tener espacios, acentos, paréntesis y otras escrituras. Los caracteres especiales de \TeX\ se escriben con barra: |\sgf{a\&b.sgf}| lee |a&b.sgf|, y lo mismo vale para |\#|, |\%|, |\$|, |\{| y |\}|. Dos espacios seguidos se leen como uno, porque \TeX\ los junta antes de que \pgfgoname\ vea el nombre.

Cada archivo se lee una sola vez en todo el documento, por más comandos que lo usen (si cambia mientras se compila, se vuelve a leer).

\subsection{Nombres para los archivos}
\label{subsec:sgfnombres}

\begin{command}{\pgfgodefinesgf\opt{\oarg{options}}\marg{name}\marg{file}}
	Le da un nombre a una partida: desde ahí, \meta{name} sirve en todos los lugares donde va un archivo SGF (|\sgf|, |\sgfinfo|, |\sgfcomment|, |\sgfdiagrams|, |\sgfproblem|, |\sgfanimation|, |sgf size| y |sgf view|). Así, la ruta del archivo se escribe una sola vez, y lo que dice qué partida es (|game| en una colección, la codificación) también. Un nombre que no está definido es un archivo, como siempre; si un nombre coincide con el de un archivo, vale el nombre. Las opciones de cada comando ganan sobre las del nombre: |\sgfinfo[game = 1]{|\meta{name}|}{PB}| lee la partida 1 aunque el nombre diga otra. Con |kernel = expl3|, el nombre se define pero no se puede usar, como los comandos SGF.

	\begin{codeexample}[]
		\pgfgodefinesgf{partida}{partida.sgf}

		\begin{goban}[sgf size = partida, scale = 0.6]
			\sgf[to = 7]{partida}
		\end{goban}

		Jugada 7: \sgfcomment{partida}{7}
	\end{codeexample}
\end{command}

\begin{key}{/pgfgo/define sgf/game = \meta{number} (initially la de cada comando)}
	En una colección, la partida que usan los comandos con este nombre, como |/pgfgo/sgf/game|.

	\begin{codeexample}[code only]
		\pgfgodefinesgf[game = 3]{tercera}{coleccion.sgf}
	\end{codeexample}
\end{key}

\begin{key}{/pgfgo/define sgf/charset = \meta{charset} (initially la de \texttt{sgf charset})}
	La codificación con la que se lee el archivo si no es UTF-8, como |/pgfgo/sgf charset|, solo para este nombre.

	\begin{codeexample}[code only]
		\pgfgodefinesgf[charset = GBK]{qipu}{partidas/qipu-sin-ca.sgf}
	\end{codeexample}
\end{key}

\subsection{Una partida en un tablero}
\label{subsec:sgftablero}

\begin{command}{\sgf\opt{\oarg{options}}\marg{file}}

	Dibuja en el |goban| la línea principal de la primera partida del archivo: las piedras iniciales (|AB| y |AW|) como con |\stones|, y las jugadas como con |\moves|. Las variantes no se dibujan. Los pases cuentan para la numeración, pero no se dibujan (con |goban/notes| se anotan, «N: pass»). Solo existe dentro de |goban|.
\end{command}

Las piedras que agregan (|AB|, |AW|) o quitan (|AE|) los nodos posteriores a la primera jugada se aplican junto con la jugada anterior. Quitar piedras necesita |captures|; sin |captures|, |\sgf| avisa una vez y no las quita.

Al terminar, la numeración y el turno quedan listos para seguir la partida con |\moves|. |\sgf| funciona con |partial|, |captures|, |name| y |resume|. Si la partida es para un tablero de otro tamaño, se produce una advertencia. Con |kernel = expl3|, |\sgf| produce un error.

En los ejemplos, los archivos se usan por su nombre (pág.~\pageref{subsec:sgfnombres}). |partida| es |partida.sgf|, una partida corta de 9$\times$9:
\begin{verbatim}
(;GM[1]FF[4]SZ[9]PB[Negras]PW[Blancas]C[Una partida corta.]
;B[ee];W[ef];B[df];W[ce];B[ff];W[gd];B[eg]C[Negras capturan.];W[dd];B[cf];W[fd])
\end{verbatim}

\begin{codeexample}[width=5cm]
	\begin{goban}[label = {at = {north, west}}]
		\sgf{partida} % Sin captures, 2 no se retira.
	\end{goban}
\end{codeexample}

\begin{codeexample}[width=5cm]
	\begin{goban}[label = {at = {north, west}}, captures]
		\sgf{partida}
	\end{goban}
\end{codeexample}

\begin{key}{/pgfgo/sgf/from = \meta{number} (initially 1)}

	Número de la primera jugada numerada. Las jugadas anteriores se dibujan sin número.

	\begin{codeexample}[width=5cm]
		\begin{goban}[captures]
			\sgf[from = 7]{partida}
		\end{goban}
	\end{codeexample}
\end{key}

\begin{key}{/pgfgo/sgf/to = \meta{number} o |last| (initially last)}

	Número de la última jugada que se juega. Las siguientes no se dibujan. Con |last| (o un número negativo), hasta la última.

	\begin{codeexample}[width=5cm]
		\begin{goban}[captures]
			\sgf[to = 6]{partida}
			\moves{F5} % Sigue con la 7 (negras), en otro lugar.
		\end{goban}
	\end{codeexample}
\end{key}


|from| y |to| cuentan las jugadas del archivo. El número que se muestra puede ser otro: la propiedad |MN| del archivo lo cambia (las jugadas siguientes siguen desde ahí), y también |restart|. En los ejemplos que siguen, |variantes| es |variantes.sgf|, una partida de 7$\times$7 con dos bifurcaciones, marcas y |MN[10]| en la cuarta jugada de la línea principal.

\begin{key}{/pgfgo/sgf/variation = \meta{números} (initially la línea principal)}
	Elige una variante: en cada bifurcación de la partida se toma, en orden, la variante indicada (1 es la principal). Por ejemplo, |variation = 2.1| toma la segunda variante en la primera bifurcación y la primera en la segunda. Cuando la lista se termina, se sigue por la primera variante. Si una variante no existe, se produce una advertencia.

	\begin{codeexample}[]
		\begin{goban}[sgf size = variantes]
			\sgf[variation = 2]{variantes}
		\end{goban}

		\begin{goban}[sgf size = variantes]
			\sgf[variation = 1.2]{variantes}
		\end{goban}
	\end{codeexample}
\end{key}

\begin{key}{/pgfgo/sgf/marks = \meta{boolean} (default true, initially true)}
	Dibuja las marcas del archivo (|TR|, |SQ|, |CR|, |MA|, |SL| y las etiquetas de |LB|), las flechas (|AR|) y las líneas (|LN|) del último nodo que se muestra, y atenúa los puntos de |DD| (que siguen atenuados hasta el próximo |DD|; pág.~\pageref{subsec:flechas}). Cada marca toma el color de la piedra que hay en ese punto, para que se vea, y reemplaza al número de la jugada. Las etiquetas son la excepción: no tapan un número, así que sobre una piedra numerada no se dibujan. El territorio del archivo (|TB| y |TW|) se dibuja como el de |\territory| (pág.~\pageref{subsec:territorio}) y cuenta para |\pgfgoscore|.

	\begin{codeexample}[]
		\begin{goban}[sgf size = variantes]
			\sgf[to = 2]{variantes}
		\end{goban}

		\begin{goban}[sgf size = variantes]
			\sgf[to = 2, marks = false]{variantes}
		\end{goban}
	\end{codeexample}
\end{key}

\begin{key}{/pgfgo/sgf/step = \meta{number} o \meta{list of numbers} (initially 100)}
	El tamaño de los tramos en que |restart| reinicia la numeración. Con una lista, cada número es el tamaño de un tramo, y el último se repite: con |step = {50, 100}|, el primer tramo tiene 50 jugadas y los demás 100. Es la misma \emph{key} que |/pgfgo/sgf diagrams/step|, donde cada tramo es un diagrama.
\end{key}

\begin{key}{/pgfgo/sgf/restart = \meta{restart options} (default true)}
	Desde el segundo tramo de |step| jugadas, cada uno numera desde 1 (o desde |restart/from|), como en los libros: con |step = 100, restart|, la jugada 101 se muestra como 1. El primer tramo conserva los números de la partida. Las mismas opciones tienen |\sgfanimation| y |\sgfdiagrams|.

	\begin{codeexample}[]
		\begin{goban}[sgf size = variantes]
			\sgf[]{variantes}
		\end{goban}

		\begin{goban}[sgf size = variantes]
			\sgf[step = 3, restart]{variantes}
		\end{goban}
	\end{codeexample}
\end{key}

\begin{key}{/pgfgo/sgf/restart/true}
	Reinicia la numeración en cada tramo.
\end{key}

\begin{key}{/pgfgo/sgf/restart/false}
	No la reinicia: las jugadas llevan los números de la partida (con |MN|).
\end{key}

\begin{key}{/pgfgo/sgf/restart/auto}
	Hace lo que pide el archivo con |PM|: con |PM[2]|, la numeración se reinicia cada 100 jugadas, sin importar |step| (la 101 se muestra como 1); si no, no se reinicia. Es el valor inicial.
\end{key}

\begin{key}{/pgfgo/sgf/restart/from = \meta{number} (initially 1)}
	El número de la primera jugada de cada tramo, desde el segundo. También prende |restart|. Con |step = 99| y |restart = {from = 0}|, la jugada 100 se muestra como 0, la 101 como 1, y así.

	\begin{codeexample}[]
		\begin{goban}[sgf size = variantes]
			\sgf[step = 3, restart = {from = 0}]{variantes}
		\end{goban}
	\end{codeexample}
\end{key}

\begin{key}{/pgfgo/sgf/numbers = \meta{true, false o auto} (default true, initially auto)}
	Con |false|, las jugadas se dibujan sin número. Con |auto|, hace lo que pide el archivo con |PM|: con |PM[0]|, no se numeran. |\sgfdiagrams| tiene su propia \emph{key} |numbers|, y |\sgfproblem| numera siempre.

	\begin{codeexample}[]
		\begin{goban}[sgf size = variantes]
			\sgf[numbers = false]{variantes}
		\end{goban}
	\end{codeexample}
\end{key}

\begin{key}{/pgfgo/sgf/last mark = \meta{mark options} o |false| (initially false)}
	Marca la última jugada que se muestra, con las opciones de una marca (|circle|, |triangle|, |custom = ...|, etc.). Un pase no se marca, y tampoco un punto que el archivo ya marca ni una jugada con número (el número ya dice cuál es la última): sirve con |numbers = false|.

	\begin{codeexample}[]
		\begin{goban}[sgf size = variantes]
			\sgf[numbers = false, last mark = circle]{variantes}
		\end{goban}
	\end{codeexample}
\end{key}

\begin{key}{/pgfgo/sgf/game = \meta{number} (initially 1)}
	Un archivo SGF puede tener varias partidas (una colección). Elige cuál se lee; si no existe, se lee la primera con un aviso. También vale para |\sgfanimation|, |\sgfdiagrams|, |\sgfcomment|, |\sgfinfo| y |sgf size|, que tienen su propia \emph{key} |game|; para todo el documento, |\pgfgoset{sgf = {game = 2}}|.
\end{key}

\begin{key}{/pgfgo/sgf/variations = \meta{none, children, siblings o auto} (default auto, initially none)}
	Marca con letras las variantes de la última jugada que se muestra: con |children|, las jugadas que pueden seguir; con |siblings|, las alternativas a esa jugada (la letra reemplaza a su número). |auto| hace lo que pide el archivo con |ST|. Solo se marcan si hay más de una.

	\begin{codeexample}[]
		\begin{goban}[sgf size = variantes]
			\sgf[to = 1, marks = false, variations = children]{variantes}
		\end{goban}
	\end{codeexample}
\end{key}

|\sgf| también lee |PL| (quién juega después: |\moves| sigue con ese color) y |CA| (la codificación del archivo; ver |sgf charset|).

\begin{key}{/pgfgo/sgf charset = \meta{charset} (initially auto)}
	Codificación de los archivos SGF que no están en UTF-8, para todos los comandos que leen archivos. Con |auto|, es la que dice el archivo con |CA| o, si no dice ninguna, Latin-1 (la de SGF por defecto). Se leen |Latin-1|, |GBK| (y |GB2312|), |Big5|, |Shift_JIS|, |EUC-JP| y |EUC-KR|; los nombres no distinguen mayúsculas ni signos (|shift-jis| y |EUC_KR| también valen). Un archivo que es UTF-8 válido se lee siempre como UTF-8, aunque |CA| diga otra cosa; uno en UTF-16 (con la marca de orden de bytes al principio) también se lee.

	Muchos archivos de servidores asiáticos no dicen su codificación; para leerlos, por ejemplo, |\pgfgoset{sgf charset = GBK}|. Si un archivo sin |CA| parece estar en una de esas codificaciones, se produce un aviso que lo recuerda. Con otra codificación se produce un aviso y el archivo se lee como Latin-1: hay que convertirlo antes a UTF-8, por ejemplo con |iconv -f KOI8-R -t UTF-8 partida.sgf > partida-utf8.sgf|.

	Para escribir los comentarios en chino, japonés o coreano hace falta además el paquete de ese idioma: |ctex|, |luatexja| o |luatexko|. Los ejemplos 14, 18 y 19 de |doc/examples| muestran una partida comentada en cada uno. Sin él, LaTeX solo anota los caracteres que faltan en el |.log|; por eso, si la fuente con la que se escribe un comentario o un |\sgfinfo| no tiene algún carácter chino, japonés o coreano, se produce un aviso, una vez por fuente. Con esos paquetes (o |xeCJK|) no se mira: escriben esos caracteres con otra fuente.
\end{key}

\begin{key}{/pgfgo/sgf = \meta{sgf options}}
	Opciones de todos los |\sgf| del documento, por ejemplo |\pgfgoset{sgf = {variations}}|.
\end{key}

\subsection{Datos y comentarios}
\label{subsec:sgfdatos}

\begin{command}{\sgfinfo\opt{\oarg{options}}\marg{file}\marg{property}}
	Escribe una propiedad de la raíz de la partida, por ejemplo |PB| y |PW| (los jugadores), |RE| (el resultado), |KM| (el komi) o |DT| (la fecha), o con |move| la de una jugada. Si no está, no escribe nada. Se puede usar fuera de |goban|.

	\begin{codeexample}[]
		\sgfinfo{variantes}{PB} (negras) contra \sgfinfo{variantes}{PW} (blancas), komi \sgfinfo{variantes}{KM}. Resultado: \sgfinfo{variantes}{RE}.
	\end{codeexample}

	El valor se escribe tal como está en el archivo (con los caracteres escapados ya resueltos y en la codificación de |CA|): |RE| queda como |B+R|, no en palabras. Las propiedades con datos de la partida, según el formato SGF, son:

	\begin{center}
		\small
		\begin{tabular}{lp{6.5cm}l}
			\emph{Propiedad} & \emph{Qué es} & \emph{Ejemplo} \\[2pt]
			|PB|, |PW| & El nombre de Negras y el de Blancas. & |Lee Sedol| \\
			|BR|, |WR| & El rango de Negras y el de Blancas. & |9p| \\
			|BT|, |WT| & El equipo de Negras y el de Blancas. & \\
			|RE| & El resultado: |B+| o |W+| y los puntos, |R| (abandono), |T| (tiempo) o |F| (falta); |0| o |Draw|, empate; |Void|, sin resultado; |?|, desconocido. & |W+3.5| \\
			|KM| & El komi. & |6.5| \\
			|HA| & El hándicap: cuántas piedras. & |2| \\
			|DT| & La fecha, |AAAA-MM-DD| (o varias, separadas por comas). & |2016-03-09| \\
			|EV| & El torneo o evento. & |Meijin| \\
			|RO| & La ronda. & |3 (final)| \\
			|PC| & El lugar. & |Tokio| \\
			|GN| & El nombre de la partida. & \\
			|GC| & Un comentario sobre la partida: los antecedentes, los jugadores. & \\
			|ON| & La apertura (fuseki). & |chinese| \\
			|RU| & Las reglas. & |Japanese| \\
			|TM| & El tiempo de cada jugador, en segundos. & |7200| \\
			|OT| & El tiempo extra (byo-yomi). & |5x30 byo-yomi| \\
			|SO| & La fuente: un libro, una revista. & \\
			|AN| & Quién comentó la partida. & \\
			|US| & Quién escribió el archivo. & \\
			|CP| & El copyright. & \\
			|AP| & El programa que escribió el archivo y su versión. & |CGoban:3| \\
			|SZ| & El tamaño del tablero (|19|, o |19:13| si no es cuadrado). & |19| \\
			|CA| & La codificación del archivo. & |UTF-8| \\
			|FF|, |GM| & La versión del formato y el juego (|1|: go). & |4| \\
		\end{tabular}
	\end{center}

	Sirve cualquier otra propiedad de la raíz (o de la jugada, con |move|). Si tiene varios valores, se escriben separados por comas: |\sgfinfo{partida}{AB}| escribe las piedras negras puestas antes de empezar en coordenadas SGF, |dd, pp|. Los nombres largos de las versiones viejas del formato (|PlayerBlack|) se leen como su abreviatura: se piden con |PB|.
\end{command}

\begin{command}{\sgfcomment\opt{\oarg{options}}\marg{file}\marg{move}}
	Escribe el comentario de una jugada del archivo (|0|: el de antes de la primera jugada): el nombre del nodo (|N|), en negrita; las anotaciones de la jugada (|TE|, |BM|, |DO|, |IT|) y de la posición (|GB|, |GW|, |DM|, |UC|); y el comentario (|C|). Las anotaciones se escriben con los textos de |/pgfgo/sgf annotations|. Los valores (|V|) y los puntos calientes (|HO|), que sirven para el árbol de la partida, no se escriben. Se puede usar fuera de |goban|. |\sgfanimation| y |\sgfdiagrams| escriben los comentarios de la misma forma.

	\begin{codeexample}[]
		Jugada 7: \sgfcomment{partida}{7}
	\end{codeexample}
\end{command}

\begin{key}{/pgfgo/sgf info/game = \meta{number} (initially 1)}
	Como |/pgfgo/sgf/game|, para |\sgfinfo|.
\end{key}

\begin{key}{/pgfgo/sgf info/move = \meta{number} (initially 0)}
	La propiedad de esa jugada: la de su nodo y la de los nodos sin jugada que lo siguen. Con |0|, la de la raíz (y los nodos anteriores a la primera jugada). Sirve, por ejemplo, para el tiempo que le queda a cada jugador en las partidas de KGS y OGS, |BL| y |WL| (en segundos): |\sgfinfo[move = 120]{partida}{WL}|.
\end{key}

\begin{key}{/pgfgo/sgf info/variation = \meta{números} (initially la línea principal)}
	Con |move|, la línea de la jugada, como |/pgfgo/sgf/variation|.
\end{key}

\begin{key}{/pgfgo/sgf comment/game = \meta{number} (initially 1)}
	Como |/pgfgo/sgf/game|.
\end{key}

\begin{key}{/pgfgo/sgf comment/variation = \meta{números} (initially la línea principal)}
	Como |/pgfgo/sgf/variation|.

	\begin{codeexample}[]
		\sgfcomment[variation = 2]{problema}{1}
	\end{codeexample}
\end{key}

\begin{key}{/pgfgo/sgf comment/annotations = \meta{boolean} (default true, initially true)}
	Escribe el nombre y las anotaciones. Con |false|, solo el comentario.

	\begin{codeexample}[]
		Con: \sgfcomment{partida}{7}
		\quad Sin: \sgfcomment[annotations = false]{partida}{7}
	\end{codeexample}
\end{key}

\begin{key}{/pgfgo/sgf comment/chat = \meta{boolean} (default true, initially true)}
	Con |false|, quita del comentario el chat de los servidores: en OGS, todo lo que sigue a la línea «|-- chat --|» (lo de antes es el comentario del que revisa la partida); en KGS, las líneas «|nombre [rango]: ...|» y, si hay alguna, también las «|nombre: ...|» del mismo comentario (las de los que no tienen rango). Si no queda nada, la jugada no tiene comentario.
\end{key}

Los textos de las anotaciones están en el idioma del documento (pág.~\pageref{subsec:idiomas}), y pueden cambiarse para todo el documento con |\pgfgoset{sgf annotations = {...}}|, por ejemplo |bad move = Mala jugada|. Las anotaciones enfatizadas de |DM| y |UC| usan el mismo texto.

\begin{key}{/pgfgo/sgf annotations/tesuji = \meta{text} (initially Tesuji o su traducción)}
	Texto de |TE|.

	\begin{codeexample}[]
		\pgfgoset{sgf annotations = {tesuji = Buena jugada}}
		\sgfcomment{partida}{7}
	\end{codeexample}
\end{key}

\begin{key}{/pgfgo/sgf annotations/very good tesuji = \meta{text} (initially Very good tesuji o su traducción)}
	Texto de |TE[2]|.
\end{key}

\begin{key}{/pgfgo/sgf annotations/bad move = \meta{text} (initially Bad move o su traducción)}
	Texto de |BM|.
\end{key}

\begin{key}{/pgfgo/sgf annotations/very bad move = \meta{text} (initially Very bad move o su traducción)}
	Texto de |BM[2]|.
\end{key}

\begin{key}{/pgfgo/sgf annotations/doubtful = \meta{text} (initially Doubtful move o su traducción)}
	Texto de |DO|.
\end{key}

\begin{key}{/pgfgo/sgf annotations/interesting = \meta{text} (initially Interesting move o su traducción)}
	Texto de |IT|.
\end{key}

\begin{key}{/pgfgo/sgf annotations/good for black = \meta{text} (initially Good for Black o su traducción)}
	Texto de |GB|.
\end{key}

\begin{key}{/pgfgo/sgf annotations/very good for black = \meta{text} (initially Very good for Black o su traducción)}
	Texto de |GB[2]|.
\end{key}

\begin{key}{/pgfgo/sgf annotations/good for white = \meta{text} (initially Good for White o su traducción)}
	Texto de |GW|.
\end{key}

\begin{key}{/pgfgo/sgf annotations/very good for white = \meta{text} (initially Very good for White o su traducción)}
	Texto de |GW[2]|.
\end{key}

\begin{key}{/pgfgo/sgf annotations/even = \meta{text} (initially Even position o su traducción)}
	Texto de |DM|.
\end{key}

\begin{key}{/pgfgo/sgf annotations/unclear = \meta{text} (initially Unclear position o su traducción)}
	Texto de |UC|.
\end{key}

\subsection{Animaciones}
\label{subsec:sgfanimacion}

\begin{command}{\sgfanimation\opt{\oarg{options}}\marg{file}}

	Recorre la partida con el paquete |animate|: dibuja un |goban| por jugada, y los botones de |animate| permiten avanzar, retroceder o reproducirla. El cuadro $k$ muestra la posición después de la jugada $k$; el cuadro 0, la posición antes de la primera jugada. Se usa fuera de |goban|. Los tableros tienen el tamaño de la partida (|SZ|), salvo que |goban| indique otro. Necesita |\usepackage{animate}| y dos compilaciones.
\end{command}

Las animaciones de |animate| funcionan en Acrobat Reader y en algunos otros visores (Foxit, PDF-XChange). En los demás se ve solo el primer cuadro. Estas opciones pueden fijarse para todo el documento con |\pgfgoset{sgf animation = {...}}|.

\begin{codeexample}[width=5cm]
	\sgfanimation[comments, goban = {captures}]{partida}
\end{codeexample}

\begin{key}{/pgfgo/sgf animation/from = \meta{number} (initially 0)}
	Primer cuadro.
\end{key}

\begin{key}{/pgfgo/sgf animation/to = \meta{number} o |last| (initially last)}
	Último cuadro. Con |last| (o |-1|), el de la última jugada.
\end{key}

\begin{key}{/pgfgo/sgf animation/fps = \meta{number} (initially 1)}
	Cuadros por segundo al reproducir.
\end{key}

\begin{key}{/pgfgo/sgf animation/goban = \meta{goban options}}
	Opciones de cada tablero, como en |\begin{goban}[...]|.
\end{key}

\begin{key}{/pgfgo/sgf animation/animate = \meta{animate options} (initially controls)}
	Opciones del entorno |animateinline| de |animate|, por ejemplo |controls, loop|.
\end{key}

\begin{key}{/pgfgo/sgf animation/numbers = \meta{boolean} (default true, initially false)}
	Numera todas las jugadas, como |\sgf|. Sin |numbers|, las jugadas no tienen número y se marca la última.

	\begin{codeexample}[width=5cm]
		\sgfanimation[from = 6, to = 8, numbers, goban = {captures}]{partida}
	\end{codeexample}
\end{key}

\begin{key}{/pgfgo/sgf animation/last mark = \meta{mark options o false} (initially circle)}
	Marca de la última jugada, con el color de quien jugó. Con |false|, la última jugada no se marca. No se usa con |numbers|, ni cuando el archivo ya marca ese punto.

	\begin{codeexample}[width=5cm]
		\sgfanimation[from = 5, to = 7, last mark = triangle, fps = 2,
			animate = {controls, loop}, goban = {captures}]{partida}
	\end{codeexample}
\end{key}

\begin{key}{/pgfgo/sgf animation/comments = \meta{comments options} (default true)}
	Muestra debajo del tablero el comentario (|C|) de la jugada. El del cuadro 0 es el de los nodos anteriores a la primera jugada. Se prende con |comments| o |comments = true| y se apaga con |false| (es el valor inicial); |height| y |font| no la prenden: |comments = {true, height = 6em}|.
\end{key}

\begin{stylekey}{/pgfgo/sgf animation/comments/true}
	Muestra los comentarios.
\end{stylekey}

\begin{stylekey}{/pgfgo/sgf animation/comments/false}
	No muestra los comentarios. Es el valor inicial.
\end{stylekey}

\begin{key}{/pgfgo/sgf animation/comments/height = \meta{dimension} (initially 4em)}
	Alto del espacio para los comentarios. Es el mismo en todos los cuadros, para que no cambien de tamaño.
\end{key}

\begin{key}{/pgfgo/sgf animation/comments/font = \meta{font commands} (initially \texttt{\string\small})}
	Fuente de los comentarios.

	\begin{codeexample}[width=5cm]
		\sgfanimation[from = 7, to = 9, goban = {captures},
			comments = {true, height = 2.5em, font = \itshape}]{partida}
	\end{codeexample}
\end{key}

\begin{key}{/pgfgo/sgf animation/variation = \meta{números} (initially la línea principal)}
	Como |/pgfgo/sgf/variation|: la animación recorre la línea que lleva a esa variante.

	\begin{codeexample}[width=5cm]
		\sgfanimation[variation = 2, numbers]{variantes}
	\end{codeexample}
\end{key}

\begin{key}{/pgfgo/sgf animation/marks = \meta{boolean} (default true, initially true)}
	Como |/pgfgo/sgf/marks|: cada cuadro muestra las marcas de su jugada.
\end{key}

\begin{key}{/pgfgo/sgf animation/step = \meta{number} o \meta{list of numbers} (initially 100)}
	Como |/pgfgo/sgf/step|.
\end{key}

\begin{key}{/pgfgo/sgf animation/restart = \meta{restart options} (default true)}
	Como |/pgfgo/sgf/restart|.
\end{key}

\begin{key}{/pgfgo/sgf animation/restart/true}
	Como |/pgfgo/sgf/restart/true|.
\end{key}

\begin{key}{/pgfgo/sgf animation/restart/false}
	Como |/pgfgo/sgf/restart/false|.
\end{key}

\begin{key}{/pgfgo/sgf animation/restart/auto}
	Como |/pgfgo/sgf/restart/auto|. Es el valor inicial.
\end{key}

\begin{key}{/pgfgo/sgf animation/restart/from = \meta{number} (initially 1)}
	Como |/pgfgo/sgf/restart/from|.
\end{key}

\begin{key}{/pgfgo/sgf animation/annotations = \meta{boolean} (default true, initially true)}
	Con |comments|, escribe también el nombre y las anotaciones de cada jugada, como |\sgfcomment|.
\end{key}

\begin{key}{/pgfgo/sgf animation/chat = \meta{boolean} (default true, initially true)}
	Con |comments|, como |/pgfgo/sgf comment/chat|: con |false|, sin el chat de KGS y OGS.
\end{key}

\begin{key}{/pgfgo/sgf animation/game = \meta{number} (initially 1)}
	Como |/pgfgo/sgf/game|.
\end{key}

\begin{key}{/pgfgo/sgf animation/view = \meta{boolean} (default true, initially true)}
	Si el archivo tiene vista (|VW|), los tableros muestran solo esa parte, como |sgf view|. Todos los cuadros tienen el mismo tablero: si la vista cambia durante la partida, el tablero abarca las vistas de todos los cuadros (y si en alguno se ve todo el tablero, se ve todo).
\end{key}

\begin{key}{/pgfgo/sgf animation/layers = \meta{boolean} (default true, initially true)}
	Arma la animación por capas: el tablero se dibuja una sola vez, y cada piedra también, en una lámina que |animate| muestra mientras la piedra sigue en el tablero (con la opción |timeline| de |animate|). Lo único que se dibuja en cada cuadro es lo que cambia: la marca de la última jugada, las marcas del archivo y el comentario. Así, el PDF crece con la cantidad de jugadas y no con su cuadrado: una partida de 200 jugadas ocupa menos de 1~MB, en lugar de más de 10~MB, y además compila más rápido. Con |layers = false|, cada cuadro es un |goban| entero. El |timeline| se guarda en un archivo auxiliar, |<jobname>-pgfgo-<n>.tln|.
\end{key}

\subsection{Diagramas como en los libros}
\label{subsec:sgfdiagramas}

\begin{command}{\sgfdiagrams\opt{\oarg{options}}\marg{file}}
	Dibuja la partida como en los libros: un diagrama cada |step| jugadas (1--50, 51--100, etc.). Cada diagrama parte de la posición real, con las jugadas anteriores sin número y sus capturas; las piedras capturadas dentro del diagrama no se quitan (se quitan en el siguiente). Debajo de cada uno van su epígrafe y sus notas «N en M» (pág.~\pageref{subsec:notas}). Se usa fuera de |goban|. Las opciones pueden fijarse para todo el documento con |\pgfgoset{sgf diagrams = {...}}|.

	\begin{codeexample}[code only]
		\sgfdiagrams[step = 5, title = Diagrama]{partida}
	\end{codeexample}
\end{command}

\sgfdiagrams[step = 5, title = Diagrama]{partida}

\begin{key}{/pgfgo/sgf diagrams/step = \meta{number} o \meta{list of numbers} (initially 50)}
	Cantidad de jugadas de cada diagrama. Con una lista, el primer diagrama tiene tantas jugadas como dice el primer número, el segundo como dice el segundo, y así; el último número se repite: con |step = {5, 10, 3}|, los diagramas tienen 5, 10, 3, 3, 3... jugadas.
\end{key}

\begin{key}{/pgfgo/sgf diagrams/from = \meta{number} (initially 1)}
	Primera jugada. Los diagramas son los mismos que sin |from|, con el mismo número: solo se dibujan los que tienen jugadas desde |from|, y el primero empieza en |from|. Con |step = 50| y |from = 120|, el primero es «Diagrama 3 (120--150)».
\end{key}

\begin{key}{/pgfgo/sgf diagrams/to = \meta{number} o |last| (initially last)}
	Última jugada. Con |last| (o |-1|), la última de la partida.
\end{key}

\begin{key}{/pgfgo/sgf diagrams/numbers = \meta{boolean} (default true, initially true)}
	Con |false|, las jugadas de cada diagrama se dibujan sin número (las notas «N at M» siguen usando los números de la partida).
\end{key}

\begin{key}{/pgfgo/sgf diagrams/last mark = \meta{mark options} o |false| (initially false)}
	Como |/pgfgo/sgf/last mark|, en cada diagrama.

	\begin{codeexample}[code only]
		\sgfdiagrams[step = 5, to = 10, numbers = false, last mark = circle]{partida}
	\end{codeexample}

	\sgfdiagrams[step = 5, to = 10, numbers = false, last mark = circle]{partida}
\end{key}

\begin{key}{/pgfgo/sgf diagrams/restart = \meta{restart options} (default true)}
	Como |/pgfgo/sgf/restart|, con un diagrama por tramo: desde el segundo diagrama, cada uno numera desde 1 (o desde |restart/from|), como algunos libros. El primer diagrama conserva los números de la partida.

	\begin{codeexample}[code only]
		\sgfdiagrams[step = {2, 3}, restart = {from = 0}]{partida}
	\end{codeexample}
\end{key}

\sgfdiagrams[step = {2, 3}, to = 8, restart = {from = 0}, title = Diagrama]{partida}

\begin{key}{/pgfgo/sgf diagrams/restart/true}
	Reinicia la numeración en cada diagrama.
\end{key}

\begin{key}{/pgfgo/sgf diagrams/restart/false}
	No la reinicia: las jugadas llevan los números de la partida.
\end{key}

\begin{key}{/pgfgo/sgf diagrams/restart/auto}
	Como |/pgfgo/sgf/restart/auto|: con |PM[2]|, cada 100 jugadas, sin importar los diagramas. Es el valor inicial.
\end{key}

\begin{key}{/pgfgo/sgf diagrams/restart/from = \meta{number} (initially 1)}
	El número de la primera jugada de cada diagrama, desde el segundo. También prende |restart|. Con |step = 99| y |restart = {from = 0}|, el primer diagrama va de 1 a 99, y en el segundo la jugada 100 se muestra como 0, la 101 como 1, y así.
\end{key}

\begin{key}{/pgfgo/sgf diagrams/notes = \meta{boolean} (default true, initially true)}
	Escribe debajo de cada diagrama sus notas: «N en M» y los pases, «N: pasa».
\end{key}

\begin{key}{/pgfgo/sgf diagrams/title = \meta{text} (initially Diagram o su traducción)}
	Palabra del epígrafe: «Diagrama 2 (51--100)».
\end{key}

\begin{key}{/pgfgo/sgf diagrams/caption = \meta{caption options}}
	Opciones del epígrafe y de las notas.
\end{key}

\begin{key}{/pgfgo/sgf diagrams/caption/font = \meta{font commands} (initially \texttt{\string\small})}
	Fuente del epígrafe y de las notas: |caption = {font = \itshape}|.

	\begin{codeexample}[code only]
		\sgfdiagrams[from = 6, title = Figura, caption = {font = \itshape},
			goban = {label = {at = {north, west}}}]{partida}
	\end{codeexample}

	\sgfdiagrams[from = 6, title = Figura, caption = {font = \itshape},
		goban = {label = {at = {north, west}}}]{partida}
\end{key}

\begin{key}{/pgfgo/sgf diagrams/caption/min width = \meta{dimension} (initially 8em)}
	Con |layout = {columns = auto}|, un epígrafe más ancho que su tablero se corta en líneas al ancho del tablero. Si el tablero es más angosto que esta medida (en la fuente del epígrafe), se corta a esta medida; nunca se corta más angosto que su palabra más larga, ni se ensancha más allá de su ancho en una línea. En la grilla, el epígrafe se corta al ancho de la celda.
\end{key}

\begin{key}{/pgfgo/sgf diagrams/goban = \meta{goban options}}
	Opciones de cada tablero. Por defecto, los tableros tienen el tamaño de la partida y usan |captures| y |notes|.
\end{key}

\begin{key}{/pgfgo/sgf diagrams/layout = \meta{layout options}}
	Cómo se disponen los diagramas: opciones de |/pgfgo/layout| (pág.~\pageref{subsec:disposicion}) solo para |\sgfdiagrams|. Por ejemplo, |layout = {columns = 3}|.
\end{key}

\begin{key}{/pgfgo/sgf diagrams/variation = \meta{números} (initially la línea principal)}
	Como |/pgfgo/sgf/variation|.
\end{key}

\begin{key}{/pgfgo/sgf diagrams/variations = \meta{variations options} (default commented)}
	Dibuja también las variantes de la partida, cada una en su diagrama, después del diagrama donde nace su primera jugada: «Diagrama 3a», «Diagrama 3b»... Las jugadas de cada variante se numeran desde 1, y el epígrafe dice qué jugada reemplaza la primera (|variation title|). Las anteriores se ven sin número, como estaban. Las opciones dicen cuáles se dibujan, y se suman: |variations = {commented, 5}|. Además de las de abajo, cualquier otra cosa es el camino de una variante, como en |variation| (|2|, |2.3.1|): esa se dibuja siempre, a cualquier profundidad. Los |1| del final no cambian la línea: |2.3.1| es la misma que |2.3|. Si el archivo no la tiene, se produce un aviso. Sus comentarios y notas se escriben como los de los demás diagramas. Solo con la línea principal (sin |variation|).

	\begin{codeexample}[]
		\sgfdiagrams[step = 2, to = 3, variations = all, goban = {scale = 0.5}]{variantes}
	\end{codeexample}
\end{key}

\begin{stylekey}{/pgfgo/sgf diagrams/variations/none}
	Ninguna. Es el valor inicial, y borra lo anterior.
\end{stylekey}

\begin{stylekey}{/pgfgo/sgf diagrams/variations/commented}
	Las que tienen algún comentario (|C|) en sus jugadas propias. Es lo que hace |variations| a secas: en las partidas comentadas, las demás suelen ser exploraciones sin explicación.
\end{stylekey}

\begin{stylekey}{/pgfgo/sgf diagrams/variations/all}
	Todas.
\end{stylekey}

\begin{key}{/pgfgo/sgf diagrams/variations/depth = \meta{number} (initially 1)}
	Hasta qué profundidad dibujan |commented| y |all|: |1|, solo las que salen de la línea principal; |2|, también las variantes de esas, y así. Las de una variante van después de ella, y su epígrafe dice qué jugada de ella reemplazan.
\end{key}

\begin{key}{/pgfgo/sgf diagrams/variation title = \meta{text} (initially «variation at» y el número, o su traducción)}
	Lo que va entre paréntesis en el epígrafe de una variante: «Diagrama 3a (variante en 26)». |\pgfgovariationmove| es el número de la jugada que reemplaza.
\end{key}

\begin{key}{/pgfgo/sgf diagrams/comments = \meta{comments options} (default true)}
	Escribe debajo de cada diagrama los comentarios de sus jugadas, como |\sgfcomment|, cada uno con el número de la jugada. Se prende con |comments| o |comments = true| y se apaga con |false| (es el valor inicial); las demás opciones no la prenden: |comments = {true, align = justify}|. Así, |\pgfgoset{sgf diagrams = {comments = {font = \itshape}}}| fija el estilo de todo el documento sin prender los comentarios.
\end{key}

\begin{stylekey}{/pgfgo/sgf diagrams/comments/true}
	Escribe los comentarios.
\end{stylekey}

\begin{stylekey}{/pgfgo/sgf diagrams/comments/false}
	No escribe los comentarios. Es el valor inicial.
\end{stylekey}

\begin{key}{/pgfgo/sgf diagrams/comments/font = \meta{font commands} (initially la del epígrafe)}
	Fuente de los comentarios. Sin esta opción, la de |caption/font|.
\end{key}

\begin{key}{/pgfgo/sgf diagrams/comments/align = \meta{left, right, center o justify} (initially left)}
	Alineación de los comentarios en el ancho del diagrama. Con |justify|, en un diagrama angosto quedan espacios grandes entre las palabras.
\end{key}

\begin{key}{/pgfgo/sgf diagrams/comments/root = \meta{boolean} (default true, initially false)}
	Escribe también el comentario de la raíz (el de los nodos anteriores a la primera jugada), sin número, antes de los del diagrama que empieza en la jugada 1.

	\begin{codeexample}[code only]
		\sgfdiagrams[step = 10, comments = {true, root, font = \itshape}]{partida}
	\end{codeexample}
\end{key}

\begin{key}{/pgfgo/sgf diagrams/comments/number = \meta{text} (initially \texttt{\string\textbf\{\#1\}:\string~})}
	Lo que va delante de cada comentario; |#1| es el número de la jugada, el que se ve en el tablero (con |restart|, el reiniciado). Por ejemplo, |number = {Jugada~#1.~}|.
\end{key}

\begin{key}{/pgfgo/sgf diagrams/comments/separator = \meta{text} (initially un párrafo por comentario)}
	Lo que va entre los comentarios. Con el valor inicial, cada uno va en su párrafo; con |separator = {\quad}|, van seguidos.

	\begin{codeexample}[code only]
		\sgfdiagrams[step = 10, comments = {true, number = {\textbf{#1}.~}, separator = {~}}]{partida}
	\end{codeexample}
\end{key}

\begin{key}{/pgfgo/sgf diagrams/comments/position = \meta{below, right, left o after} (initially below)}
	Dónde van los comentarios: debajo del epígrafe, al costado del diagrama, como en los libros de partidas comentadas, o con |after| todos juntos después de los diagramas, a lo ancho de la página, cada grupo con el epígrafe de su diagrama como encabezado (se pueden cortar entre páginas). Al costado empiezan a la altura del borde de arriba del tablero; el epígrafe, las notas y el resultado siguen debajo del tablero.

	\begin{codeexample}[code only]
		\sgfdiagrams[step = 10, comments = {true, position = right}]{partida}
	\end{codeexample}
\end{key}

\begin{key}{/pgfgo/sgf diagrams/comments/width = \meta{dimension} o auto (initially auto)}
	Ancho de los comentarios al costado. Con |auto|, lo que queda: fuera de la grilla, el resto de la línea (cada diagrama va en su línea); en la grilla (|layout = {columns = ...}|), el resto de la celda, y con |fit| el tablero ocupa la mitad de la celda y los comentarios la otra mitad. Con una dimensión, en la grilla el tablero ocupa lo que queda de la celda.
\end{key}

\begin{key}{/pgfgo/sgf diagrams/comments/sep = \meta{dimension} (initially 1em)}
	Espacio entre el diagrama y los comentarios al costado.

	\begin{codeexample}[code only]
		\sgfdiagrams[step = 10, layout = {columns = 2, fit},
			comments = {true, position = left, sep = 2em}]{partida}
	\end{codeexample}
\end{key}

\begin{key}{/pgfgo/sgf diagrams/comments/height = \meta{dimension} o none (initially none)}
	Altura máxima de los comentarios de un diagrama, debajo o al costado. Lo que no entra va después de los diagramas, como con |after|, con el epígrafe del diagrama como encabezado. Sirve para los comentarios muy largos, que estiran la fila de la grilla.

	\begin{codeexample}[code only]
		\sgfdiagrams[step = 30, layout = {columns = 2},
			comments = {true, height = 6\baselineskip}]{partida}
	\end{codeexample}
\end{key}

\begin{key}{/pgfgo/sgf diagrams/annotations = \meta{boolean} (default true, initially true)}
	Con |comments|, escribe también el nombre y las anotaciones de cada jugada.

	\begin{codeexample}[code only]
		\sgfdiagrams[step = 10, comments]{partida}
		\sgfdiagrams[step = 10, comments, annotations = false]{partida}
	\end{codeexample}

	\sgfdiagrams[step = 10, comments]{partida}
	\sgfdiagrams[step = 10, comments, annotations = false]{partida}
\end{key}

\begin{key}{/pgfgo/sgf diagrams/chat = \meta{boolean} (default true, initially true)}
	Con |comments|, como |/pgfgo/sgf comment/chat|: con |false|, sin el chat de KGS y OGS. Una jugada a la que solo le quedaba chat no aparece.
\end{key}

\begin{key}{/pgfgo/sgf diagrams/figures = \meta{boolean} (default true, initially true)}
	Si el archivo marca figuras (|FG|), cada figura es un diagrama, con su nombre como epígrafe, en lugar de uno cada |step| jugadas.
\end{key}

\begin{key}{/pgfgo/sgf diagrams/view = \meta{boolean} (default true, initially true)}
	Si el archivo tiene vista (|VW|), cada diagrama muestra solo esa parte, como |sgf view|. Si la vista cambia durante la partida, cada diagrama abarca las vistas de sus jugadas.
\end{key}

\begin{key}{/pgfgo/sgf diagrams/game = \meta{number} (initially 1)}
	Como |/pgfgo/sgf/game|.
\end{key}

\begin{key}{/pgfgo/sgf diagrams/result = \meta{result options} (default true)}
	Escribe debajo del último diagrama, si llega al final de la partida, el resultado (|RE|) y, si el archivo marca el territorio (|TB|, |TW|), el puntaje de cada jugador: «Resultado: Ganan Negras por 0.5 puntos (Negras 6, Blancas 5.5)». El puntaje se cuenta como con |\pgfgoscore| (pág.~\pageref{subsec:territorio}), con el komi del archivo (|KM|). Se apaga con |result = false|; |title| y |format| no lo prenden.
\end{key}

\begin{stylekey}{/pgfgo/sgf diagrams/result/true}
	Escribe el resultado. Es el valor inicial.
\end{stylekey}

\begin{stylekey}{/pgfgo/sgf diagrams/result/false}
	No escribe el resultado.
\end{stylekey}

\begin{key}{/pgfgo/sgf diagrams/result/title = \meta{text} (initially Result o su traducción)}
	Palabra que precede al resultado.
\end{key}

\begin{key}{/pgfgo/sgf diagrams/result/format = \meta{text o raw} (initially text)}
	Cómo se escribe el resultado. |text|: en palabras, en el idioma del documento (pág.~\pageref{subsec:idiomas}), «Ganan Negras por abandono»; |raw|: como está en el archivo, «B+R». Con |text| se entienden |0| o |Draw| (empate), |Void| (sin resultado), |?| (desconocido) y |B+| o |W+| seguido de nada, de la diferencia, de |R| o |Resign|, de |T| o |Time|, o de |F| o |Forfeit|. Otro resultado se escribe como está.
\end{key}

\begin{key}{/pgfgo/sgf diagrams/black name = \meta{text} (initially Black o su traducción)}
	Nombre de las negras en el puntaje y en el resultado.
\end{key}

\begin{key}{/pgfgo/sgf diagrams/white name = \meta{text} (initially White o su traducción)}
	Nombre de las blancas en el puntaje y en el resultado.

	\begin{codeexample}[code only]
		\sgfdiagrams[result = {title = Final, format = raw}]{variantes}
		\sgfdiagrams[black name = Ana, white name = Bruno]{variantes}
	\end{codeexample}

	\sgfdiagrams[result = {title = Final, format = raw}]{variantes}
	\sgfdiagrams[black name = Ana, white name = Bruno]{variantes}
\end{key}

\begin{command}{\pgfgowinner}
	En los textos del resultado (|\pgfgodefinelanguage|), el nombre del ganador: |black name| o |white name|.
\end{command}

\begin{command}{\pgfgovariationmove}
	En el epígrafe de una variante (|variation title|) y en su texto en cada idioma (|variation at|), el número de la jugada que reemplaza.
\end{command}

\begin{command}{\pgfgopoints}
	En los textos del resultado, la diferencia de puntos, como está en el archivo.
\end{command}

\subsection{Problemas}
\label{subsec:sgfproblemas}

\begin{command}{\sgfproblem\opt{\oarg{options}}\marg{file}}
	Dibuja un problema (tsumego). Primero, el enunciado: la posición inicial con sus marcas, quién juega (|PL|, o el color de la primera jugada) y el comentario del principio. Después, un diagrama por cada línea del árbol del archivo (cada variante de la solución, correcta o no), numerada desde 1, con sus notas «N en M» y los comentarios y anotaciones de sus jugadas, que suelen decir si la variante es correcta. Los diagramas se dibujan como los de |\sgfdiagrams|. Si el archivo no tiene vista (|VW|), se muestra solo la parte del tablero que usa el problema. Se usa fuera de |goban|. Las opciones pueden fijarse para todo el documento con |\pgfgoset{sgf problem = {...}}|. En los ejemplos, |problema| es |problema.sgf|, un problema de vida y muerte: la jugada correcta y dos equivocadas, con sus comentarios.

	\begin{codeexample}[code only]
		\sgfproblem[title = Problema 1, variation title = Variante,
			black to play = Juegan negras]{problema}
	\end{codeexample}
\end{command}

\sgfproblem[title = Problema 1, variation title = Variante,
	black to play = Juegan negras]{problema}

\begin{key}{/pgfgo/sgf problem/statement = \meta{true, false o alone} (default true, initially true)}
	Dibuja el enunciado. Con |alone|, el enunciado va solo en su fila, centrado, y las variantes debajo.
\end{key}

\begin{key}{/pgfgo/sgf problem/solutions = \meta{boolean} (default true, initially true)}
	Dibuja las variantes. Con |solutions = false|, queda solo el enunciado, para poner las soluciones en otro lugar (por ejemplo, con |statement = false| al final del capítulo).

	\begin{codeexample}[code only]
		\sgfproblem[solutions = false, title = Problema 2]{problema}
	\end{codeexample}

	\sgfproblem[solutions = false, title = Problema 2]{problema}
\end{key}

\begin{key}{/pgfgo/sgf problem/title = \meta{text} (initially vacío)}
	Comienzo del epígrafe del enunciado: «Problema 1: Juegan negras». Con varias partidas (|game = 1:20|), va seguido del número de cada una; vacío, es |Problem| o su traducción: «Problema 7: Juegan negras». Las soluciones dicen también de qué problema son: «Problema 7: Variante 1».
\end{key}

\begin{key}{/pgfgo/sgf problem/to play = \meta{auto, black, white o none} (initially auto)}
	Quién juega, para el epígrafe del enunciado. Con |auto|, lo que dice |PL| o, si no está, el color de la primera jugada; si no hay ninguna de las dos, no se dice. Muchas colecciones no lo dicen en el archivo: |to play = black|. Con |none|, el epígrafe no lo dice.
\end{key}

\begin{key}{/pgfgo/sgf problem/verdict = \meta{boolean} (default true, initially false)}
	El epígrafe de cada solución dice si es correcta o incorrecta («Variante 2 (incorrecta)»), con los textos |correct| y |wrong| del idioma. Se decide con las anotaciones de la línea: el último |GB| o |GW| (la posición es buena para negras o para blancas), según quién juega; si no hay, la primera jugada con |TE| o |BM| (las jugadas impares son de quien juega: un |BM| suyo es un error, uno del rival no); si no, la palabra |RIGHT| en un comentario, como en los archivos de goproblems.com. Si no hay nada de eso, no dice nada.

	\begin{codeexample}[]
		\sgfproblem[verdict, statement = false, comments = false,
			goban = {scale = 0.5}]{problema}
	\end{codeexample}
\end{key}

\begin{key}{/pgfgo/sgf problem/variation title = \meta{text} (initially Variation o su traducción)}
	Palabra del epígrafe de cada variante: «Variante 2».
\end{key}

\begin{key}{/pgfgo/sgf problem/black to play = \meta{text} (initially Black to play o su traducción)}
	Texto del enunciado cuando juegan las negras.
\end{key}

\begin{key}{/pgfgo/sgf problem/white to play = \meta{text} (initially White to play o su traducción)}
	Texto del enunciado cuando juegan las blancas.
\end{key}

\begin{key}{/pgfgo/sgf problem/comments = \meta{comments options} (default true)}
	Escribe el comentario del principio debajo del enunciado y los de cada variante debajo de su diagrama, como |/pgfgo/sgf diagrams/comments|. Acá están prendidos al empezar; se apagan con |comments = false|. Las opciones son |true|, |false|, |font|, |align|, |number|, |separator|, |position|, |width|, |sep| y |height|, como en |\sgfdiagrams|; |font| y |align| valen también para el comentario del enunciado.
\end{key}

\begin{stylekey}{/pgfgo/sgf problem/comments/true}
	Escribe los comentarios. Es el valor inicial.
\end{stylekey}

\begin{stylekey}{/pgfgo/sgf problem/comments/false}
	No escribe los comentarios.
\end{stylekey}

\begin{key}{/pgfgo/sgf problem/comments/font = \meta{font commands} (initially la del epígrafe)}
	Como |/pgfgo/sgf diagrams/comments/font|.
\end{key}

\begin{key}{/pgfgo/sgf problem/comments/align = \meta{left, right, center o justify} (initially left)}
	Como |/pgfgo/sgf diagrams/comments/align|.
\end{key}

\begin{key}{/pgfgo/sgf problem/comments/number = \meta{text} (initially \texttt{\string\textbf\{\#1\}:\string~})}
	Como |/pgfgo/sgf diagrams/comments/number|.
\end{key}

\begin{key}{/pgfgo/sgf problem/comments/separator = \meta{text} (initially un párrafo por comentario)}
	Como |/pgfgo/sgf diagrams/comments/separator|.
\end{key}

\begin{key}{/pgfgo/sgf problem/comments/position = \meta{below, right, left o after} (initially below)}
	Como |/pgfgo/sgf diagrams/comments/position|.
\end{key}

\begin{key}{/pgfgo/sgf problem/comments/width = \meta{dimension} o auto (initially auto)}
	Como |/pgfgo/sgf diagrams/comments/width|.
\end{key}

\begin{key}{/pgfgo/sgf problem/comments/sep = \meta{dimension} (initially 1em)}
	Como |/pgfgo/sgf diagrams/comments/sep|.
\end{key}

\begin{key}{/pgfgo/sgf problem/comments/height = \meta{dimension} o none (initially none)}
	Como |/pgfgo/sgf diagrams/comments/height|.
\end{key}

\begin{key}{/pgfgo/sgf problem/annotations = \meta{boolean} (default true, initially true)}
	Con |comments|, escribe también el nombre y las anotaciones de cada jugada (|TE|, |BM|, etc.).
\end{key}

\begin{key}{/pgfgo/sgf problem/chat = \meta{boolean} (default true, initially true)}
	Con |comments|, como |/pgfgo/sgf comment/chat|: con |false|, sin el chat de KGS y OGS.
\end{key}

\begin{key}{/pgfgo/sgf problem/notes = \meta{boolean} (default true, initially true)}
	Escribe las notas «N en M» de cada variante.
\end{key}

\begin{key}{/pgfgo/sgf problem/caption = \meta{caption options}}
	Opciones de los epígrafes, las notas y los comentarios.
\end{key}

\begin{key}{/pgfgo/sgf problem/caption/font = \meta{font commands} (initially \texttt{\string\small})}
	Fuente de los epígrafes, las notas y los comentarios.

	\begin{codeexample}[code only]
		\sgfproblem[statement = alone, comments = false, notes = false,
			caption = {font = \itshape}]{problema}
	\end{codeexample}

	\sgfproblem[statement = alone, comments = false, notes = false,
		caption = {font = \itshape}]{problema}
\end{key}

\begin{key}{/pgfgo/sgf problem/caption/min width = \meta{dimension} (initially 8em)}
	Como en |\sgfdiagrams|: con |columns = auto|, el ancho mínimo al que se corta un epígrafe más ancho que su tablero («Problema 3: Juegan Blancas» debajo de una esquina chica).
\end{key}

\begin{key}{/pgfgo/sgf problem/goban = \meta{goban options}}
	Opciones de cada tablero, como |/pgfgo/sgf diagrams/goban|.
\end{key}

\begin{key}{/pgfgo/sgf problem/layout = \meta{layout options}}
	Cómo se disponen el enunciado y las variantes: opciones de |/pgfgo/layout| (pág.~\pageref{subsec:disposicion}) solo para |\sgfproblem|.
\end{key}

\begin{key}{/pgfgo/sgf problem/view = \meta{boolean} (default true, initially true)}
	Si el archivo tiene vista (|VW|) al principio, todos los diagramas muestran esa parte.
\end{key}

\begin{key}{/pgfgo/sgf problem/crop = \meta{boolean} (default true, initially true)}
	Si no se usa la vista del archivo, muestra solo la parte del tablero donde hay piedras, jugadas o marcas en alguna línea del problema, con |margin| líneas más. Un lado que queda a dos líneas o menos del borde llega hasta el borde. Es la misma en todos los diagramas. Con |crop = false|, se ve todo el tablero.
\end{key}

\begin{key}{/pgfgo/sgf problem/margin = \meta{number} (initially 1)}
	Líneas que |crop| agrega alrededor del problema.

	\begin{codeexample}[code only]
		\sgfproblem[solutions = false, margin = 3]{problema}
	\end{codeexample}

	\sgfproblem[solutions = false, margin = 3]{problema}
\end{key}

\begin{key}{/pgfgo/sgf problem/game = \meta{numbers} (initially 1)}
	Como |/pgfgo/sgf/game|: en una colección, el problema que se dibuja. También varios, con números y rangos como los de las coordenadas: |game = 1:20| o |game = {1:4, 9}|. Entonces van primero todos los enunciados, numerados (|title|), y después las soluciones de los que las tienen. Así se arma una hoja de problemas:

	\begin{codeexample}[code only]
		\sgfproblem[game = 1:24, to play = black, comments = false,
			layout = {columns = 6, fit}]{coleccion.sgf}
	\end{codeexample}
\end{key}

\subsection{Disposición de los diagramas}
\label{subsec:disposicion}

|\sgfdiagrams| y |\sgfproblem| dibujan varios diagramas. Por defecto van en el párrafo, uno al lado del otro, y pasan a la línea siguiente cuando no entran. Con |columns|, van en una grilla de columnas del mismo ancho: cada fila es una línea aparte (la página puede cortarse entre filas), los tableros de una fila quedan alineados arriba y una fila incompleta deja vacías las últimas columnas. Con |fit|, además, los tableros se escalan para llenar el ancho de su columna.

\begin{codeexample}[code only]
	\sgfdiagrams[step = 2, to = 6, layout = {columns = 3, fit}]{partida}
\end{codeexample}

\sgfdiagrams[step = 2, to = 6, layout = {columns = 3, fit}]{partida}

\begin{stylekey}{/pgfgo/layout = \meta{layout options}}
	Opciones de la disposición para todo el documento, por ejemplo |\pgfgoset{layout = {columns = 3, fit}}|. Para un solo comando, |/pgfgo/sgf diagrams/layout| y |/pgfgo/sgf problem/layout|; también para todo el documento, |\pgfgoset{sgf diagrams = {layout = {...}}}|.
\end{stylekey}

\begin{key}{/pgfgo/layout/columns = \meta{number o auto} (initially auto)}
	Cantidad de columnas de la grilla. Con |auto|, los diagramas van en el párrafo, sin grilla.
\end{key}

\begin{key}{/pgfgo/layout/width = \meta{dimension} o |natural| (initially \texttt{\string\linewidth})}
	Ancho de la grilla. Si es menor que el de la línea, la grilla se centra. Cada columna mide lo que queda al restarle los |column sep| a |width|, repartido entre las columnas; un tablero más angosto queda en su columna según |align|, así que entre dos tableros hay más espacio que |column sep|.

	Con |natural|, la grilla se ajusta a los diagramas: las columnas tienen el ancho del diagrama más ancho (con su epígrafe, que se corta como sin grilla, ver |caption/min width|) y entre dos tableros del mismo ancho queda exactamente |column sep|. La grilla mide eso por la cantidad de columnas (o de diagramas, si son menos) y se centra en la línea; si es más ancha que la línea, se sale del margen. Los diagramas se dibujan todos juntos al final (o al terminar cada grupo de |per page| o |per frame|). |fit| no sirve con |natural|: el paquete avisa y los tableros quedan con su escala.

	\begin{codeexample}[code only]
		\begin{pgfgolayout}[columns = 2, width = natural, column sep = 0em, row sep = 0em]
			\begin{goban}[partial = A9:E5, scale = 0.5]\stone[b]{C7}\end{goban}
			\begin{goban}[partial = E9:I5, scale = 0.5]\stone[w]{G7}\end{goban}
			\begin{goban}[partial = A1:E5, scale = 0.5]\stone[b]{C3}\end{goban}
			\begin{goban}[partial = E1:I5, scale = 0.5]\stone[w]{G3}\end{goban}
		\end{pgfgolayout}
	\end{codeexample}

	\begin{pgfgolayout}[columns = 2, width = natural, column sep = 0em, row sep = 0em]
		\begin{goban}[partial = A9:E5, scale = 0.5]\stone[b]{C7}\end{goban}
		\begin{goban}[partial = E9:I5, scale = 0.5]\stone[w]{G7}\end{goban}
		\begin{goban}[partial = A1:E5, scale = 0.5]\stone[b]{C3}\end{goban}
		\begin{goban}[partial = E1:I5, scale = 0.5]\stone[w]{G3}\end{goban}
	\end{pgfgolayout}
\end{key}

\begin{key}{/pgfgo/layout/column sep = \meta{dimension} (initially 1em)}
	Espacio entre columnas.
\end{key}

\begin{key}{/pgfgo/layout/row sep = \meta{skip} (initially \texttt{\string\medskipamount})}
	Espacio entre filas.
\end{key}

\begin{key}{/pgfgo/layout/fit = \meta{boolean} (default true, initially false)}
	Escala cada tablero para que llene el ancho de su columna; reemplaza el |scale| del tablero. Las líneas no cambian de grosor. En |\sgfdiagrams| y |\sgfproblem|, todos los tableros usan la escala que llena la columna con el primero, así las piedras tienen el mismo tamaño en todos. Si con esa escala un tablero no entra en la columna (un problema más ancho que el primero, por ejemplo), solo ese tablero se achica hasta llenarla. Solo con |columns|, y no con |width = natural|.
\end{key}

\begin{key}{/pgfgo/layout/per page = \meta{number} o |false| (initially 0)}
	Cada \meta{number} diagramas, empieza una página nueva (con |\newpage|): con |per page = 1|, un diagrama por página. En la grilla, una fila a medias se dibuja antes del corte; conviene que \meta{number} sea un múltiplo de |columns|. Con |0| o |false|, no hay cortes. En |\sgfproblem| con |statement = alone|, el enunciado cuenta como un diagrama.

	\begin{codeexample}[code only]
		\sgfdiagrams[step = 50, layout = {columns = 2, fit, per page = 4}]{lee}
	\end{codeexample}
\end{key}

\begin{key}{/pgfgo/layout/per frame = \meta{number} o |false| (default 1, initially 0)}
	Dentro de un |frame| de |beamer|, cada \meta{number} diagramas pasan a la diapositiva siguiente: con |per frame| sola, un diagrama por diapositiva, cada uno con su epígrafe. Los diagramas usan |\only<+>|, así que se suman a las demás capas del |frame|: después de un |\pause|, el primer diagrama va en la diapositiva siguiente. Fuera de |beamer|, el paquete avisa y la \emph{key} no se usa. En |beamer|, si también hay |per page|, vale |per frame|.

	\begin{codeexample}[code only]
		\begin{frame}{La partida}
			\sgfdiagrams[step = 50, goban = {scale = 0.6}, layout = {per frame}]{lee}
		\end{frame}
	\end{codeexample}
\end{key}

\begin{key}{/pgfgo/layout/align = \meta{center, left o right} (initially center)}
	Dónde va cada diagrama dentro de su columna, con su epígrafe y sus notas. Los comentarios van siempre alineados a la izquierda.

	\begin{codeexample}[code only]
		\sgfdiagrams[step = 2, to = 6, layout = {columns = 2, width = 0.6\linewidth,
			column sep = 3em, row sep = 2em, align = left}]{partida}
	\end{codeexample}

	\sgfdiagrams[step = 2, to = 6, layout = {columns = 2, width = 0.6\linewidth,
		column sep = 3em, row sep = 2em, align = left}]{partida}
\end{key}

\begin{environment}{{pgfgolayout}\opt{\oarg{layout options}}}
	Dispone en una grilla los |goban| que contiene, y los diagramas de los |\sgfdiagrams| y |\sgfproblem| que contiene, en el orden en que aparecen. Así se pueden juntar en una grilla tableros dibujados a mano y diagramas de distintos archivos. Las opciones son las de |/pgfgo/layout|; las opciones |layout| de los comandos de adentro no se usan. Adentro no debe haber texto, solo los tableros y los comandos. Con |fit|, cada |goban| se dibuja varias veces para medirlo: su contenido no debe tener efectos fuera del tablero (por ejemplo, contadores que aumentan).

	\begin{codeexample}[code only]
		\begin{pgfgolayout}[columns = 4, fit]
			\begin{goban}[5]\stones[b]{C3}\end{goban}
			\begin{goban}[5]\stones[w]{C3}\end{goban}
			\sgfdiagrams[step = 5]{partida}
		\end{pgfgolayout}
	\end{codeexample}
\end{environment}

\begin{pgfgolayout}[columns = 4, fit]
	\begin{goban}[5]\stones[b]{C3}\end{goban}
	\begin{goban}[5]\stones[w]{C3}\end{goban}
	\sgfdiagrams[step = 5]{partida}
\end{pgfgolayout}

	\section{Jugadores y perfiles}
\label{sec:jugadores}

Cada piedra, jugada y marca pertenece a uno de tres jugadores: el primero (|black|), el segundo (|white|) o el neutral (|neutral|). Su aspecto no está en cada comando sino en el jugador: el árbol de \emph{keys} de cada uno reúne las opciones de sus piedras, sus jugadas, sus marcas y su territorio, que en las secciones anteriores aparecen bajo |stone|, |move| y |mark|. Aquí están los jugadores y, después, los perfiles.

\subsection{Jugadores}

Las opciones del tablero se cambian con |\pgfgoset{goban = {...}}|. Las de los jugadores, con estas \emph{keys}:

\begin{key}{/pgfgo/players = \meta{players options}}
	Opciones de los jugadores: |neutral|, |black| y |white|. No hay que confundir jugadores con perfiles: un perfil se carga en los jugadores, y después estas \emph{keys} los cambian.
\end{key}

Las opciones de los jugadores pueden configurarse todas juntas con la \emph{key} |all|, o por separado. Los cambios realizados aquí no afectarán la definición de los perfiles. Es decir, si cargamos el perfil |pgfgo| y modificamos al jugador negro, esto no afectar\'a a la definici\'on de |pgfgo-black|, sino que a las opciones propias del jugador primero.

\begin{key}{/pgfgo/players/all = \meta{all player options}}
	Aplica las opciones a los tres jugadores.
\end{key}

\begin{key}{/pgfgo/players/neutral = \meta{neutral player options}}
	Opciones del jugador neutral (las piedras sin jugador).
\end{key}

\begin{key}{/pgfgo/players/black = \meta{first player options}}
	Opciones del primer jugador.
\end{key}

\begin{key}{/pgfgo/players/white = \meta{second player options}}
	Opciones del segundo jugador.

	\begin{codeexample}[width=4cm]
		\pgfgoset{players = {
			all = {stone = {line = {color = red}}},
			neutral = {stone = {fill = {color = yellow}}},
			black = {stone = {fill = {color = blue!60!black}}}}}

		\begin{goban}[3]
			\stones[b]{A1, B2}
			\stones[w]{C3}
			\stone{A3}
		\end{goban}
	\end{codeexample}
\end{key}

Para no escribir la ruta completa, hay cuatro comandos:

\begin{command}{\pgfgosetplayer\marg{player}\marg{player options}}
	Es lo mismo que |\pgfgoset{players = {<player> = {<player options>}}}|.
\end{command}

\begin{command}{\pgfgosetallplayers\marg{player options}}
	Aplica las opciones a los tres jugadores. Es lo mismo que |\pgfgoset{players = {all = {...}}}|.

	\begin{codeexample}[width=4cm]
		\pgfgosetallplayers{stone = {radius = 0.3cm}}
		\pgfgosetplayer{white}{stone = {fill = {color = orange}}}

		\begin{goban}[3]
			\stones[b]{A1, B2}
			\stones[w]{C3}
			\stone{A3}
		\end{goban}
	\end{codeexample}
\end{command}

\begin{command}{\pgfgocopyplayer\marg{target player}\marg{source player}}
	Copia las opciones actuales de |source player| en |target player|.

	\begin{codeexample}[width=3cm]
		\pgfgocopyplayer{neutral}{white}

		\begin{goban}[partial = A1:C3]
			\stone{B2} % Neutral, pero con el aspecto de las blancas.
			\stones[b]{A1,C3}
		\end{goban}
	\end{codeexample}
\end{command}

\begin{command}{\pgfgocopyplayers\marg{neutral source}\marg{black source}\marg{white source}}
	Copia, en este orden, |neutral source| en el jugador neutral, |black source| en el negro y |white source| en el blanco (con |\pgfgocopyplayer|). Como las copias se hacen una tras otra, no sirve para intercambiar jugadores.

	\begin{codeexample}[width=4cm]
		\pgfgocopyplayers{neutral}{white}{black}
		% No intercambia: las dos son blancas.

		\begin{goban}[3]
			\stones[b]{A1, B2}
			\stones[w]{C3}
		\end{goban}
	\end{codeexample}
\end{command}

\label{subsec:jugadores}Cada jugador tiene el mismo árbol de \emph{keys}. Cada rama se documenta con el objeto al que se aplica, sin el jugador en la ruta: para el jugador blanco, |/pgfgo/stone/radius| es |/pgfgo/players/white/stone/radius|. Las tres ramas de abajo, aunque mencionen |white|, valen para los tres jugadores.

\begin{key}{/pgfgo/players/white/stone = \meta{stone options}}
	Opciones de sus piedras y de sus prisioneros (pág.~\pageref{sec:piedras}).
\end{key}

\begin{key}{/pgfgo/players/white/move = \meta{move options}}
	Opciones de sus jugadas (pág.~\pageref{subsec:movimientos}).
\end{key}

\begin{key}{/pgfgo/players/white/mark = \meta{mark options}}
	El estilo de sus marcas (pág.~\pageref{sec:marcas}). La forma (|triangle|, |custom|, \dots) no es parte del estilo: va en las opciones de cada marca (como en |\marks[w, custom = {?}]{A3}|) o en |/pgfgo/mark|. Dentro de |\pgfgoset{players = ...}| da error.
\end{key}

La cuarta rama, |territory|, es el estilo de su territorio (pág.~\pageref{subsec:territorio}).


\subsection{Perfiles}
\label{subsec:perfiles}

Los perfiles son equivalentes a los estilos de \tikzname. En resumen, permiten guardar una configuración de tablero y piedras con un nombre y llamarla cuando se desee.

\begin{command}{\pgfgodefinegoban\opt{\oarg{options}}\marg{goban name}\marg{goban options}}
	Le da un nombre a un grupo de opciones de |goban|. Tiene dos opciones.
	
	\begin{key}{/pgfgo/profiles/goban/from = \meta{goban name}}
		Define el goban a partir de una copia de otro.
		\begin{codeexample}[code only]
			\pgfgodefinegoban[from = gobanA]{gobanB}{...}
		\end{codeexample}
		
		Es posible copiarse a sí mismo para añadir opciones, aunque no es la mejor alternativa.
		\begin{codeexample}[code only]
			\pgfgodefinegoban[from = gobanB]{gobanB}{...}
		\end{codeexample}
	\end{key}
	
	\begin{key}{/pgfgo/profiles/goban/append}
		Agrega opciones a un goban ya definido.
		\begin{codeexample}[code only]
			\pgfgodefinegoban[append]{gobanB}{...}
		\end{codeexample}
		Es equivalente a:
		\begin{codeexample}[code only]
			\pgfgodefinegoban[from = gobanB]{gobanB}{...}
		\end{codeexample}
	\end{key}
\end{command}

\begin{command}{\pgfgodefineplayer\opt{\oarg{options}}\marg{player name}\marg{player options}}
	Lo mismo, para un jugador. Sus opciones son las de |\pgfgodefinegoban|.

	\begin{key}{/pgfgo/profiles/player/from = \meta{player name}}
		Define el jugador a partir de una copia de otro.
	\end{key}

	\begin{key}{/pgfgo/profiles/player/append}
		Amplía las opciones de un jugador ya definido. Es lo mismo que |from = |\meta{player name} con su propio nombre.
	\end{key}
\end{command}

\begin{key}{/pgfgo/profile = \meta{name}}
	Aplica el perfil |name|: el goban |name| y los jugadores |name-neutral|, |name-black| y |name-white|. Se usa en |\pgfgoset|, en las opciones de un |goban| y en las de una vista (p.~\pageref{subsec:vistas}).

	Antes de aplicarlo, el tablero y los jugadores vuelven al estilo inicial (el del perfil |default|): así un perfil no hereda lo que no define, y una sombra o una imagen puesta antes se pierde. Lo que no es estilo queda como estaba: el tamaño, |partial|, |scale|, dónde van las coordenadas y cuáles son, |hoshi/at|, las capturas y las notas. Lo que va después del perfil se suma a él.

	\begin{codeexample}[width=3cm]
		\pgfgoset{stone = {shadow}}
		\begin{goban}[3, profile = default]
			\stones[b]{A1, B2}
			\stones[w]{C3}
		\end{goban}
	\end{codeexample}

	Es necesario que existan el goban |name| y los jugadores |name-neutral|, |name-black| y |name-white|. Por ejemplo, |profile = pgfgo| funciona porque existe un goban llamado |pgfgo|, y jugadores llamados |pgfgo-neutral|, |pgfgo-black| y |pgfgo-white|.
\end{key}

\begin{key}{/pgfgo/goban profile = \meta{goban name}}
	Aplica solo el goban |goban name|, después de volver el tablero al estilo inicial.
\end{key}

\begin{key}{/pgfgo/players profile = \meta{name}}
	Aplica solo los jugadores |name-neutral|, |name-black| y |name-white|, cada uno después de volver al estilo inicial.

	\begin{codeexample}[]
		\begin{goban}[3, goban profile = pgfgo]
			\stones[b]{A1, B2}
			\stones[w]{C3}
		\end{goban}

		\begin{goban}[3, players profile = pgfgo]
			\stones[b]{A1, B2}
			\stones[w]{C3}
		\end{goban}
	\end{codeexample}
\end{key}

\begin{key}{/pgfgo/players/black/profile = \meta{player name}}
	Aplica el jugador |player name| a uno de los jugadores (|neutral|, |black| o |white|), después de volverlo al estilo inicial. Con |all|, el mismo a los tres.
	\begin{codeexample}[code only]
		\pgfgoset{goban profile = pgfgo, players = {black = {profile = wain-black}}}
	\end{codeexample}
\end{key}

\begin{codeexample}[width=3cm]
	\pgfgodefinegoban{example}{background = {fill = {color = blue!05}}}
	\pgfgodefineplayer[from = default-neutral]{example-neutral}
	{stone = {
			fill = {color = gray},
			line = {color = magenta!20!black}}
	}
	
	\pgfgodefineplayer[from = example-neutral]{example-white}
	{stone = {fill = {color = yellow}}}
	
	\pgfgodefineplayer[from = example-neutral]{example-black}
	{stone = {fill = {color = blue}}}
	
	\pgfgoset{profile = example}
	
	\begin{goban}[partial = A1:C3]
		\moves[numbers = false]{A1:3}
		\stones[b]{B1,B3,C2}
		\stones[w]{C1,C3}
		\stone{B2}
	\end{goban}
\end{codeexample}

\begin{command}{\pgfgodefineprofile\opt{\oarg{options}}\marg{name}\marg{profile options}}
	Define de una vez un perfil completo para |profile|: el goban |name| y los jugadores |name-neutral|, |name-black| y |name-white|. Es lo mismo que el ejemplo de arriba, en un solo comando.

	\begin{codeexample}[width=3cm]
		\pgfgodefineprofile[from = default]{ejemplo2}{
			goban = {background = {fill = {color = blue!05}}},
			all = {stone = {line = {color = magenta!20!black}}},
			white = {stone = {fill = {color = yellow}}},
			black = {stone = {fill = {color = blue}}}}

		\pgfgoset{profile = ejemplo2}

		\begin{goban}[partial = A1:C3]
			\moves[numbers = false]{A1:3}
			\stones[b]{B1,B3,C2}
			\stones[w]{C1,C3}
			\stone{B2}
		\end{goban}
	\end{codeexample}

	\begin{key}{/pgfgo/profiles/profile/goban = \meta{goban options}}
		Las opciones del goban.
	\end{key}

	\begin{key}{/pgfgo/profiles/profile/all = \meta{player options}}
		Opciones de los tres jugadores. Van antes de las de cada uno.
	\end{key}

	\begin{key}{/pgfgo/profiles/profile/neutral = \meta{player options}}
		Opciones del jugador neutral.
	\end{key}

	\begin{key}{/pgfgo/profiles/profile/black = \meta{player options}}
		Opciones del jugador negro.
	\end{key}

	\begin{key}{/pgfgo/profiles/profile/white = \meta{player options}}
		Opciones del jugador blanco.
	\end{key}

	\begin{key}{/pgfgo/profiles/profile/from = \meta{profile name}}
		Opción del comando: cada parte empieza con la del perfil \meta{profile name} (el goban \meta{profile name}, el jugador \meta{profile name}|-neutral|, etc.).
	\end{key}
\end{command}

	\section{Cat\'alogo}
\label{sec:catalogo}


\subsection{Perfiles}
\subsubsection{default}
\begin{codeexample}[width=4cm]
	\pgfgouseprofile{default}
	
	\begin{goban}[partial = {from = A1 to G3}]
		\moves*[from = 1]{A:F 3}
		\stones[w]{A1:2}
		\stones[b]{B1:2}
		\stones[n]{C1:2}
		\stone[w, mark = circle]{D1}
		\mark[w, circle]{D2}
		\stone[b, mark = circle]{E1}
		\mark[b, circle]{E2}
		\stone[n, mark = circle]{F1}
		\mark[n, circle]{F2}
		\mark[w, tsquare]{F1}
		\mark[b, tcircle]{F1}
		\mark[n, tsquare]{F1}
	\end{goban}
\end{codeexample}


\subsubsection{pgfgo}
\begin{codeexample}[width=4cm]
	\pgfgouseprofile{pgfgo}
	
	\begin{goban}[partial = {from = A1 to G3}]
		\moves*[from = 1]{A:F 3}
		\stones[w]{A1:2}
		\stones[b]{B1:2}
		\stones[n]{C1:2}
		\stone[w, mark = circle]{D1}
		\mark[w, circle]{D2}
		\stone[b, mark = circle]{E1}
		\mark[b, circle]{E2}
		\stone[n, mark = circle]{F1}
		\mark[n, circle]{F2}
		\mark[w, tsquare]{F1}
		\mark[b, tcircle]{F1}
		\mark[n, tsquare]{F1}
	\end{goban}
\end{codeexample}



\subsubsection{pgfgosmooth}
\begin{codeexample}[width=4cm]
	\pgfgouseprofile{pgfgosmooth}
	
	\begin{goban}[partial = {from = A1 to G3}]
		\moves*[from = 1]{A:F 3}
		\stones[w]{A1:2}
		\stones[b]{B1:2}
		\stones[n]{C1:2}
		\stone[w, mark = circle]{D1}
		\mark[w, circle]{D2}
		\stone[b, mark = circle]{E1}
		\mark[b, circle]{E2}
		\stone[n, mark = circle]{F1}
		\mark[n, circle]{F2}
		\mark[w, tsquare]{F1}
		\mark[b, tcircle]{F1}
		\mark[n, tsquare]{F1}
	\end{goban}
\end{codeexample}



\subsubsection{wain}
\begin{codeexample}[width=4cm]
	\pgfgouseprofile{wain}
	
	\begin{goban}[partial = {from = A1 to G3}]
		\moves*[from = 1]{A:F 3}
		\stones[w]{A1:2}
		\stones[b]{B1:2}
		\stones[n]{C1:2}
		\stone[w, mark = circle]{D1}
		\mark[w, circle]{D2}
		\stone[b, mark = circle]{E1}
		\mark[b, circle]{E2}
		\stone[n, mark = circle]{F1}
		\mark[n, circle]{F2}
		\mark[w, tsquare]{F1}
		\mark[b, tcircle]{F1}
		\mark[n, tsquare]{F1}
	\end{goban}
\end{codeexample}


%\subsection{Colores}
%
%
%\pgfgoname\ a\~nade los siguientes nombres para colores
%\begin{verbatim}
%	\definecolor{pgfgo}{RGB}{220,179,92}
%	\colorlet{tablero2}{red!10!blue!10}
%	\colorlet{tablero3}{blue!10!green!10}
%	\definecolor{tablero4}{HTML}{9B8465}
%\end{verbatim}
%
%
%
%\begin{codeexample}[width=7cm]
%	\pgfgoset{goban = {partial = {from = A1 to C6}}}
%	
%	\begin{goban}[background = {fill = {color = pgfgo}}]
%	\end{goban}
%	
%	\begin{goban}[background = {fill = {color = tablero2}}]
%	\end{goban}
%	
%	\begin{goban}[background = {fill = {color = tablero3}}]
%	\end{goban}
%	
%	\begin{goban}[background = {fill = {color = tablero4}}]
%	\end{goban}
%\end{codeexample}

	\appendix
	\section{Novedades}
\label{sec:novedades}


\begin{itemize}
	\item \pgfgoversion:
		\begin{itemize}
			\item El cálculo (coordenadas, numeración de jugadas, memoria, etc.) se ha separado del dibujo. Lo hace un \emph{kernel}, que se elige con la opción |kernel| del paquete (pág.~\pageref{subsec:kernel}). Además del kernel escrito en L3, se ha añadido un kernel en Lua, que es el que se usa por defecto con Lua\LaTeX.
			\item Con el kernel en Lua se han añadido las capturas (|captures|, pág.~\pageref{subsec:capturas}) y la lectura de partidas en SGF (|\sgf|, pág.~\pageref{subsec:sgf}). Con el paquete |animate|, |\sgfanimation| permite recorrer una partida con botones.
			\item Se han añadido las notas «N en M» para las jugadas numeradas en puntos ya jugados (|notes| y |\pgfgomovenotes|, pág.~\pageref{subsec:notas}).
			\item SGF: variantes (|variation|), marcas del archivo (|marks|), los datos de la partida (|\sgfinfo|), la numeración de |MN| y |restart|, y el tamaño del tablero tomado del archivo (|sgf size|).
			\item SGF: colecciones de partidas (|game|), la vista (|VW|, |sgf view|), las figuras (|FG|) en |\sgfdiagrams|, las variantes como letras (|variations|), |PL| y los archivos en Latin-1.
			\item SGF: el nombre de los nodos y las anotaciones (|TE|, |BM|, |GB|, etc.) se escriben con los comentarios. |\sgfcomment| escribe el comentario de una jugada, y |\sgfdiagrams| puede escribir los comentarios debajo de cada diagrama (|comments|).
			\item Se han añadido flechas y líneas (|\arrow| y |\line|, pág.~\pageref{subsec:flechas}) y la marca |dim|. |\sgf| dibuja además |AR|, |LN|, |SL| y |DD|.
			\item El tablero dibuja los puntos de estrella (|hoshi|, pág.~\pageref{subsec:hoshi}).
			\item Se han añadido |\territory|, que marca el territorio (con un polígono, o solo con |auto| y el kernel en Lua), y |\pgfgoscore|, que cuenta los puntos.
			\item |\sgfanimation| arma la animación por capas (|layers|): cada piedra se dibuja una sola vez, y el PDF de una partida larga ocupa mucho menos.
			\item Con |captured marks = remove|, las marcas sobre una piedra capturada se van con ella.
			\item Se han añadido |\pgfgoprisoners|, que cuenta los prisioneros de cada jugador, y |\sgfdiagrams|, que dibuja una partida como serie de diagramas, como en los libros.
			\item Se ha añadido |\sgfproblem|, que dibuja un problema (tsumego): el enunciado y un diagrama por cada variante de la solución, con sus comentarios.
			\item SGF: la numeración que pide |PM| (|numbers| y |restart = auto|), y la vista (|VW|) cuando cambia durante la partida (|sgf view move|; |\sgfdiagrams| y |\sgfanimation| la usan solos). |\sgfdiagrams| escribe el resultado y el puntaje debajo del último diagrama (|result|).
			\item Las coordenadas del borde pueden ser letras, números o kanji en cada eje, contados desde cualquier lado (|columns|, |rows|, |style = japanese|, etc.).
			\item SGF: se leen archivos en GBK, Big5, Shift\_JIS, EUC-JP y EUC-KR (|CA| o |sgf charset|).
			\item Dos perfiles con imágenes, |shell| y |yunzi| (pág.~\pageref{subsec:perfilshell}): un tablero de madera con piedras de pizarra y concha, o mates. Las imágenes vienen con el paquete, en la carpeta |img|.
		\end{itemize}
	\item 0.2025.03.13: Corrección de errores. Cada |goban| empieza en la jugada 1 con negras, |\mark| y |\marks| solo existen dentro de |goban|, y lo que cae fuera de un tablero parcial ya no se dibuja. La lista completa está en el archivo |CHANGELOG.md|.
	\item 0.2025.03.10:
		\begin{itemize}
			\item Se han a\~nadido las \emph{keys} |inner| y |outer| para el \emph{rounded corner} del tablero.
			\item Se han a\~nadido tres nuevos operadores de coordenadas: |*| para todas las intersecciones, |L| para toda la columna de coordenadas L y |N| para toda la fila de coordenadas N.
		\end{itemize}
	\item 0.2025.03.07: Primer lanzamiento.
\end{itemize}

\end{document}